\documentclass[11pt]{amsart}

\usepackage{amsmath,amssymb,amsthm}
\usepackage[a4paper,margin=25mm,includeheadfoot]{geometry}
\usepackage[T1]{fontenc}
\usepackage{lmodern}
\usepackage[hidelinks]{hyperref}

\newenvironment{romanlist}{\begin{enumerate}\renewcommand{\labelenumi}{(\roman{enumi})}%
  \renewcommand{\theenumi}{\roman{enumi}}}{\end{enumerate}}
\newenvironment{alphalist}{\begin{enumerate}\renewcommand{\labelenumi}{(\alph{enumi})}%
  \renewcommand{\theenumi}{\alph{enumi}}}{\end{enumerate}}

\newtheorem{theorem}{Theorem}[section]
\newtheorem{proposition}[theorem]{Proposition}
\newtheorem{lemma}[theorem]{Lemma}
\newtheorem{corollary}[theorem]{Corollary}
\newtheorem{fact}{Fact}
\renewcommand{\thefact}{AK\arabic{fact}}
\theoremstyle{definition}
\newtheorem{definition}[theorem]{Definition}
\theoremstyle{remark}
\newtheorem{remark}[theorem]{Remark}

\newcommand{\C}{\mathbb{C}}
\newcommand{\R}{\mathbb{R}}
\newcommand{\Rp}{\mathbb{R}_{>0}}
\newcommand{\ord}{\operatorname{ord}}
\newcommand{\code}[1]{\begingroup\def\_{\textunderscore\allowbreak}\texttt{#1}\endgroup}
\setlength{\emergencystretch}{3em}

\title[Five-body central configurations]{Finiteness of planar central configurations of five
bodies for all positive masses}

\author{Zhe Chen}
\pdfinfo{/Author (Zhe Chen)}

\date{Draft of \today}

\begin{document}

\begin{abstract}
We prove that for every choice of five positive masses the planar five-body problem has only
finitely many central configurations up to rotation, translation and scaling.  This answers
Smale's sixth problem for five bodies.  Albouy and Kaloshin proved the statement for all masses
outside a closed algebraic subset of codimension two of the mass space.  We keep their
classification of singular sequences and replace the construction of the exceptional set by a
direct exclusion: we use finite boundary charts for seven of the eight diagrams that can carry
an unbounded four-product, and an exhaustive order census for Diagram 5.  We derive the
necessary mass conditions where needed and show that no germ with an unbounded four-product
lies on a curve of normalized central configurations through a positive configuration.  Five
diagrams are excluded for all complex germs, two on curves through real
configurations, and one on curves through positive configurations.  Parts of the argument are
computer-assisted: exact algebra, interval arithmetic and exact linear programming.
\end{abstract}

\maketitle

\tableofcontents


\section{Introduction}\label{sec:intro}

A configuration of $n$ point masses $m_1,\dots,m_n>0$ at positions $q_1,\dots,q_n\in\R^2$ is a
\emph{central configuration} if there is $\lambda\in\R$ with
\[
  \lambda (q_k-c) = \sum_{l\neq k} m_l\,\frac{q_l-q_k}{|q_l-q_k|^3},\qquad k=1,\dots,n,
\]
where $c$ is the centre of mass.  Central configurations give the relative equilibria of the
planar $n$-body problem and the homographic solutions; they are invariant under rotations,
translations and scalings, and one counts them modulo these symmetries.  Euler and Lagrange
found the five classes for $n=3$; Moulton~\cite{Moulton} showed that there are exactly $n!/2$
collinear classes for every $n$.  Wintner~\cite{Wintner} and Smale~\cite{Smale} asked whether
the number of classes is finite for every choice of positive masses; this is Smale's sixth
problem.  Hampton and Moeckel~\cite{HM} proved finiteness for $n=4$.  Albouy and
Kaloshin~\cite{AK} gave a new proof for $n=4$ and proved the following for $n=5$.

\begin{theorem}[{\cite[Thm.~2]{AK}}]
There is a closed algebraic subset $A$ of codimension two such that for every
$(m_1,\dots,m_5)\in\Rp^5\setminus A$ there are finitely many positive normalized central
configurations of the planar five-body problem.
\end{theorem}
Our main result removes the exceptional set.

\begin{theorem}\label{thm:main}
For every $(m_1,\dots,m_5)\in\Rp^5$ there are finitely many positive normalized central
configurations of the planar five-body problem.  Equivalently, the planar five-body problem
has finitely many central configurations up to rotation, translation and scaling.
\end{theorem}

\subsection*{Related work}
The generic finiteness of \cite{AK} for five bodies has been reproved by tropical
geometry~\cite{JL} and by an algorithmic version of the method of \cite{AK}~\cite{CC1,CC2},
which also reduces the planar six-body problem to finitely many diagrams.  In the spatial
five-body problem finiteness is known outside an explicit set of exceptional
masses~\cite{HJ}, and four-body central configurations are finite for all nonzero
masses~\cite{YZ}.  Computer-assisted methods classify the central configurations of five equal
masses~\cite{LS,MZ19} and, recently, give exact counts (between $294$ and $450$ classes) for
finitely many mass vectors on the loci (23), (29) and (32) of \cite[\S8]{AK}, the last of which
lies in the exceptional set~\cite{MZ26}; these computations treat individual mass vectors.
For all positive masses, finiteness is known for special classes such as the equilateral
pentagonal central configurations~\cite{DH}.  We are not aware of a previous proof of
Theorem~\ref{thm:main} on the whole exceptional set.

\subsection*{Strategy}
Albouy and Kaloshin embed the central configuration equations in a polynomial system and
study \emph{singular sequences}: unbounded sequences of complex normalized solutions.  To such a
sequence they attach a two-coloured \emph{diagram} recording which terms of the equations are
of maximal order, and they prove that every singular sequence approaches one of sixteen
diagrams~\cite[Fig.~11]{AK}.  If a mass vector admits infinitely many normalized central
configurations, the solutions contain an algebraic curve; along it some four-product of squared
distances is unbounded, and a singular sequence along which it diverges approaches one of the
first eight diagrams~\cite[Props.~8--9 and proof of Thm.~6]{AK}.  Each diagram forces a polynomial
relation among the masses.  Albouy and Kaloshin then push a second quantity to a second singular
limit on the same curve to obtain an independent relation, which yields the codimension-two set
$A$.

We do not construct $A$.  We show instead that, at every positive mass vector, none of the eight
diagrams can be the first limit on a curve of normalized central configurations through a
positive configuration.  Our route is as follows.
\begin{romanlist}
  \item For Diagrams 1--4 and 6--8, build charts (weighted blow-ups of the AK system) and prove
  from the estimates of \cite[\S3]{AK} that every germ converges to a finite chart boundary;
  for Diagram 5, classify the orders of the outer bodies instead.
  \item Derive the necessary mass conditions from the boundary equations, or from the isolated
  triangle in Diagram 5, where the exclusion needs one (Diagram 7 needs none); several conditions
  in \cite[\S8.3--8.8]{AK} are stated there with short or no proofs.
  \item Exclude every germ with an unbounded four-product, by a finite-order obstruction, an
  order census, or an identity on the whole curve.
\end{romanlist}
The exclusion holds at three levels.  For Diagrams 1, 2, 4, 6 and 8 no germ exists at all, even
in the complex domain.  On curves through real configurations $U=I$ is a nonzero constant
(Lemma~\ref{lem:U-constant}; compare \cite[Lemma~2]{AK}).  The order census for Diagram 5
shows that its remaining germs on such curves have only bounded four-products, while the
inertia identity for Diagram 7 forces $I\to0$ for germs with an unbounded four-product; neither
can supply the required first limit.  For Diagram 3 the real level is not enough: the modified
Roberts continuum~\cite[p.~543]{AK} is a curve of normalized central configurations at positive
masses that contains an arc of real, not positive, configurations and carries a Diagram 3 germ.
There we use positivity through an identity $\delta_{15}+2\delta_{13}=0$ that holds on the whole
curve carrying the remaining germs.

\subsection*{Computer assistance}
Several steps are computer-assisted: polynomial identities, resultants and Gr\"obner bases
over $\mathbb{Q}$ and over number fields, Sturm sequences, interval and ball arithmetic with
explicit bounds on the rounding errors, and exact rational linear programming.  The statements
that depend on such computations are marked as computer-assisted, and
Appendix~\ref{app:verification} describes the computations.

\subsection*{What we use from \texorpdfstring{\cite{AK}}{Albouy and Kaloshin}}
We use without re-proving it the classification of singular sequences (Figure~11 of
\cite{AK} and its completeness), Proposition~8 of \cite{AK}, the rules and estimates of
\cite[\S3]{AK}, and Lemma~1 and Proposition~1 of \cite{AK}.  These are listed precisely in
Section~\ref{sec:ak}.  We re-derive the mass conditions of all diagrams that need one
(Diagrams~1--6 and~8), several of which \cite[\S\S8.3--8.8]{AK} state with short or no proofs,
and we prove directly that $U$ is constant along curves of solutions.

\subsection*{Organization}
Section~\ref{sec:ak} recalls the framework of \cite{AK}.  Section~\ref{sec:reduction} reduces
Theorem~\ref{thm:main} to the exclusion of germs with an unbounded four-product.
Section~\ref{sec:exclusion} states the exclusion theorem and proves Theorem~\ref{thm:main}.
Sections~\ref{sec:d1}--\ref{sec:d8} treat the eight diagrams.  Appendix~\ref{app:verification}
describes the computer-assisted parts and Appendix~\ref{app:census} an independent cross-check
through the exceptional set of \cite{AK}.


\section{The framework of Albouy and Kaloshin}\label{sec:ak}

\subsection{Normalized central configurations}
Fix masses $m=(m_1,\dots,m_5)\in\Rp^5$.  Following \cite[\S2]{AK} we use the complex coordinates
$z_k=x_k+iy_k$, $w_k=x_k-iy_k$ and inverse distances $\delta_{kl}=\delta_{lk}$, and write
$z_{kl}=z_l-z_k$, $w_{kl}=w_l-w_k$.  The normalized system (system (4) of \cite{AK}, written as
system (10) there) is
\begin{equation}\label{eq:system}
\begin{gathered}
  z_k = \sum_{l\ne k} m_l\,\delta_{kl}^3\,(z_k-z_l),\qquad
  w_k = \sum_{l\ne k} m_l\,\delta_{kl}^3\,(w_k-w_l)\qquad (k=1,\dots,5),\\
  \delta_{kl}^2\, z_{kl}\,w_{kl} = 1\quad (1\le k<l\le 5),\qquad z_{12} = w_{12}.
\end{gathered}
\end{equation}
Its zero set $A_m\subset\C^{20}$ is the set of \emph{normalized central configurations}.  A
solution is \emph{real} if $w_k=\bar z_k$ for all $k$ and \emph{positive} if moreover
$\delta_{kl}>0$ for all $k<l$ \cite[Def.~2]{AK}.  Positive normalized central configurations are
in two-to-one correspondence with central configurations modulo rotation, translation and
scaling, so Theorem~\ref{thm:main} asserts that $A_m$ has finitely many positive points.
The map $\sigma(z,w,\delta)=(\bar w,\bar z,\bar\delta)$ preserves \eqref{eq:system} and fixes
the real solutions.  System \eqref{eq:system} is symmetric under $z\leftrightarrow w$.  A
permutation $\pi$ of the bodies, with the masses permuted, maps it to the same system with the
gauge $z_{12}=w_{12}$ replaced by $z_{\pi(1)\pi(2)}=w_{\pi(1)\pi(2)}$; a rotation
$(z,w)\mapsto(\lambda z,\lambda^{-1}w)$ restores the gauge and preserves the rest of the system,
$U$, $I$, every $r_{kl}^2$, and the orders and diagrams of germs.  We call both operations
symmetries of \eqref{eq:system}.

\subsection{Potential and inertia}
Let $U=\sum_{k<l}m_km_l\delta_{kl}$ and $I=\sum_k m_kz_kw_k$.

\begin{lemma}\label{lem:U-constant}
Every solution of \eqref{eq:system} satisfies $U=I$, and $U$ is constant on every irreducible
curve $\Gamma\subset A_m$.
\end{lemma}

\begin{proof}
Let $B\subset\C^{20}$ be the set where $\delta_{kl}^2z_{kl}w_{kl}=1$ for all $k<l$, and let
$f=2U+I$.  On $B$ no $z_{kl}w_{kl}$ vanishes, so $B$ is smooth, $(z,w)$ are local coordinates
on $B$, and $d\delta_{kl}=-\frac12\delta_{kl}^3(w_{kl}\,dz_{kl}+z_{kl}\,dw_{kl})$.  Hence, on
$B$,
\[
  \frac{\partial f}{\partial z_k}=m_k\Bigl(w_k-\sum_{l\ne k}m_l\delta_{kl}^3(w_k-w_l)\Bigr),\qquad
  \frac{\partial f}{\partial w_k}=m_k\Bigl(z_k-\sum_{l\ne k}m_l\delta_{kl}^3(z_k-z_l)\Bigr),
\]
so every point of $A_m$ is a critical point of $f|_B$ (the gauge $z_{12}=w_{12}$ plays no
role).  The scaling $(z,w,\delta)\mapsto(sz,sw,s^{-1}\delta)$ preserves $B$ and maps $f$ to
$2s^{-1}U+s^2I$; differentiating at $s=1$ gives $0=-2U+2I$ at every critical point of $f|_B$, so
$U=I$ on $A_m$.  Let $\Gamma\subset A_m$ be an irreducible curve.  At a smooth point of
$\Gamma$ the tangent line lies in the tangent space of $B$, on which $df$ vanishes; so $f$ is
locally constant on the smooth locus of $\Gamma$, which is connected, and hence constant on
$\Gamma$.  As $f=3U$ on $A_m$, $U$ is constant on $\Gamma$.
\end{proof}
The identity $U=I$
is \cite[p.~541]{AK}, and \cite[Lemma~2]{AK} shows more: $U(A_m)$ is finite.

\begin{fact}[{\cite[Lemma~1]{AK}}]\label{fact:lemma1}
If $X\subset\C^N$ is a closed algebraic subset and $f$ a polynomial, then $f(X)$ is finite or
the complement of a finite set (``$f$ is dominating'').
\end{fact}

\subsection{Singular sequences and diagrams}
Let $\mathcal Z=(z_1,\dots,z_5,Z_{12},\dots,Z_{45})$ and $\mathcal W=(w_1,\dots,w_5,W_{12},
\dots,W_{45})$ with $Z_{kl}=z_{kl}\,\delta_{kl}^3$ and $W_{kl}=w_{kl}\,\delta_{kl}^3$; these are,
up to mass factors, the terms of the equations \eqref{eq:system}.  Let $Z=\max_p|\mathcal Z_p|$
and $W=\max_q|\mathcal W_q|$.  A \emph{singular sequence} \cite[Def.~3]{AK} is a sequence of
normalized central configurations with $Z$ or $W$ unbounded, extracted so that the maximal
components are fixed and $Z^{-1}\mathcal Z$, $W^{-1}\mathcal W$ converge.  Its \emph{diagram}
has a $z$-circle at $k$ if $\lim z_k/Z\ne0$, a $z$-stroke between $k$ and $l$ if
$\lim Z_{kl}/Z\ne 0$, and likewise in the second colour \cite[\S3]{AK}.  After multiplying $z$
by $a>0$ and $w$ by $a^{-1}$ one may assume $Z=W=\epsilon^{-2}$ with $\epsilon\to0$
\cite[p.~547]{AK}.  We write $a\prec b$, $a\preccurlyeq b$, $a\approx b$, $a\sim b$ as in
\cite[\S3]{AK}.

\begin{fact}[{\cite[\S3]{AK}}]\label{fact:rules}
Rules 1a--1e and 2a--2e and Estimates 1--2 of \cite[\S3]{AK} hold for every singular sequence.
\end{fact}

\begin{fact}[{\cite[Prop.~1]{AK}}]\label{fact:prop1}
Proposition 1 of \cite{AK} (localization of the origin in a cluster) holds.
\end{fact}

\begin{fact}[{\cite[\S7]{AK}}]\label{fact:fig11}
Every singular sequence of the five-body problem approaches, after renumbering the bodies and
possibly exchanging $z$ and $w$, one of the sixteen diagrams of \cite[Fig.~11]{AK}.
\end{fact}

For pairwise distinct $i,j,k,l,m$ with $\{i,j,k,l,m\}=\{1,\dots,5\}$ the \emph{four-product} is
$p_{ij}=r_{ij}^2r_{kl}^2r_{lm}^2r_{mk}^2$, where $r_{kl}^2=z_{kl}w_{kl}$.

\begin{fact}[{\cite[Prop.~8]{AK}}]\label{fact:prop8}
Along a singular sequence approaching one of Diagrams 9--16 every four-product is bounded.
\end{fact}

Facts~AK1--AK5 are the only results of \cite{AK} that the proof uses.  The mass conditions of
\cite[\S8]{AK} for Diagrams 1--6 and 8 are re-derived below
(Propositions~\ref{prop:d1-first}, \ref{prop:d2-first}, \ref{prop:d3-first},
\ref{prop:d45-first}, \ref{prop:d6-first} and \ref{prop:d8-first}).


\section{Reduction to germs with an unbounded four-product}\label{sec:reduction}

\begin{lemma}\label{lem:distances}
Let $s\in(\C^*)^{10}$.  Only finitely many points of $A_m$ satisfy $z_{kl}w_{kl}=s_{kl}$ for all
$k<l$.
\end{lemma}

\begin{proof}
On such a point $\delta_{kl}^2=s_{kl}^{-1}$, so $\delta$ takes finitely many values; fix one.
Multiplying the $z$-equation of body $k$ by $m_k$ and summing gives $\sum_km_kz_k=0$ by
antisymmetry, and likewise $\sum_km_kw_k=0$.  Put $v_k=(z_k-z_1,w_k-w_1)\in\C^2$ and let $B$ be
the symmetric bilinear form with $B((a,b),(a,b))=ab$.  Then $B(v_k-v_l,v_k-v_l)=s_{kl}$, so the
Gram matrix $B(v_k,v_l)=\frac12(s_{1k}+s_{1l}-s_{kl})$ is fixed.  The span $V$ of the $v_k$ is
nondegenerate for $B$: either $V=\C^2$, or $V=\C v_2$ with $B(v_2,v_2)=s_{12}\ne0$.  If
$(v'_k)$ is a second family with the same Gram matrix, its span $V'$ is nondegenerate by the
same argument, and $\sum_kc_kv_k=0$ holds if and only if $\sum_kc_kB(v_k,v_l)=0$ for all $l$,
and likewise for the $v'_k$; so the $v_k$ and the $v'_k$ satisfy the same linear relations, and
$v_k\mapsto v'_k$ is an isometry of nondegenerate subspaces; by Witt's theorem it extends to an
element of the orthogonal group of $B$, which consists of the maps $(a,b)\mapsto(\lambda a,
\lambda^{-1}b)$ and $(a,b)\mapsto(\lambda b,\lambda^{-1}a)$, $\lambda\in\C^*$.  The centre of
mass determines $(z_1,w_1)$ from $(v_k)$.  For a fixed reference family $(v^0_k)$ with first
components $a^0_2$ and second components $b^0_2$ of $v^0_2$ (so $a^0_2b^0_2=s_{12}\ne0$), the gauge
$z_{12}=w_{12}$ leaves $\lambda^2=b^0_2/a^0_2$ or $\lambda^2=a^0_2/b^0_2$: at most four values.
Hence there are finitely many solutions for each $\delta$.
\end{proof}

\begin{lemma}[Arc lemma]\label{lem:arc}
If $A_m$ has infinitely many positive points, there is an irreducible algebraic curve
$\Gamma\subset A_m$ such that
\begin{romanlist}
  \item $\Gamma$ contains a nonconstant real-analytic arc of positive points;
  \item $\sigma(\Gamma)=\Gamma$;
  \item $U=I$ is constant on $\Gamma$, with value $c\in(0,\infty)$.
\end{romanlist}
\end{lemma}

\begin{proof}
In the real coordinates $(x,y,\delta)\in\R^{20}$ the positive points form a semialgebraic set.
It has finitely many connected components, so one of them is infinite, hence of positive
dimension, and the curve selection lemma \cite[Thm.~2.5.5]{BCR} gives a nonconstant Nash arc in
it.  Its Zariski
closure $\Gamma\subset A_m$ (in the complexified coordinates; passing to $(z,w)$ is linear) is an
irreducible curve, which proves (i).  Since \eqref{eq:system} has real coefficients, $\sigma(
\Gamma)$ is an irreducible curve containing the arc, so $\sigma(\Gamma)=\Gamma$.  By
Lemma~\ref{lem:U-constant}, $U=I$ is constant on $\Gamma$, and on the arc
$U=\sum m_km_l\delta_{kl}>0$.
\end{proof}

\begin{lemma}\label{lem:germ}
Let $\Gamma\subset A_m$ be an irreducible curve and $(q_n)$ a singular sequence on $\Gamma$.  There
are a Puiseux parametrization $\gamma$ of a branch of $\Gamma$ at infinity, $\gamma(t)\to\infty$ as
$t\to 0$, and $t_n\to0$ with $\gamma(t_n)$ a subsequence of $(q_n)$.  Every coordinate and every
component of $\mathcal Z$, $\mathcal W$ has a Puiseux expansion $c\,t^{\alpha}(1+o(1))$ along
$\gamma$ (or vanishes identically), and the diagram of every singular sequence extracted from
$(\gamma(t_n))$ is determined by these orders and leading coefficients: the circles and strokes
of colour $z$ are the components of $\mathcal Z$ of minimal order, and the relations $\prec$,
$\preccurlyeq$, $\approx$, $\sim$ of \cite[\S3]{AK} become relations between orders and leading
coefficients.  In particular a four-product unbounded along $(q_n)$ has negative order along
$\gamma$.
\end{lemma}

\begin{proof}
The closure of $\Gamma$ in a projective compactification has finitely many points at infinity,
each with finitely many analytic branches; a sequence tending to infinity has a subsequence on
one branch, and the normalization of the branch gives the parameter $t$.  Polynomial functions
of the coordinates are meromorphic in $t$, which gives the expansions.  If $\alpha$ is the least
order among the components of $\mathcal Z$, then $Z=\max_p|\mathcal Z_p|$ has order $\alpha$, and
$\mathcal Z_p/Z\to0$ exactly when $\ord\mathcal Z_p>\alpha$; this does not depend on the
subsequence.  The renormalization of \cite[p.~547]{AK} multiplies all $z$-quantities by one
positive factor $a(t)\approx|t|^{(\beta-\alpha)/2}$ and all $w$-quantities by $a^{-1}$, where
$\beta$ is the least order in $\mathcal W$.  This preserves every relation within one colour and
every product $z_{kl}w_{kl}$.  Replacing $a$ by the Puiseux series $\lambda(t)=(\mathcal W_q(t)/
\mathcal Z_p(t))^{1/2}$ built from maximal components, with $|\lambda|/a$ bounded above and
below, gives an analytic renormalization with the same relations; after reparametrizing
$t=s^N$, $\epsilon$ is comparable to a power of $s$.  A meromorphic function that is unbounded
along $t_n\to0$ has a pole.
\end{proof}

\begin{definition}\label{def:germ}
A \emph{Diagram-$r$ germ} is a Puiseux branch $\gamma$ of a curve in $A_m$ whose diagram, in the
sense of Lemma~\ref{lem:germ}, is Diagram $r$ of \cite[Fig.~11]{AK} after renumbering the bodies
and possibly exchanging $z$ and $w$.  It has an \emph{unbounded four-product} if some $p_{ij}$ has
negative order along $\gamma$.
\end{definition}

\begin{lemma}\label{lem:petersen}
On every irreducible curve $\Gamma\subset A_m$ some four-product is not constant.
\end{lemma}
\begin{proof}
Index the squared distances and the four-products by the ten pairs.  The pairs disjoint from
$ij$ are the sides of the triangle complementary to $ij$, so the exponent vector of $p_{ij}$ in
the variables $r_{ab}^2$ is the column $ij$ of $E=I+P$, where $P$ is the adjacency matrix of the
Petersen graph (two pairs are adjacent if they are disjoint).  The eigenvalues of $E$ are $4$,
$2$ (five times) and $-1$ (four times), so $\det E=128$, and $N=128E^{-1}$ is an integer
matrix.  As no $r_{ab}^2$ vanishes on $A_m$, $(r_{ab}^2)^{128}=\prod_{ij}p_{ij}^{N_{ab,ij}}$ there
If every $p_{ij}$ were constant on $\Gamma$, every $r_{ab}^2$ would take finitely many values on
$\Gamma$ and hence be constant, and $\Gamma$ would be finite by Lemma~\ref{lem:distances}.
\end{proof}
The lemma replaces \cite[Prop.~9]{AK}, whose proof inspects the diagrams of \cite[Fig.~11]{AK}.

\begin{proposition}[Reduction]\label{prop:reduction}
Let $m\in\Rp^5$.  If $A_m$ has infinitely many positive points, there are $r\in\{1,\dots,8\}$, a
curve $\Gamma$ as in Lemma~\ref{lem:arc}, and a Diagram-$r$ germ on $\Gamma$ with an unbounded
four-product.
\end{proposition}

\begin{proof}
Let $\Gamma$ be given by Lemma~\ref{lem:arc}.  By Lemma~\ref{lem:petersen} some four-product
$p_{ij}$ is not constant on $\Gamma$, and by Fact~\ref{fact:lemma1} it is dominating there.  A
sequence on $\Gamma$ with $p_{ij}\to\infty$ is unbounded, and a singular sequence extracted from
it approaches one of Diagrams 1--8 by Facts~\ref{fact:fig11} and~\ref{fact:prop8}.
Lemma~\ref{lem:germ} turns it into a germ.
\end{proof}

\begin{remark}\label{rem:levels}
Germs at infinity never lie on the positive arc itself: on a positive arc $U=I=c$ bounds every
$r_{kl}$ above and below.  So the signs of the boundary values of a germ say nothing about
positivity; positivity can only enter through identities valid on all of $\Gamma$.  Accordingly
we exclude germs at three levels: all complex germs; germs on curves through a real point,
where $I=U$ is a positive constant; and germs on curves with a positive arc.
\end{remark}


\section{The exclusion theorem}\label{sec:exclusion}

\begin{theorem}\label{thm:exclusion}
Let $m\in\Rp^5$, $r\in\{1,\dots,8\}$ and let $\gamma$ be a Diagram-$r$ germ with an unbounded
four-product on an irreducible curve $\Gamma\subset A_m$.
\begin{alphalist}
  \item If $r\in\{1,2,4,6,8\}$, no such germ exists.
  \item If $r\in\{5,7\}$, then $\Gamma$ contains no real point.
  \item If $r=3$, then $\Gamma$ contains no arc of positive points.
\end{alphalist}
\end{theorem}

\begin{proof}[Proof of Theorem~\ref{thm:main}]
If $A_m$ had infinitely many positive points, Proposition~\ref{prop:reduction} would give a germ
contradicting Theorem~\ref{thm:exclusion}: the curve of Lemma~\ref{lem:arc} contains a positive
arc, hence real points.
\end{proof}

The proof of Theorem~\ref{thm:exclusion} occupies Sections~\ref{sec:d1}--\ref{sec:d8}.
Table~\ref{tab:exclusions} summarizes it.

\begin{table}[ht]
\centering
\small
\begin{tabular}{cllp{5.4cm}}
\hline
$r$ & mass condition & level & method \\
\hline
1 & (23) on the triangle & complex & triangle virial against a Morse--Bott estimate \\
2 & $K_j$, $L_j$, $P_j$ & complex & isolation of every boundary label at orders 1--3,
with folds of corank one and two \\
3 & (29) $m_1m_3=m_2m_4$ & positive & orders 3--6; identity $\delta_{15}+2\delta_{13}=0$ \\
4 & (23) & complex & $w$- and $z$-force balances at orders one and three; exact census \\
5 & (23) & real & tropical census of the outer bodies \\
6 & (32) $m_1=m_2$, $m_3=m_4$ & complex & finite wing atlas \\
7 & none needed & real & inertia identity forces $I\to0$ \\
8 & (40) & complex & order-six and order-nine pairings; elimination on one curve of masses \\
\hline
\end{tabular}
\caption{The eight exclusions.  Numbers in parentheses refer to the equations of
\cite[\S8]{AK}.}\label{tab:exclusions}
\end{table}

Diagrams 1--4 and 6--8 use charts with coverage propositions and, where needed, mass conditions.
Diagram 5 instead classifies the orders of the outer bodies and uses an exact linear-programming
census to exclude unbounded four-products on curves through real configurations.
Throughout, a statement about a diagram is meant for every renumbering of the bodies (with the
masses permuted) and for both colourings; both operations are symmetries of
\eqref{eq:system}.


\section{Diagram 1}\label{sec:d1}

Diagram~1 is one of the two disconnected diagrams left by \cite[\S7.1, p.~565]{AK}: the
$z$-stroke $12$ with $z$-circles at $1$ and $2$, and the $w$-triangle with strokes $34,35,45$ and
no circle; there are no other strokes or circles.
We keep the numbering of \cite[\S8.1]{AK}: $P=\{1,2\}$, $T=\{3,4,5\}$, the letter $i$ denotes a
body of $P$ and $k,l$ denote bodies of $T$.  Put $M_P=m_1+m_2$, $M_T=m_3+m_4+m_5$,
$Z_P=M_P^{-1}(m_1z_1+m_2z_2)$, $Z_T=M_T^{-1}\sum_km_kz_k$, and define $W_P,W_T$ likewise.  On
$A_m$ we have $\sum_jm_jz_j=\sum_jm_jw_j=0$ (proof of Lemma~\ref{lem:distances}), so
$M_PZ_P+M_TZ_T=0=M_PW_P+M_TW_T$.  The coordinate equations of \eqref{eq:system} read
$E^z_j=E^w_j=0$, where
\[
  E^z_j=z_j-\sum_{l\ne j}m_l\delta_{jl}^3(z_j-z_l),\qquad
  E^w_j=w_j-\sum_{l\ne j}m_l\delta_{jl}^3(w_j-w_l).
\]
Along a germ, $a\prec b$, $a\preccurlyeq b$ and $a\approx b$ in the sense of \cite[\S3]{AK} mean
$\ord a>\ord b$, $\ord a\ge\ord b$ and $\ord a=\ord b$ (Lemma~\ref{lem:germ}).

We prove that Diagram~1 germs do not exist (Theorem~\ref{thm:d1}).  The pair torque excludes a
triangle body at the $w$-centre of the triangle (Lemma~\ref{lem:d1-central}); then every germ
converges in a weighted chart (Proposition~\ref{prop:d1-cover}); its triangle tends to the
zero-multiplier triangle, which forces relation~(23) of \cite[\S8.1]{AK}
(Proposition~\ref{prop:d1-first}); and the potential of the triangle is too large for the
Morse--Bott estimate at that triangle.

\subsection{The chart}
Let $\gamma$ be a Diagram~1 germ at $m\in\Rp^5$; renumbering the bodies and exchanging $z$ and $w$
if necessary (symmetries of \eqref{eq:system}), its diagram is the one described above.  The
coordinate equations are homogeneous of
degree one in $z$ and in $w$ separately, and the ties involve only the products $z_{jl}w_{jl}$;
so $(z,w)\mapsto(cz,c^{-1}w)$ preserves them, and it does not change the diagram
\cite[p.~547]{AK}.
By Rule~2e \cite[p.~552]{AK} the three $z$-separations of the triangle have one order, and so do
its three $w$-separations.  Taking for $c$ a power of the parameter $t$ we assume that the
$w$-separations of the triangle have order $0$; the gauge $z_{12}=w_{12}$ is not used below.
Define $\lambda,\mu,\rho,\sigma$ by
\[
  \ord z_{12}=-\lambda,\qquad \ord z_{kl}=\mu,\qquad \ord Z_P=\mu+\rho,\qquad \ord W_P=\sigma,
\]
with $\rho=+\infty$ (resp.\ $\sigma=+\infty$) if $Z_P$ (resp.\ $W_P$) vanishes identically, in
which case the terms carrying $\hat Z$ (resp.\ $\hat W$) below are absent.  Here $\sigma$ is an
exponent, not the conjugation of Section~\ref{sec:ak}.
After a finite ramification of $t$ we write
\begin{equation}\label{eq:d1-chart}
\begin{gathered}
  z_i=t^{\mu+\rho}\hat Z+t^{-\lambda}A_i,\qquad w_i=t^{\sigma}\hat W+t^{\lambda}B_i
  \qquad(i\in P),\\
  z_k=-\tfrac{M_P}{M_T}\,t^{\mu+\rho}\hat Z+t^{\mu}\zeta_k,\qquad
  w_k=-\tfrac{M_P}{M_T}\,t^{\sigma}\hat W+\eta_k\qquad(k\in T),\\
  \delta_{12}=D_{12},\qquad\delta_{ik}=t^{\lambda/2}D_{ik},\qquad\delta_{kl}=t^{-\mu/2}D_{kl},
\end{gathered}
\end{equation}
where $(A_1,A_2)=(m_2,-m_1)\,\pi$, $(B_1,B_2)=(m_2,-m_1)\,\varpi$ and
$\sum_km_k\zeta_k=\sum_km_k\eta_k=0$.  The chart coordinates are $\hat Z,\hat W,\pi,\varpi$, the
mass-centred triangle offsets $\zeta=(\zeta_k)$, $\eta=(\eta_k)$, and the ten $D_{jl}$.  The
equations get the weights $t^{\lambda}$ for $E^z_i$, $t^{-\lambda}$ for $E^w_i$, $t^{\mu/2}$ for
$E^z_k$ and $t^{3\mu/2}$ for $E^w_k$; the ten ties keep weight $1$.
The exponents of interest form the open cone
\begin{equation}\label{eq:d1-cone}
  C=\{\lambda>0,\ \mu>0,\ \sigma-\lambda>0,\ \rho+\tfrac32\mu>0,\ \rho+\lambda+\mu>0\}.
\end{equation}
For the triangle block put, with sums over $T$,
\[
  \phi_k=\sum_{l\ne k}m_lD_{kl}^3(\zeta_k-\zeta_l),\qquad
  \psi_k=\sum_{l\ne k}m_lD_{kl}^3(\eta_k-\eta_l),\qquad
  \chi_{kl}=D_{kl}^2(\zeta_l-\zeta_k)(\eta_l-\eta_k)-1,
\]
$U_T=\sum_{k<l}m_km_lD_{kl}$ and $I_T=\sum_km_k\zeta_k\eta_k$; note
$\sum_km_k\phi_k=\sum_km_k\psi_k=0$ identically.  The crossed forces are
$X^z_k=\sum_{i\in P}m_i\delta_{ik}^3(z_k-z_i)$ and $X^w_k=\sum_{i\in P}m_i\delta_{ik}^3(w_k-w_i)$.

\begin{lemma}[computer-assisted]\label{lem:d1-chart}
Substitute \eqref{eq:d1-chart} into the ten coordinate equations and the ten ties and apply the
weights.  Each of the twenty results is a finite sum $\sum_ec_e\,t^e$, where the exponents $e$
are linear forms in $(\lambda,\mu,\rho,\sigma)$, $e=0$ occurs, and every $c_e$ is a polynomial in
the chart coordinates.  Every exponent $e\ne0$ that occurs is a combination of the five forms in
\eqref{eq:d1-cone} with nonnegative rational coefficients, at least one of them positive.  No
$c_0$ involves $\hat Z$ or $\hat W$.  For $k\in T$,
\[
  t^{\mu/2}E^z_k=t^{\mu/2}(z_k-X^z_k)-\phi_k,\qquad
  t^{3\mu/2}E^w_k=t^{3\mu/2}(w_k-X^w_k)-\psi_k,
\]
and the exponents occurring in $t^{\mu/2}(z_k-X^z_k)$, respectively $t^{3\mu/2}(w_k-X^w_k)$, are
\[
\begin{gathered}
  \tfrac32\mu,\ \tfrac12(\lambda+\mu),\ \tfrac32(\lambda+\mu),\ \tfrac32\mu+\rho,\
  \tfrac32(\lambda+\mu)+\rho,\\
  \text{respectively}\\
  \tfrac32\mu,\ \tfrac32(\lambda+\mu),\ \tfrac52\lambda+\tfrac32\mu,\ \sigma+\tfrac32\mu,\
  \tfrac32(\lambda+\mu)+\sigma .
\end{gathered}
\]
\end{lemma}
We carry out the substitution exactly, with the masses as indeterminates and the exponents as
linear forms in $(\lambda,\mu,\rho,\sigma)$.  The computation determines the exponents that
occur, writes each $e\ne0$ as a combination of the five forms with explicit nonnegative rational
coefficients, and checks the displayed identities for the triangle equations by expansion
(Appendix~\ref{app:verification}).

A direct computation gives the coefficients $c_0$.  They are $A_i(1-M_PD_{12}^3)$ and
$B_i(1-M_PD_{12}^3)$ for the equations of the pair, $-\phi_k$ and $-\psi_k$ for those of the
triangle, and $M_P^2D_{12}^2\pi\varpi-1$, $-D_{ik}^2A_i\eta_k-1$ and $\chi_{kl}$ for the ties of
the pair, of the six crossed pairs and of the triangle.
The \emph{boundary fibre} $\mathcal B_1$ is their common zero set in the coordinates other than
$\hat Z,\hat W$.  On $\mathcal B_1$ we have $D_{12}^3=M_P^{-1}$ and $D_{ik}^2A_i\eta_k=-1$, and the
triangle coordinates $(\zeta,\eta,D_{34},D_{35},D_{45})$ solve the \emph{zero-multiplier triangle
system} $\phi=\psi=0$, $\chi=0$.

\subsection{Coverage}
The coverage proof uses Rules 1e and 2e and Estimates 1 and 2 of \cite[\S3]{AK}
(Fact~\ref{fact:rules}) and \cite[Prop.~1]{AK} (Fact~\ref{fact:prop1}).

\begin{lemma}\label{lem:d1-syzygies}
The following are identities of polynomials:
\begin{gather*}
  m_1E^w_1+m_2E^w_2=m_1w_1+m_2w_2-\sum_{i\in P,\,k\in T}m_im_k\delta_{ik}^3(w_i-w_k),\\
  \sum_{k\in T}m_kE^z_k=\sum_{k\in T}m_kz_k-\sum_{i\in P,\,k\in T}m_im_k\delta_{ik}^3(z_k-z_i),\\
  T_P:=\sum_{i\in P,\,k\in T}m_im_k\delta_{ik}^3(z_iw_k-w_iz_k)
  =-\sum_{i\in P}m_i\big(w_iE^z_i-z_iE^w_i\big).
\end{gather*}
In particular the torque $T_P$ of the triangle on the pair vanishes on $A_m$.
\end{lemma}

\begin{proof}
In each sum the terms between two bodies of $P$, or between two bodies of $T$, cancel in pairs.
\end{proof}

\begin{lemma}\label{lem:d1-orders}
Let $\gamma$ be a Diagram~1 germ, normalized as above.
\begin{romanlist}
  \item $\ord\delta_{12}=0$ and $\ord w_{12}=\lambda>0$; the least order of the components of
  $\mathcal Z$ is $\ord z_1=\ord z_2=\ord Z_{12}=-\lambda$.
  \item $\mu>0$, $\ord\delta_{kl}=\ord Z_{kl}=-\mu/2$, and $\lambda>\mu/2$.
  \item $\ord z_{ik}=-\lambda$, $\ord w_{ik}<\lambda$ and $\ord w_i\ge\lambda$; hence
  $\ord W_T=\ord W_P\ge\lambda$.
  \item For $k\in T$ let $o_k$ be the coefficient of $t^0$ in $w_k-W_T$.  At most one $o_k$
  vanishes.  If $o_k\ne0$, then $\ord w_{ik}=0$ and $\ord\delta_{ik}=\lambda/2$.  If $o_k=0$, then
  $f=\ord(w_k-W_T)$ satisfies $0<f<\lambda$, and $\ord w_{ik}=f$, $\ord\delta_{ik}=(\lambda-f)/2$.
\end{romanlist}
\end{lemma}

\begin{proof}
(i) The isolated component $\{1,2\}$ of the $z$-diagram has circled vertices, so it has a
maximal $z$-stroke by Rule~1e \cite[p.~546]{AK}, which can only be $12$.  Estimate~2
\cite[p.~548]{AK} gives $r_{12}\approx1$, so $\ord\delta_{12}=0$ and $\ord w_{12}=-\ord
z_{12}=\lambda$.  Estimate~1 \cite[p.~548]{AK} with $z$ and $w$ exchanged gives
$w_{12}\approx\epsilon^2$, and Estimate~2 with $z$ and $w$ exchanged gives $w_{34}\succcurlyeq\epsilon$
for the $w$-stroke $34$.  Hence $w_{12}\prec w_{34}$; a comparison within one colour is not
affected by the normalizations, so $\lambda>\ord w_{34}=0$.  Finally $1,2$ are $z$-circled and
$12$ is a $z$-stroke, and $\ord Z_{12}=\ord z_{12}+3\ord\delta_{12}=-\lambda$.

(ii) Bodies $1,2$ carry no $w$-stroke, so $T$ is an isolated component of the $w$-diagram
without $w$-circles.  By Rule~1e it has no maximal $w$-stroke.  Estimate~2 with $z$ and $w$
exchanged gives $r_{kl}\preccurlyeq1$, with $r_{kl}\approx1$ only for a maximal $w$-edge, so
$r_{kl}\prec1$; thus $\mu=\ord(z_{kl}w_{kl})=2\ord r_{kl}>0$, $\ord\delta_{kl}=
-\mu/2$ and $\ord Z_{kl}=\mu-\frac32\mu=-\frac12\mu$.  Diagram~1 has no $z$-stroke inside $T$, so
$\ord Z_{kl}>-\lambda$, that is, $\lambda>\mu/2$.

(iii) The bodies of $T$ are not $z$-circled, so $\ord z_k>-\lambda=\ord z_i$ and $\ord
z_{ik}=-\lambda$.  Since $Z_{ik}^2=(z_{ik}w_{ik}^3)^{-1}$, $\ord Z_{ik}=\frac12(\lambda-3\ord
w_{ik})$, and the absence of a $z$-stroke $ik$ gives $\ord Z_{ik}>-\lambda$, i.e.\ $\ord
w_{ik}<\lambda$.  By Estimate~1 ($z$ and $w$ exchanged) $w_{2k}\succ\epsilon^2\approx w_{12}$,
because $12$ is the only maximal $z$-edge, and $z_{2k}\approx z_{12}$.  So
\cite[Prop.~1, p.~559]{AK} with $z$ and $w$ exchanged, applied to bodies $1,2$ as in the proof of
\cite[Prop.~4, p.~564]{AK}, gives $w_1,w_2\preccurlyeq w_{12}$, i.e.\ $\ord w_i\ge\lambda$.  This
also follows from Lemma~\ref{lem:d1-syzygies}: every term of $M_PW_P=\sum m_im_k\delta_{ik}^3
(w_i-w_k)$ has order $\frac12(3\lambda-\ord w_{ik})>\lambda$, and $w_i-W_P$ has order
$\ord w_{12}=\lambda$.

(iv) Since $\sum_km_ko_k=0$, and the $o_k$ do not all vanish because $\ord w_{kl}=0$, at most one
$o_k$ vanishes.  Now $w_k-w_i=(w_k-W_T)+(W_T-w_i)$ with $\ord(W_T-w_i)\ge\lambda>0$ by (iii).  If
$o_k\ne0$ this gives $\ord w_{ik}=0$.  If $o_k=0$ then $f>0$; $f\ge\lambda$ would give $\ord
w_{ik}\ge\lambda$, against (iii), so $f<\lambda$ and $\ord w_{ik}=f$.  In both cases
$\ord\delta_{ik}=-\frac12(\ord z_{ik}+\ord w_{ik})$.
\end{proof}

\begin{lemma}[computer-assisted]\label{lem:d1-torque}
Let $\lambda,\mu,f>0$ with $f<\lambda$ and let $\rho,\sigma$ be real.  Modify
\eqref{eq:d1-chart} by $(\eta_3,\eta_4,\eta_5)=(t^f\omega,\ m_5\vartheta,\
-m_3m_5^{-1}t^f\omega-m_4\vartheta)$ and $\delta_{i3}=t^{(\lambda-f)/2}d_i$ $(i=1,2)$, all other
entries unchanged.  Then
\[
  T_P=t^{(\lambda-f)/2}\Big(m_3\sum_{i=1,2}m_id_i\,\big(d_i^2A_i\omega\big)+\sum_{g\in G}c_g\,t^g\Big),
\]
where the $c_g$ are polynomials in the chart coordinates and $G$ consists of the thirteen
exponents
\[
\begin{gathered}
  \tfrac12f,\ \tfrac32f,\ \sigma-f,\ \sigma+\tfrac12f,\ 2\lambda+\mu-f,\
  2\lambda+\mu+\tfrac12f,\ \lambda+\mu+\rho,\ \lambda+\mu+\sigma-f,\ 2\lambda+\mu+\rho-f,\\
  \lambda+\mu+\rho+\tfrac12f,\ \lambda+\mu+\sigma+\tfrac12f,\ 2\lambda+\mu+\rho+\tfrac12f,\
  \lambda+\mu+\rho+\tfrac32f,
\end{gathered}
\]
each a combination of $\lambda,\ \mu,\ \sigma-\lambda,\ \rho+\lambda+\mu,\ f,\ \lambda-f$ with
nonnegative rational coefficients, at least one of them positive.  Moreover, if
$d_i^2A_i\omega=-1$ for $i=1,2$, then $(m_1d_1)^2-(m_2d_2)^2=-(m_1^3+m_2^3)/(m_1A_1\omega)$.
\end{lemma}
We expand $T_P$ exactly, with the masses as indeterminates; no relation among the masses is used,
in particular not (23).  Each exponent in $G$ is written as a combination of the six forms with
explicit nonnegative rational coefficients, and the last identity is checked by rational
simplification (Appendix~\ref{app:verification}).

\begin{lemma}\label{lem:d1-central}
Let $m\in\Rp^5$.  Along every Diagram~1 germ all three leading triangle $w$-offsets $o_k$ are
nonzero.
\end{lemma}

\begin{proof}
Suppose not.  By Lemma~\ref{lem:d1-orders}(iv), after renumbering $T$, $o_3=0\ne o_4,o_5$,
$0<f<\lambda$, $\ord\delta_{i3}=\frac12(\lambda-f)$ and $\ord\delta_{ik}=\frac12\lambda$ for $k=4,5$.
In the first two identities of Lemma~\ref{lem:d1-syzygies} the crossed terms have orders
$\frac12(3\lambda-\ord w_{ik})$ and $\frac12(\lambda-3\ord w_{ik})$, so
$\sigma\ge\frac12(3\lambda-f)$ and $\mu+\rho\ge\frac12(\lambda-3f)$, that is,
$\sigma-\lambda\ge\frac12(\lambda-f)>0$ and $\rho+\lambda+\mu\ge\frac32(\lambda-f)>0$.
Hence $\gamma$ is described by the chart of Lemma~\ref{lem:d1-torque}, with chart coordinates of
order $\ge0$, of order $0$ for $d_i,A_i,\omega$ (so these are nonzero at $t=0$), and every
exponent in $G$ is positive.
As $\ord(Z_T-Z_P)=\mu+\rho>-\lambda$ and $\sigma>\lambda>f$, we have $z_3-z_i=-t^{-\lambda}A_i(1+o(1))$
and $w_3-w_i=t^f\omega(1+o(1))$, so the ties $\delta_{i3}^2(z_3-z_i)(w_3-w_i)=1$ give
$d_i^2A_i\omega=-1$ at $t=0$.
By Lemma~\ref{lem:d1-torque}, $t^{-(\lambda-f)/2}T_P$ tends along $\gamma$ to
$-m_3(m_1d_1+m_2d_2)$, evaluated at $t=0$, and this is nonzero because
$(m_1d_1)^2-(m_2d_2)^2=-(m_1^3+m_2^3)/(m_1A_1\omega)\ne0$.  But $T_P=0$ along $\gamma$ by
Lemma~\ref{lem:d1-syzygies}.
\end{proof}

\begin{proposition}[Coverage]\label{prop:d1-cover}
Let $m\in\Rp^5$ and let $\gamma$ be a Diagram~1 germ, normalized as above.
\begin{romanlist}
  \item $\lambda>\mu/2>0$.
  \item Every leading triangle $w$-offset $o_k$ is nonzero; hence $\ord w_{ik}=0$ and
  $\ord\delta_{ik}=\lambda/2$ for $i\in P$, $k\in T$.
  \item $\sigma\ge\frac32\lambda$ and $\mu+\rho\ge\frac12\lambda$.  In particular
  $(\lambda,\mu,\rho,\sigma)\in C$, with $\sigma-\lambda\ge\frac12\lambda$,
  $\rho+\frac32\mu\ge\frac12(\lambda+\mu)$ and $\rho+\lambda+\mu\ge\frac32\lambda$.
  \item Along $\gamma$ every chart coordinate of \eqref{eq:d1-chart} converges as $t\to0$; the
  limits of $\pi$, $\varpi$, of the $D_{jl}$, of the $\eta_k$ and of the differences
  $\zeta_k-\zeta_l$, $\eta_k-\eta_l$ $(k\ne l)$ are nonzero; and the limit point lies in the
  boundary fibre $\mathcal B_1$.
\end{romanlist}
\end{proposition}

\begin{proof}
(i) is Lemma~\ref{lem:d1-orders}(i),(ii), and (ii) follows from Lemmas~\ref{lem:d1-central}
and~\ref{lem:d1-orders}(iv).  (iii) On $\gamma$, Lemma~\ref{lem:d1-syzygies} gives
$M_PW_P=\sum m_im_k\delta_{ik}^3(w_i-w_k)$ and $M_TZ_T=\sum m_im_k\delta_{ik}^3(z_k-z_i)$.  By (ii)
and Lemma~\ref{lem:d1-orders}(iii) the terms of the first sum have order $\frac32\lambda$ and
those of the second $\frac32\lambda-\lambda=\frac12\lambda$; as $\ord Z_T=\ord Z_P=\mu+\rho$ the
two bounds follow, and the three facet bounds are arithmetic.
(iv) By the definitions of $\lambda,\mu,\rho,\sigma$, by Lemma~\ref{lem:d1-orders} and by (ii),
every chart coordinate is a Puiseux series of order $\ge0$, of order exactly $0$ for the listed
ones; so all converge.  By Lemma~\ref{lem:d1-chart} and (iii), each weighted equation equals
$c_0+\sum_{e>0}c_et^e$ along $\gamma$ with bounded $c_e$; it vanishes identically, so $c_0$
vanishes at the limit point.
\end{proof}

\begin{remark}\label{rem:d1-ak}
(a) Lemma~\ref{lem:d1-orders} and Proposition~\ref{prop:d1-cover} give the clustering scheme
$z:1\ldots345\ldots2$, $w:3\ldots12\ldots5\ldots4$, which
\cite[\S8.1, p.~572]{AK} derives from Rule~2e, Estimate~2 and \cite[Prop.~4]{AK}.
Lemma~\ref{lem:d1-central} is the torque argument of the proof of \cite[Prop.~4]{AK}, carried out
in the chart with every competing term controlled.  The proof of Theorem~\ref{thm:d1} does not
use the clustering scheme of \cite[\S8.1]{AK}.
(b) If $o_3=0$, then $w_{34}:w_{45}:w_{53}=m_5:-(m_4+m_5):m_4$ at leading order.  Relation (21)
of \cite{AK}, which the triangle $w$-rows give at leading order, together with the ties and
$z_{34}+z_{45}+z_{53}=0$, then forces $m_3^2(m_4+m_5)^4=(m_4m_5)^3$.  On the sheet of (23) written, up to scale, as
$(m_3,m_4,m_5)=(b^2,b^2(b+1)^2,(b+1)^2)$, this relation is
$b^4(b-1)^2(b+1)^8(b^2+b+1)(b^4+b^3+4b^2+b+1)=0$, and with body $4$ or $5$ in place of body $3$ it
has no positive root.  So under (23) a vanishing $o_k$ can occur only on the central-body slice
$m_4=m_5=4m_3$.  We do not use this: Lemma~\ref{lem:d1-central} holds at every positive mass.
\end{remark}

\subsection{The mass condition}

\begin{lemma}[computer-assisted]\label{lem:d1-shape}
Let $m_3,m_4,m_5>0$ and let $(\zeta,\eta,D_{34},D_{35},D_{45})$, with
$\sum_km_k\zeta_k=\sum_km_k\eta_k=0$, solve the zero-multiplier triangle system.  Put
$\alpha_{kl}=\zeta_l-\zeta_k$ and $\beta_{kl}=\eta_l-\eta_k$ for $k<l$ in $T$.  Then
\begin{alphalist}
  \item $\beta_{kl}=\Lambda^{-1}\alpha_{kl}$ for one $\Lambda\in\C^*$, and $(D_{kl}\alpha_{kl})^2=
  \Lambda$ for all $k<l$;
  \item $(D_{34}\alpha_{34},D_{35}\alpha_{35},D_{45}\alpha_{45})=s\,(-1,1,-1)$ with $s^2=\Lambda$;
  the other six sign patterns admit no solution;
  \item $m_5\alpha_{34}^2=m_3\alpha_{45}^2=m_4\alpha_{35}^2$ and $Q(m_3,m_4,m_5)=0$, where
  \[
    Q=2m_3m_4m_5(m_3+m_4+m_5)-m_3^2m_4^2-m_4^2m_5^2-m_5^2m_3^2 .
  \]
\end{alphalist}
\end{lemma}
The proof is an exact computation with $m_3,m_4,m_5$ as indeterminates.  Explicit cofactor
identities give (a), and in each sign pattern the relation $\alpha_{35}=\alpha_{34}+\alpha_{45}$
eliminates the separations and leaves a polynomial relation among the masses: for the six
patterns excluded in (b) the polynomial is a square plus a polynomial with positive
coefficients, and for the pattern $s\,(-1,1,-1)$ it is $-Q$ (Appendix~\ref{app:verification}).
For positive masses $Q=0$ is equivalent to (23): with $x_k=m_k^{-1/2}$,
$Q=(m_3m_4m_5)^2(x_3+x_4+x_5)(-x_3+x_4+x_5)(x_3-x_4+x_5)(x_3+x_4-x_5)$, and the first factor is
positive.

The torus $\mathbb T=\{(\theta_1,\theta_2,\theta_3)\in(\C^*)^3:\theta_3^2=\theta_1\theta_2\}$ acts by
$(\zeta,\eta,D)\mapsto(\theta_1\zeta,\theta_2\eta,\theta_3^{-1}D)$.  It multiplies $\phi$ by
$\theta_1\theta_3^{-3}$, $\psi$ by $\theta_2\theta_3^{-3}$ and $I_T$ by $\theta_1\theta_2$, and it
fixes $\chi$.

\begin{lemma}[computer-assisted]\label{lem:d1-mb}
Let $b,c>0$, $(m_3,m_4,m_5)=(b^2c^2,\,b^2(b+c)^2,\,c^2(b+c)^2)$ and $\kappa>0$.  Let $y^*$ be the
point with $\eta=\zeta$, $\sum_km_k\zeta_k=0$, $\zeta_4-\zeta_3=\kappa m_5^{-1/2}$,
$\zeta_5-\zeta_3=-\kappa m_4^{-1/2}$, $D_{34}=m_5^{1/2}/\kappa$, $D_{35}=m_4^{1/2}/\kappa$ and
$D_{45}=-m_3^{1/2}/\kappa$.  In the coordinates $y=(\zeta_3,\zeta_4,\eta_3,\eta_4,D_{34},D_{35},
D_{45})$ let $F=(\phi_3,\phi_4,\psi_3,\psi_4,\chi_{34},\chi_{35},\chi_{45})$.  Then $F(y^*)=0$; the
Jacobian matrix $F'(y^*)$ has rank $5$, and its kernel is spanned by the generators
$(\zeta_3,\zeta_4,0,0,-\frac12D)$ and $(0,0,\eta_3,\eta_4,-\frac12D)$ of the $\mathbb T$-action; and
\[
  I_T(y^*)=\frac{\kappa^2}{m_3+m_4+m_5}\Big(\frac{m_3m_4}{m_5}+\frac{m_3m_5}{m_4}
  +\frac{m_4m_5}{m_3}\Big)>0 .
\]
\end{lemma}
The computation is exact in the field $\mathbb Q(b,c,\kappa)$: it gives $F(y^*)=0$, shows that
$F'(y^*)$ annihilates the two generators, and factors a $5\times5$ minor of $F'(y^*)$ and the
value $I_T(y^*)$ into factors with positive coefficients, so that both are positive for
$b,c,\kappa>0$ (Appendix~\ref{app:verification}).  One checks that every positive solution of
$m_3^{-1/2}=m_4^{-1/2}+m_5^{-1/2}$ is of the stated form.
By the equivariance above, the same rank and kernel statements hold at every point of
$O=\mathbb T\cdot y^*$, and $I_T\ne0$ on $O$.

\begin{proposition}[First condition]\label{prop:d1-first}
Let $m\in\Rp^5$ and let $\gamma$ be a Diagram~1 germ.  Then $Q(m_3,m_4,m_5)=0$, that is, the
triangle masses satisfy relation (23) of \cite[\S8.1]{AK}.  After renumbering $T$ so that
$m_3^{-1/2}=m_4^{-1/2}+m_5^{-1/2}$, the limit $y_0$ of
$y(t)=(\zeta_3,\zeta_4,\eta_3,\eta_4,D_{34},D_{35},D_{45})(t)$ lies on the orbit $O$ of
Lemma~\ref{lem:d1-mb}.
\end{proposition}

\begin{proof}
By Proposition~\ref{prop:d1-cover}(iv) the limit lies in $\mathcal B_1$, so its triangle part
solves the zero-multiplier triangle system.  In the notation of Lemma~\ref{lem:d1-shape},
$\phi_4=m_3\alpha_{34}D_{34}^3-m_5\alpha_{45}D_{45}^3$ and
$\phi_5=m_4\alpha_{45}D_{45}^3+m_3\alpha_{35}D_{35}^3$, and $m_3\phi_3=-m_4\phi_4-m_5\phi_5$;
likewise for $\psi$.  So $\phi=0$ and $\psi=0$ are relations (22) and (21) of \cite[\S8.1]{AK} in
cross-multiplied form, and Lemma~\ref{lem:d1-shape} gives $Q=0$.
After the renumbering, Lemma~\ref{lem:d1-shape}(c) and $\alpha_{35}=\alpha_{34}+\alpha_{45}$
leave only $(\alpha_{34},\alpha_{45},\alpha_{35})=\kappa(x_5,-x_3,-x_4)$ with $\kappa\in\C^*$, and
Lemma~\ref{lem:d1-shape}(a),(b) then determine $\eta$ and $D$ from $(\kappa,s)$.  The torus
$\mathbb T$ acts transitively on these pairs $(\kappa,s)$, and $y^*$ is the solution with
$\kappa>0$, $s=-1$.  Hence the solutions form the single orbit $O$, and $y_0\in O$.
\end{proof}

\subsection{Exclusion}

\begin{lemma}[computer-assisted]\label{lem:d1-virial}
For arbitrary $\zeta,\eta\in\C^3$ and $D_{34},D_{35},D_{45}\in\C$,
\begin{gather*}
  \sum_{k\in T}m_k\zeta_k\psi_k=U_T+\sum_{k<l}m_km_lD_{kl}\chi_{kl},\\
  2\,dU_T+\sum_{k\in T}m_k\big(\psi_k\,d\zeta_k+\phi_k\,d\eta_k\big)
  =\sum_{k<l}m_km_l\big(D_{kl}\,d\chi_{kl}-2\chi_{kl}\,dD_{kl}\big).
\end{gather*}
\end{lemma}
Both identities are checked by exact expansion, with the masses as indeterminates
(Appendix~\ref{app:verification}).
On the tie manifold $\{\chi=0\}$ they reduce to $U_T=\sum_km_k\zeta_k\psi_k$ and
$dU_T=-\frac12\sum_km_k(\psi_k\,d\zeta_k+\phi_k\,d\eta_k)$; both vanish on $O$.

\begin{lemma}\label{lem:d1-estimate}
Every $y_0\in O$ has a neighbourhood $V\subset\C^7$ with a constant $c_V$ such that
$|U_T(y)|\le c_V\,|(\phi_3,\phi_4,\psi_3,\psi_4)(y)|^2$ for every $y\in V$ with $\chi(y)=0$.
\end{lemma}

\begin{proof}
By Lemma~\ref{lem:d1-mb} and equivariance, five components $F_5$ of $F$ have independent
differentials at $y_0$.  The orbit map $\theta\mapsto\theta y_0$ is an immersion of a
neighbourhood of $1\in\mathbb T$ onto a $2$-dimensional piece of $O$ inside the $2$-dimensional
manifold $F_5^{-1}(0)$ near $y_0$, so $F_5^{-1}(0)\cap V\subset O$ for small $V$, and
$\operatorname{dist}(y,O)\le c\,|F_5(y)|\le c\,|(\phi_3,\phi_4,\psi_3,\psi_4)(y)|$ for $y\in V$
with $\chi(y)=0$ (implicit function theorem).  Near $y_0$ the manifold $\{\chi=0\}$ is the graph
of $D$ over $(\zeta_3,\zeta_4,\eta_3,\eta_4)$, since $\partial\chi_{kl}/\partial D_{kl}=
2D_{kl}^{-1}(\chi_{kl}+1)\ne0$.  In these coordinates $U_T$ and its differential vanish on $O$
(Lemma~\ref{lem:d1-virial}), so Taylor's formula at a nearest point of $O$ gives $|U_T(y)|\le
c'\operatorname{dist}(y,O)^2$.
\end{proof}

\begin{theorem}\label{thm:d1}
Let $m\in\Rp^5$.  No curve in $A_m$ has a Diagram~1 germ (Definition~\ref{def:germ}).  In
particular there is no Diagram~1 germ with an unbounded four-product, which is the case $r=1$ of
Theorem~\ref{thm:exclusion}\,(a).
\end{theorem}

\begin{proof}
Let $\gamma$ be a Diagram~1 germ, normalized as above and with $T$ renumbered as in
Proposition~\ref{prop:d1-first}.  Along $\gamma$ the ties of the triangle give $\chi(y(t))=0$
exactly, and $y(t)\to y_0\in O$.

\emph{Forcing.}  Along $\gamma$ we have $E^z_k=E^w_k=0$, so Lemma~\ref{lem:d1-chart} gives
$\phi_k=t^{\mu/2}(z_k-X^z_k)$ and $\psi_k=t^{3\mu/2}(w_k-X^w_k)$, with the exponents listed
there and bounded coefficients.  By Proposition~\ref{prop:d1-cover}(iii),
$\frac32\mu+\rho\ge\frac12(\lambda+\mu)$.  Hence $|\phi(y(t))|=O(|t|^r)$ with
$r=\min(\frac32\mu,\frac12(\lambda+\mu))>0$, and $|\psi(y(t))|=O(|t|^{3\mu/2})$.

\emph{Morse--Bott.}  By Lemma~\ref{lem:d1-estimate}, $|U_T(y(t))|=O(|t|^{2r})$.

\emph{Virial.}  By Lemma~\ref{lem:d1-virial} and $\chi=0$,
$U_T(y(t))=\sum_km_k\zeta_k\psi_k=t^{3\mu/2}\sum_km_k\zeta_k(w_k-X^w_k)$.  Since
$w_k=W_T+\eta_k$ and $\sum_km_k\zeta_k=0$, we get $\sum_km_k\zeta_kw_k=I_T(y(t))$, while
$\ord X^w_k\ge\frac32\lambda$ by Proposition~\ref{prop:d1-cover}(ii).  Hence
$U_T(y(t))=t^{3\mu/2}\big(I_T(y_0)+o(1)\big)$ with $I_T(y_0)\ne0$.

\emph{Contradiction.}  The last two steps give $\frac32\mu\ge2r=\min(3\mu,\lambda+\mu)$.  Since
$\mu>0$, this means $\lambda+\mu\le\frac32\mu$, i.e.\ $\lambda\le\frac12\mu$, which contradicts
Proposition~\ref{prop:d1-cover}(i).
\end{proof}

\begin{remark}\label{rem:d1-wall}
The window $\lambda>\mu/2$ is used only in the last step.  In the chart family
\eqref{eq:d1-chart} the outer bodies are $z$-circled exactly when $\lambda\ge\mu/2$.  The wall
$\lambda=\mu/2$ is the Diagram~4 chart after a rotation; there the orders $\frac32\mu$ and
$\lambda+\mu$ coincide and the argument gives nothing, which is why Diagram~4 needs the analysis
of Section~\ref{sec:d45}.  For $\lambda<\mu/2$ the diagram is Diagram~5.  No bound on
four-products is used, so Theorem~\ref{thm:d1} excludes Diagram~1 germs as first and as second
limits.
\end{remark}



\section{Diagram 2}\label{sec:d2}

Diagram~2 of \cite[Fig.~11]{AK} consists of the $z$-segment $12$ and the $w$-triangle $345$,
and every vertex carries a circle of the colour of its edges \cite[\S8.2, p.~572]{AK}.
Exchanging $z$ and $w$ and renumbering the bodies are symmetries of \eqref{eq:system}, so we
assume that a Diagram~2 germ has exactly this diagram.
Throughout this section $A,A'$ denote the pair bodies, $\{A,A'\}=\{1,2\}$, and $B,B'$ bodies of
the triangle $\{3,4,5\}$; $M=m_1+m_2$, $M_T=m_3+m_4+m_5$ and $m'=(m_3,m_4,m_5)$.  The letter $t$
is a chart coordinate; the parameter of a germ (called $t$ in Lemma~\ref{lem:germ}) is called
$s$, and $\ord$ is the order in $s$.  The letters $x,y,u$ are those of \cite[\S8.2]{AK}.

The proof has four steps.  We blow up \eqref{eq:system} along the diagram; every Diagram~2 germ
converges to a point of the boundary fibre $t=0$, a \emph{label}
(Proposition~\ref{prop:d2-cover}).  A torque identity of the pair forces $\Phi=0$ at the limit
label, which puts the masses on the loci $L_j$, $K_j$, $P_j$ of \cite[\S8.2]{AK}
(Proposition~\ref{prop:d2-first}).  Where the selected twenty equations define a graph over
$t$, the two remaining equations reduce to $tq(t)=0$, with $q(0)=\Phi$ and $q'(0)$ given by a
quadratic form of the triangle (Lemmas~\ref{lem:d2-isolation} and~\ref{lem:d2-second}).  At
fold labels without such a graph we instead use one- or two-parameter local reductions.
Finally every label with $\Phi=0$ at positive masses is shown not to be a limit of germs
(\S\ref{sec:d2-cases}), in most regular cases because $q'(0)\ne0$.  No reality condition is used.

\subsection{The chart and its labels}\label{sec:d2-chart}

The chart has coordinates $(t,Z,W,D)\in\C^{21}$, $Z=(Z_1,\dots,Z_5)$, $W=(W_1,\dots,W_5)$,
$D=(D_{kl})_{k<l}$, related to the coordinates of \eqref{eq:system} by
\begin{equation}\label{eq:d2-chart}
\begin{gathered}
  z_A=t^{-1}Z_A,\quad w_A=tW_A,\quad z_B=tZ_B,\quad w_B=t^{-1}W_B,\\
  \delta_{12}=D_{12},\quad \delta_{BB'}=D_{BB'},\quad \delta_{AB}=tD_{AB}.
\end{gathered}
\end{equation}
The single-index capitals $Z_k,W_k$ are chart coordinates, unrelated to the two-index quantities
$Z_{kl},W_{kl}$ used for diagrams.  Multiplying the coordinate equations of \eqref{eq:system} for
$z_A,w_A,z_B,w_B$ by $t,t^{-1},t^{-1},t$ gives the polynomials
\begin{equation}\label{eq:d2-rows}
\begin{aligned}
  F^Z_A&=Z_A-m_{A'}D_{12}^3(Z_A-Z_{A'})-t^3\textstyle\sum_Bm_BD_{AB}^3(Z_A-t^2Z_B),\\
  F^W_A&=W_A-m_{A'}D_{12}^3(W_A-W_{A'})+t\textstyle\sum_Bm_BD_{AB}^3(W_B-t^2W_A),\\
  F^Z_B&=Z_B-\textstyle\sum_{B'\ne B}m_{B'}D_{BB'}^3(Z_B-Z_{B'})
         +t\textstyle\sum_Am_AD_{AB}^3(Z_A-t^2Z_B),\\
  F^W_B&=W_B-\textstyle\sum_{B'\ne B}m_{B'}D_{BB'}^3(W_B-W_{B'})
         -t^3\textstyle\sum_Am_AD_{AB}^3(W_B-t^2W_A),
\end{aligned}
\end{equation}
and the distance equations become $T_{12}=D_{12}^2(Z_2-Z_1)(W_2-W_1)-1$,
$T_{BB'}=D_{BB'}^2(Z_{B'}-Z_B)(W_{B'}-W_B)-1$ ($B<B'$) and
$T_{AB}=D_{AB}^2(t^2Z_B-Z_A)(W_B-t^2W_A)-1$.  The gauge of \eqref{eq:system} is replaced by
$G_1=Z_2-Z_1-1$ and $G_2=W_5-1$.
Let $\mathcal X_m\subset\C^{21}$ be the zero set of these $22$ polynomials and
$\mathcal B_m=\mathcal X_m\cap\{t=0\}$; the points of $\mathcal B_m$ are the \emph{labels}.

\begin{definition}[Shape factors]\label{def:d2-shapes}
For $\varsigma\in\{\pm1\}^3$ let $N_\varsigma(w_3,w_4,w_5)$ be the numerator of
$w_3/(\varsigma_3(w_5-w_4)^2)+w_4/(\varsigma_4(w_3-w_5)^2)+w_5/(\varsigma_5(w_4-w_3)^2)$.  The
quintics $S_1,S_2,S_3$ are $N_\varsigma$ for $\varsigma=(1,-1,-1),(1,-1,1),(1,1,-1)$, and
$N_{(1,1,1)}$ is a constant multiple of $S_4S_5$, where
\begin{gather*}
  S_4=w_3^3+w_4^3+w_5^3-w_3^2w_4-w_3w_4^2-w_3^2w_5-w_3w_5^2-w_4^2w_5-w_4w_5^2+3w_3w_4w_5,\\
  S_5=w_3^2+w_4^2+w_5^2-w_3w_4-w_3w_5-w_4w_5 .
\end{gather*}
The quartics $S_6,S_7,S_8$ are the factors of
$\operatorname{Res}_p\bigl((w_4-w_3)^3-p^3(w_5-w_4)^3,\ (w_4-w_3)^3-(-1-p)^3(w_3-w_5)^3\bigr)$
(which is $S_5^3S_6S_7S_8$ up to a constant) whose coefficients of $w_3^3w_4$ are $-2$, $1$ and
$-5$.  These are the factors of the complex three-body central configuration locus of
\cite[Appendix, pp.~585--586]{AK}.
\end{definition}

\begin{lemma}[computer-assisted]\label{lem:d2-fibre}
Let $m\in\Rp^5$.  A point $p=(0,Z,W,D)$ is a label if and only if
\begin{alphalist}
  \item $Z_1=-m_2/M$, $Z_2=m_1/M$, $D_{12}^3=1/M$, $W_1=-m_2D_{12}$ and $W_2=m_1D_{12}$;
  \item $F^Z_B=F^W_B=T_{BB'}=0$ at $t=0$ (system \eqref{eq:system} for the bodies $3,4,5$,
  without its gauge) and $W_5=1$;
  \item $D_{AB}^2Z_AW_B=-1$ for all $A,B$.
\end{alphalist}
At a label $\sum_Bm_BZ_B=\sum_Bm_BW_B=0$, every $W_B\ne0$, the $W_B$ are pairwise distinct, and
the \emph{triangle} $(Z_B,W_B,D_{BB'})$ is of one of two kinds.
\begin{romanlist}
  \item \emph{Lagrange:} $D_{BB'}=M_T^{-1/3}\omega^{a_{BB'}}$, $\omega=e^{2\pi i/3}$,
  $a_{BB'}\in\{0,1,2\}$ ($27$ classes), and $\kappa=(W_4-W_3)/(W_5-W_4)$ is a root of
  $\omega^{a_{45}}\kappa^2+(\omega^{a_{34}}+\omega^{a_{45}}-\omega^{a_{35}})\kappa+\omega^{a_{34}}$.
  The root is double exactly for the six classes with $a_{34},a_{35},a_{45}$ distinct; these are
  the \emph{fold} triangles.  $(W_3,W_4,W_5)$ is a zero of $S_5$ if $a_{34}=a_{35}=a_{45}$ or the
  $a_{BB'}$ are distinct, of $S_6$ if $a_{34}=a_{35}\ne a_{45}$, of $S_7$ if
  $a_{34}=a_{45}\ne a_{35}$, and of $S_8$ if $a_{35}=a_{45}\ne a_{34}$.
  \item \emph{Euler:} $Z_B=c^{-2}W_B$ for some $c\in\C^*$, $D_{BB'}=e_{BB'}c/(W_{B'}-W_B)$ for
  $B<B'$ with $e_{BB'}=\pm1$, and $(W_3,W_4,W_5)$ is a zero of $S_1$, $S_2$, $S_3$ or $S_4S_5$
  according to the class of $(e_{34},e_{35},e_{45})$ modulo a global sign.  An Euler triangle on
  $S_5$ is a Lagrange fold triangle.
\end{romanlist}
\end{lemma}

\begin{proof}
At $t=0$ the crossed terms of \eqref{eq:d2-rows} vanish and $T_{AB}=-D_{AB}^2Z_AW_B-1$; this
gives (b), (c) and, with $G_1$, the system $Z_A=m_{A'}D_{12}^3(Z_A-Z_{A'})$,
$W_A=m_{A'}D_{12}^3(W_A-W_{A'})$, $T_{12}=0$, whose solutions are (a).  Summing $m_BF^Z_B$ and
$m_BF^W_B$ gives the centre-of-mass relations.  If $W_B=0$ then $T_{1B}=-1$, and if
$W_B=W_{B'}$ then $T_{BB'}=-1$.  If the vectors $(1)_B,(Z_B)_B,(W_B)_B$ are linearly independent,
then $D_{BB'}^3=1/M_T$ for all three edges.  Indeed, the linear map
$X\mapsto(\sum_{B'\ne B}m_{B'}D_{BB'}^3(X_B-X_{B'}))_B$ of $\C^3$ fixes $(Z_B)$ and $(W_B)$ by (b)
and kills $(1)_B$, and so does $X\mapsto X-M_T^{-1}(\sum_Bm_BX_B)(1)_B$, by the centre-of-mass
relations; as the three vectors form a basis, the two maps agree, and their entries at
$(B,B')$, $B\ne B'$, give $m_{B'}D_{BB'}^3=m_{B'}/M_T$.  Then the coordinate equations
reduce to the centre-of-mass relations, the ties give
$Z_{B'}-Z_B=M_T^{2/3}\omega^{a_{BB'}}/(W_{B'}-W_B)$, and $Z_5-Z_3=(Z_4-Z_3)+(Z_5-Z_4)$ is the
quadratic for $\kappa$.  Otherwise $Z=\alpha W+\beta$; the centre of mass gives $\beta=0$, the
ties give $\alpha\ne0$, and with $\alpha=c^{-2}$ they give the form of $D_{BB'}$.  For the double
roots write $c_{BB'}=\omega^{a_{BB'}}$: if two of the $c_{BB'}$ agree, the discriminant is
$b(b-4a)\ne0$ with $|a|=|b|=1$ ($a$ the repeated value, $b$ the remaining one), while if they are
distinct, $c_{34}+c_{45}=-c_{35}$ and $c_{34}c_{45}=c_{35}^2$, so the discriminant vanishes.  The
assignment of the shape factors is computer-assisted.
\end{proof}

The assignment is an exact computation over $\mathbb Q(\omega)$.  For each of the $27$ Lagrange
classes we reduce the factor $S_j$ of the class, as a polynomial in $\kappa$, modulo the quadratic
for $\kappa$, and for a fold class we evaluate it at the double root; for each of the eight sign
vectors $(e_{34},e_{35},e_{45})$ we identify the corresponding $N_\varsigma$ with one of $\pm S_1$,
$\pm S_2$, $\pm S_3$, $\pm S_4S_5$; and for the Euler triangles on $S_5$ we check that the three
numbers $D_{BB'}^3$ coincide modulo $S_5$ (Appendix~\ref{app:verification}).

\begin{definition}[Mass polynomials]\label{def:d2-mass}
For positions $W_3,W_4,W_5$ put $f_B=m_B^2/W_B$,
$\Lambda=f_3^2+f_4^2+f_5^2-2f_3f_4-2f_3f_5-2f_4f_5$ and
$\Theta=3(f_3^2+f_4^2+f_5^2)+2f_3f_5+2f_4f_5-6f_3f_4$.  Equations (25) and (28) of
\cite[\S8.2]{AK} are
\begin{equation}\label{eq:d2-25-28}
  \Lambda=0,\qquad
  \Lambda\,y^2+2(f_3-f_4+f_5)(f_3-f_4-f_5)\,y+\tfrac12(\Theta-\Lambda)=0 .
\end{equation}
For a rational expression $H$ in $W_3,W_4,W_5$ whose denominator is a monomial, let $H^\flat$ be
the polynomial in $W_4$ obtained by clearing the position denominators, substituting $W_5=1$ and
$W_3=-(m_4W_4+m_5)/m_3$ (centre of mass), and clearing the resulting power of $m_3$.  Define
$L_j(m')$ and $P_j(y)$ by
\[
  \operatorname{Res}_{W_4}(S_j^\flat,\Lambda^\flat)=c_j\,L_j,\qquad
  \operatorname{Res}_{W_4}\bigl(S_j^\flat,(28)^\flat\bigr)=m_3^{2\deg S_j}P_j(y),
\]
where $c_j$ has no zero on $\Rp^3$; for $j=5$ the first resultant is
$m_3^4M_T^2L_5L_5'$.  Write $P_j(y)=E_j(y^2)+yO_j(y^2)$ and $K_j=P_j(0)=E_j(0)$.
\end{definition}

These are the factors $L_1,\dots,L_8,L_5'$ and the polynomials $P_1,\dots,P_8$ of
\cite[pp.~573--575]{AK}, reconstructed from the appendix factors; the displayed coefficients and
published symmetries of \cite{AK} are reproduced.  In particular
$2L_5=(m_3-m_4)^2+(m_4-m_5)^2+(m_5-m_3)^2$, and $L_5'$ has degree $12$ and $91$ coefficients, all
positive.
Permutations of $(m_3,m_4,m_5)$ permute $\{L_1,L_2,L_3\}$ and $\{L_6,L_7,L_8\}$ up to nonzero
constants; the exchange of bodies $3$ and $4$ sends $S_1$ to $-S_2$, exchanges $P_1\leftrightarrow
P_2$ and $P_6\leftrightarrow P_7$, and fixes $P_3,P_4,P_5,P_8$.
These statements are proved by exact computation over $\mathbb Q$, with the masses and $y$ as
indeterminates: the resultants defining $L_j$ and $L_5'$ are computed and factored, those
defining $P_j$ are determined by evaluation on unisolvent grids of points modulo three primes,
using a priori bounds on their degrees and coefficients, and the decompositions $P_j=E_j+yO_j$,
the leading coefficients and the symmetries are checked by expansion and substitution
(Appendix~\ref{app:verification}).
\subsection{Coverage}\label{sec:d2-cover}

\begin{proposition}[Coverage]\label{prop:d2-cover}
Let $m\in\Rp^5$ and let $\gamma$ be a Diagram~2 germ on a curve in $A_m$.  Along $\gamma$ put
$t=(z_{12}w_5)^{-1/2}$ \textup{(}either branch\textup{)} and $\lambda=(tz_{12})^{-1}$, and let
$(Z,W,D)$ be given by \eqref{eq:d2-chart} applied to the rotated configuration
$(\lambda z,\lambda^{-1}w,\delta)$.  Then $t,Z,W,D$ are Puiseux series in $s$ with $\ord t>0$ and
nonnegative orders, $(t,Z,W,D)\in\mathcal X_m$ for $s\ne0$, and the limit $p$ as $s\to0$ is a label at
which $Z_1,Z_2$, $W_3,W_4,W_5$ and all $D_{kl}$ are nonzero.
\end{proposition}

\begin{proof}
The equations of \eqref{eq:system} other than its gauge are invariant under the rotations
$(z,w)\mapsto(\lambda z,\lambda^{-1}w)$, the chart rows are these equations times powers of $t$,
and $G_1=\lambda tz_{12}-1=0$, $G_2=\lambda^{-1}tw_5-1=t^2z_{12}w_5-1=0$; so the chart point lies
in $\mathcal X_m$.  Rotating $\gamma$ beforehand changes neither $t$ nor the chart point, so for the
orders we may normalize as in Lemma~\ref{lem:germ} and \cite[p.~547]{AK}: the maximal components
of $\mathcal Z$ and $\mathcal W$ have one order, and we use units in which it is $-2$.  By
Fact~\ref{fact:rules}:
\begin{enumerate}
  \item The pair is an isolated component of the $z$-diagram with $z$-circles, so its $z$-stroke
  is maximal (Rule~1e): $\ord z_{12}=-2$; by Estimate~2, $r_{12}\approx1$, so $\ord w_{12}=2$ and
  $\ord\delta_{12}=0$.
  \item The triangle is an isolated $w$-component with $w$-circles; it has a maximal $w$-stroke
  (Rule~1e), and by Rule~2e all three are maximal: $\ord w_{BB'}=-2$, $\ord z_{BB'}=2$,
  $\ord\delta_{BB'}=0$.
  \item The circles give $\ord z_A=\ord w_B=-2$.  The bodies $3,4,5$ carry no $z$-stroke and
  $1,2$ no $w$-stroke, hence no circle of that colour (Rule~1a): $\ord z_B>-2$, $\ord w_A>-2$.
  \item Hence $\ord z_{AB}=\ord w_{AB}=-2$ and $\ord\delta_{AB}=2$.
  \item In $w_A=\sum_{l\ne A}m_l\delta_{Al}^3(w_A-w_l)$ the pair term has order $2$ and the
  crossed terms order $4$, so $\ord w_A\ge2$; likewise $\ord z_B\ge2$.
  \item So $\ord t=2$ and $\ord\lambda=0$; $Z_A=\lambda tz_A$, $W_B=\lambda^{-1}tw_B$, $D_{12}$,
  $D_{BB'}$ and $D_{AB}=\delta_{AB}/t$ have order exactly $0$, and $W_A=\lambda^{-1}w_A/t$,
  $Z_B=\lambda z_B/t$ have order $\ge0$.
\end{enumerate}
Hence every chart coordinate converges, $t\to0$, and the coordinates of order exactly $0$ have
nonzero limits.  The limit lies in the closed set $\mathcal X_m$ and has $t=0$.
\end{proof}

\subsection{First conditions}\label{sec:d2-first}

\begin{lemma}[Syzygies]\label{lem:d2-syzygy}
Identically in $(t,Z,W,D)$,
\begin{gather*}
  \textstyle\sum_{k=1}^5m_k(W_kF^Z_k-Z_kF^W_k)=0,\qquad
  \sum_Am_A(W_AF^Z_A-Z_AF^W_A)=t\,Q,\\
  Q=\textstyle\sum_{A,B}m_Am_BD_{AB}^3(t^4W_AZ_B-Z_AW_B).
\end{gather*}
At $t=0$ the only rows with a nonzero $t$-derivative are $F^W_A$ and $F^Z_B$, with
$\partial_tF^W_A=\sum_Bm_BD_{AB}^3W_B$ and $\partial_tF^Z_B=\sum_Am_AD_{AB}^3Z_A$.  At a label,
$Q=\Phi:=\sum_Bm_B\varphi_B$ with $\varphi_B=m_1D_{1B}+m_2D_{2B}$, and
$\partial_tF^Z_B=-\varphi_B/W_B$.
\end{lemma}

\begin{proof}
Since $W_AF^Z_A=w_AF^z_A$, $Z_AF^W_A=z_AF^w_A$, $W_BF^Z_B=w_BF^z_B$ and $Z_BF^W_B=z_BF^w_B$ for
the coordinate equations $F^z_k,F^w_k$ of \eqref{eq:system}, the first identity is the torque
identity $\sum_km_k(w_kF^z_k-z_kF^w_k)=0$, which holds by antisymmetry.  In the second sum the
terms with $D_{12}$ cancel by antisymmetry and the crossed terms give $tQ$.  The derivatives are
read off from \eqref{eq:d2-rows}, and Lemma~\ref{lem:d2-fibre}(c) gives
$D_{AB}^3Z_AW_B=-D_{AB}$ and $D_{AB}^3Z_A=-D_{AB}/W_B$.
\end{proof}

\emph{Crossed data.}  At a label fix $v_B$ with $v_B^2=1/W_B$.  By Lemma~\ref{lem:d2-fibre}(a),(c),
$D_{1B}^2=M/(m_2W_B)$ and $D_{2B}^2=-M/(m_1W_B)$, so $D_{1B}=s_B\sqrt{M/m_2}\,v_B$ and
$D_{2B}=is'_B\sqrt{M/m_1}\,v_B$ with $s_B,s'_B\in\{\pm1\}$ (positive square roots); these are the
$64$ crossed-sign assignments.  Put $\sigma_B=s_Bs'_B$,
$c_\pm=\sqrt M\,(m_1/\sqrt{m_2}\pm i\,m_2/\sqrt{m_1})$ and $x=i(m_2/m_1)^{3/2}$.  Then
$c_+c_-=M(m_1^3+m_2^3)/(m_1m_2)>0$, $\varphi_B=c_{\sigma_B}s_Bv_B$, and
$m_2D_{2B}/(m_1D_{1B})=\sigma_Bx$ with $x^2=-m_2^3/m_1^3$.  The label is \emph{same-sign} if
$\sigma_3=\sigma_4=\sigma_5$ and \emph{one-different} otherwise (two and six of the eight
patterns); then the \emph{odd body} is the $B$ whose $\sigma_B$ differs from the majority sign
$\sigma$.  In both cases put $\tilde v_B=\varphi_B/c_\sigma$, so that $\tilde v_B=s_Bv_B$ if
$\sigma_B=\sigma$ and $\tilde v_k=\rho\,s_kv_k$ for the odd body $k$, where
$\rho=c_{-\sigma}/c_\sigma$ satisfies $|\rho|=1$, $\rho\ne\pm1$.  Thus
$\Phi=c_\sigma\sum_Bm_B\tilde v_B$.
\begin{proposition}[First conditions]\label{prop:d2-first}
Let $m\in\Rp^5$ and let $p$ be the limit label of a Diagram~2 germ
\textup{(}Proposition~\ref{prop:d2-cover}\textup{)}.  Then $\Phi(p)=0$, and:
\begin{alphalist}
  \item if $p$ is same-sign, its triangle satisfies \textup{(25)}, $\Lambda=0$, and
  $L_j(m')=0$ for every $j\ne5$ with $S_j(W_3,W_4,W_5)=0$ \textup{(}$L_5L_5'=0$ if $j=5$\textup{)};
  \item if $p$ is one-different with odd body $5$ and majority sign $\sigma$, its triangle satisfies
  \textup{(28)} with $y=\sigma(x+x^{-1})/2$, and $u:=y^2=-(T-1)^2/(4T)$, $T=(m_1/m_2)^3$.  Hence
  $P_j(y)=0$ for every $j$ with $S_j(W_3,W_4,W_5)=0$: if $m_1=m_2$ then $K_j(m')=0$, and if
  $m_1\ne m_2$ then $u<0$ and $E_j(u)=O_j(u)=0$.
\end{alphalist}
\end{proposition}

\begin{proof}
Along the germ all rows vanish and $t\not\equiv0$, so Lemma~\ref{lem:d2-syzygy} gives $Q\equiv0$
along it, and $\Phi(p)=Q(p)=0$ by continuity.  No Jacobian rank is used, so this holds at every
label, including fold labels.
In case (a), $\Phi=c_\sigma\sum_Bm_Bs_Bv_B$ with $c_\sigma\ne0$, and the product of
$m_3v_3\pm m_4v_4\pm m_5v_5$ over the four sign choices is $\Lambda$, since
$(m_Bv_B)^2=f_B$.  In case (b) put $r_B=m_1m_BD_{1B}$, so $r_B^2=(m_1^2M/m_2)f_B$, and
$x'=\sigma x$; then $\Phi=(r_3+r_4)(1+x')+r_5(1-x')$, which is (26) of \cite[\S8.2]{AK}, and the
product over the four sign changes of $r_3,r_4$ is $4x'^2(m_1^2M/m_2)^2$ times the left side of
(28) with $y=(x'+1/x')/2$.
In both cases $W_4$ is a common root of $S_j^\flat$ and $\Lambda^\flat$, resp.\ $(28)^\flat$, so
the resultants of Definition~\ref{def:d2-mass} vanish, and $c_j\ne0$, $m_3\ne0$.  Finally
$y=iv$ with $v=\sigma\bigl((m_2/m_1)^{3/2}-(m_1/m_2)^{3/2}\bigr)/2\in\R$, which vanishes iff
$m_1=m_2$, and $E_j,O_j$ have real coefficients; so $P_j(iv)=E_j(-v^2)+ivO_j(-v^2)=0$ gives the
last assertion.
\end{proof}

\subsection{The isolation mechanism}\label{sec:d2-isolation}

Let $F_{\mathrm{sel}}$ denote the twenty rows other than $F^Z_2$ and $F^Z_5$.  A label $p$ is
\emph{excluded} if no holomorphic curve germ $c\colon(\C,0)\to(\mathcal X_m,p)$ has
$t\circ c\not\equiv0$; an isolated point of $\mathcal X_m$ is excluded.

\begin{lemma}[Local analytic lemma]\label{lem:d2-isolation}
Let $m\in\Rp^5$ and $p\in\mathcal B_m$.
\begin{alphalist}
  \item If $p$ is excluded, it is not the limit label of a Diagram~2 germ.
  \item Suppose $dF_{\mathrm{sel}}(p)$ has rank $20$, so that near $p$ the set
  $\{F_{\mathrm{sel}}=0\}$ is a smooth curve $C$.  Then near $p$ the set $\mathcal X_m$ consists of the
  points of $C$ at which $tQ=0$, and $p$ is excluded if $Q|_C\not\equiv0$.  If moreover
  $\partial F_{\mathrm{sel}}/\partial(Z,W,D)(p)$ is invertible, then $C$ is the graph of an
  analytic map $t\mapsto\phi(t)$, and $p$ is an isolated point of $\mathcal X_m$ as soon as
  $q(t)=Q(t,\phi(t))\not\equiv0$, in particular if $q(0)=\Phi(p)\ne0$ or $q'(0)\ne0$.
  \item Suppose some twenty of the twenty-two polynomials have a differential of rank $20$ at
  $p$, so that near $p$ their common zeros form a smooth curve $C'$.  If one of the other two
  polynomials does not vanish identically on $C'$, then $p$ is an isolated point of $\mathcal X_m$.
\end{alphalist}
\end{lemma}

\begin{proof}
(a) By Proposition~\ref{prop:d2-cover}, a Diagram~2 germ with limit $p$ gives, after the
substitution $s=\tau^N$ that clears the denominators of the Puiseux exponents, a holomorphic germ
$\tau\mapsto(t,Z,W,D)$ in $\mathcal X_m$ at $p$ with $t\not\equiv0$.
(b) On $C$ all rows except $F^Z_2,F^Z_5$ vanish and $W_5=1$, so Lemma~\ref{lem:d2-syzygy} gives
$m_2W_2F^Z_2=tQ=-m_5F^Z_5$ on $C$; and $W_2(p)=m_1D_{12}(p)\ne0$.  This describes $\mathcal X_m$
near $p$.  A germ $c$ in $\mathcal X_m$ at $p$ lies in $C$, and if $t\circ c\not\equiv0$ then
$Q\circ c\equiv0$; when
$Q|_C\not\equiv0$ its zeros on the curve $C$ are isolated, so $c$ is constant, which contradicts
$t\circ c\not\equiv0$.  In the graph case $\mathcal X_m$ is, near $p$, the set of points $(t,\phi(t))$ with
$tq(t)=0$, which is $\{p\}$ when $q\not\equiv0$.
(c) Near $p$, $\mathcal X_m$ lies in $C'$ and in the zero set of the other polynomial, whose restriction
to the smooth curve $C'$ is analytic and not identically zero, hence has an isolated zero at $p$.
\end{proof}

At a few labels the Jacobian of the twenty-two polynomials has rank $19$.  There we use the
following formal obstruction, a Newton--Puiseux argument on the blow-up of a Lyapunov--Schmidt
reduction.

\begin{lemma}[Corank-two isolation]\label{lem:d2-corank-two}
Let $m\in\Rp^5$ and $p\in\mathcal B_m$.  Let $F_R$ be nineteen of the twenty-two chart polynomials,
$F_{R'}$ the other three, and $y\ne t$ a chart coordinate such that the derivative of
$(F_R,t,y)\colon\C^{21}\to\C^{21}$ at $p$ is invertible.  By the implicit function theorem,
near $p$ the set $\{F_R=0\}$ is the image of an injective analytic map $\phi$, defined near
$0\in\C^2$, with $\phi(0)=p$ and $(y,t)\circ\phi(a,t)=(y(p)+a,t)$, and near $p$
\[
  \mathcal X_m=\phi\bigl(g^{-1}(0)\bigr),\qquad g=F_{R'}\circ\phi\colon(\C^2,0)\to(\C^3,0).
\]
Let $g^{(j)}$ be the homogeneous part of degree $j$ of $g$, and for $\xi\in\C^2$ let
$\Pi_\xi=Dg^{(2)}(\xi)$, a $3\times2$ matrix.  Suppose:
\begin{alphalist}
  \item $g^{(1)}=0$; equivalently, the Jacobian at $p$ of all twenty-two polynomials has
  rank~$19$;
  \item the components of $g^{(2)}$ have only finitely many common zeros $[\xi_1],\dots,[\xi_r]$
  in $\mathbb P^1$, and $\Pi_{\xi_j}\ne0$ for every $j$;
  \item for every $j$ there is $N\ge3$ such that the vector $r_N(\xi_j)$ defined below exists
  and does not lie in $\Pi_{\xi_j}\C^2$.
\end{alphalist}
Then no nonconstant formal arc $x\in\C[[\tau]]^2$ with $x(0)=0$ satisfies $g(x)=0$.  Hence $p$
is an isolated point of $\mathcal X_m$; in particular $p$ is excluded.  The same holds if
\textup{(a)--(c)} hold for $Kg(L\,\cdot)$, $K\in GL_3(\C)$, $L\in GL_2(\C)$, in place
of $g$.

The vectors $r_N(\xi)$ are defined for a common zero $\xi$ of $g^{(2)}$ and a fixed
$e\in\C^2\setminus\C\xi$.  As $\Pi_\xi\xi=2g^{(2)}(\xi)=0$ and $\Pi_\xi\ne0$, the vector
$\varpi=\Pi_\xi e$ spans the line $\Pi_\xi\C^2$.  Put $\beta_1=0$.  For $n\ge1$, if $\beta_n$ is
defined, let $r_{n+2}(\xi)$ be the coefficient of $\lambda^{n+2}$ in
$g(\lambda\xi+\beta_n(\lambda)e)$; if moreover $r_{n+2}(\xi)=-b_{n+1}\varpi$ with $b_{n+1}\in\C$,
put $\beta_{n+1}(\lambda)=\beta_n(\lambda)+b_{n+1}\lambda^{n+1}$.  Thus
\[
  r_3(\xi)=g^{(3)}(\xi),\qquad
  r_4(\xi)=g^{(4)}(\xi)+b_2\,Dg^{(3)}(\xi)e+b_2^2\,g^{(2)}(e).
\]
Condition (c) with $N=3$ says that the $3\times3$ matrix $\bigl(\Pi_{\xi_j}\,|\,r_3(\xi_j)\bigr)$
has rank two; with $N=4$ it says that $\bigl(\Pi_{\xi_j}\,|\,r_3(\xi_j)\bigr)$ has rank one and
$\bigl(\Pi_{\xi_j}\,|\,r_4(\xi_j)\bigr)$ has rank two.
\end{lemma}

\begin{proof}
The points of $\{F_R=0\}$ near $p$ are the points $\phi(a,t)$, and $\phi(a,t)\in\mathcal X_m$ if
and only if $g(a,t)=0$.  If $p$ were not isolated in $\mathcal X_m$, the curve selection lemma
\cite[Thm.~2.5.5]{BCR}, applied to the real algebraic set $\mathcal X_m\subset\C^{21}=\R^{42}$,
would give a nonconstant semialgebraic arc $\gamma\colon[0,\varepsilon)\to\mathcal X_m$ with
$\gamma(0)=p$.  Its
coordinates are convergent Puiseux series, so $x=(y\circ\gamma-y(p),t\circ\gamma)$ is, after a
substitution $s=\tau^{N'}$, a nonconstant arc in $\C[[\tau]]^2$ with $g(x)=0$.  So the first
assertion implies the second.  The zero set of $Kg(L\,\cdot)$ is $L^{-1}g^{-1}(0)$, which gives
the last assertion.

Let $x$ be a nonconstant formal arc with $x(0)=0$ and $g(x)=0$, and write
$x=\tau^k\xi'+O(\tau^{k+1})$ with $\xi'\ne0$.  By (a),
$g(x)=\tau^{2k}g^{(2)}(\xi')+O(\tau^{2k+1})$, so $\xi'\in\C^*\xi$ for one of the $\xi=\xi_j$
of (b).  Write $x=u\xi+\tilde\theta e$; then $\ord u=k<\ord\tilde\theta$.  Take $\lambda=u(\tau)$ as
the new parameter; then $\tau$ is a power series in $\lambda^{1/k}$ with nonzero linear term,
and $x=\lambda\xi+\theta e$ with $\theta\in\C[[\lambda^{1/k}]]$ and $\ord\theta>1$ ($\theta=0$ is
allowed).

We show by induction on $n\ge1$ that $\beta_n$ is defined, that
$g(\lambda\xi+\beta_ne)=O(\lambda^{n+2})$ and that $o_n=\ord(\theta-\beta_n)>n$.  For $n=N-2$
this contradicts (c).  For $n=1$ the claim holds by (a), $g^{(2)}(\xi)=0$ and $\ord\theta>1$.
Taylor's formula in the direction $e$ gives
\[
  0=g(\lambda\xi+\theta e)
   =g(\lambda\xi+\beta_ne)+(\theta-\beta_n)\bigl(\lambda\varpi+O(\lambda^2)\bigr)
    +O(\lambda^{2o_n}),
\]
because $Dg(\lambda\xi+\beta_ne)e=\lambda\Pi_\xi e+\beta_n\Pi_ee+O(\lambda^2)
=\lambda\varpi+O(\lambda^2)$ by (a) and $\ord\beta_n\ge2$.  If $o_n<n+1$, the term
$(\theta-\beta_n)\lambda\varpi$ has order $o_n+1<n+2$, every other term has larger order
($2o_n>o_n+1$), and its leading coefficient is a nonzero multiple of $\varpi\ne0$; this is
impossible.  Hence $o_n\ge n+1$, and the coefficient of $\lambda^{n+2}$ gives
$r_{n+2}(\xi)+b\varpi=0$, where $b$ is the coefficient of $\lambda^{n+1}$ in $\theta$.  So
$\beta_{n+1}$ is defined, $b_{n+1}=b$ since $\varpi\ne0$, and $o_{n+1}>n+1$.  The same expansion
with $\beta_{n+1}$ in place of $\theta$ gives
$g(\lambda\xi+\beta_{n+1}e)=(r_{n+2}(\xi)+b_{n+1}\varpi)\lambda^{n+2}+O(\lambda^{n+3})
=O(\lambda^{n+3})$.
\end{proof}

\begin{remark}\label{rem:d2-corank-two}
The $\xi_j$ are the points where the strict transform of $g^{-1}(0)$ could meet the exceptional
curve of the blow-up of $0\in\C^2$, and the induction is the Newton--Puiseux algorithm in the
blow-up chart at $\xi_j$.  Arcs inside the boundary fibre $\{t=0\}$ need no separate argument:
they are tangent to the line $t=0$, which is then one of the $\xi_j$, and the same induction
excludes them.
\end{remark}

\emph{The triangle block and the form $Q_T$.}  For a label $p$ let $J_T(p)$ be the Jacobian at
$p$ of the nine rows $F^Z_3,F^W_3,F^Z_4,F^W_4,F^W_5,T_{34},T_{35},T_{45},G_2$ with respect to the
nine triangle coordinates $Z_B,W_B,D_{BB'}$.  If $J_T(p)$ is invertible and $\xi\in\C^3$, let
$(\zeta,\eta,\vartheta)$ solve $J_T(p)(\zeta,\eta,\vartheta)=r(\xi)$, where $r(\xi)$ has the
entries $\xi_3/W_3$ and $\xi_4/W_4$ in the rows $F^Z_3$, $F^Z_4$ and $0$ elsewhere, and put
\[
  Q_T(\xi)=-\tfrac12\textstyle\sum_Bm_B\,\xi_B\,\eta_B/W_B .
\]
$Q_T$ is a quadratic form determined by $m'$ and the triangle of $p$.

\begin{lemma}[Second-order coefficient]\label{lem:d2-second}
Let $p$ be a label with $\Phi(p)=0$ and $J_T(p)$ invertible.  Then
$\partial F_{\mathrm{sel}}/\partial(Z,W,D)(p)$ is invertible and, in the notation of
Lemma~\ref{lem:d2-isolation} and of the crossed data,
\[
  q(0)=0,\qquad q'(0)=c_\sigma^2\,Q_T(\tilde v).
\]
\end{lemma}

\begin{proof}
At $t=0$ the pair rows involve only $Z_A,W_A,D_{12}$, the triangle rows only the triangle
coordinates, and $T_{AB}$ involves no crossed coordinate other than $D_{AB}$, with
$\partial T_{AB}/\partial D_{AB}=-2D_{AB}Z_AW_B=2/D_{AB}$; so the Jacobian is block lower
triangular.  The pair block (rows $F^Z_1,F^W_1,F^W_2,T_{12},G_1$) is invertible for positive
masses (computer-assisted).  Differentiate $F_{\mathrm{sel}}(t,\phi(t))=0$ at $t=0$, using
Lemma~\ref{lem:d2-syzygy}.  The pair rows are forced by
$\partial_tF^W_A=-Z_A^{-1}\sum_Bm_BD_{AB}$, and
$\partial_tF^W_1-\partial_tF^W_2=M\Phi/(m_1m_2)=0$; for equal forcings the solution has
$Z_A'(0)=D_{12}'(0)=0$ (computer-assisted).  The triangle rows give
$J_T(p)(Z',W',D')_T(0)=r(\varphi)$, hence $W_B'(0)=c_\sigma\eta_B(\tilde v)$.  On $C$ the ties
$T_{AB}=0$ give $Q=\sum_{A,B}m_Am_BD_{AB}+O(t^2)$ and
$D_{AB}'=-\frac12D_{AB}(Z_A'/Z_A+W_B'/W_B)$.  Therefore
$q'(0)=\sum_{A,B}m_Am_BD'_{AB}(0)=-\frac12\sum_Bm_B\varphi_BW_B'(0)/W_B=c_\sigma^2Q_T(\tilde v)$.
\end{proof}

The two computer-assisted steps are exact computations with the masses as indeterminates, at
the pair part of Lemma~\ref{lem:d2-fibre}(a) and modulo $MD_{12}^3=1$: the determinant of the pair
block is a power of $D_{12}$ times factors of fixed sign for positive masses, and for equal
forcings the pair rows are solved with $Z_A'(0)=D_{12}'(0)=0$ (Appendix~\ref{app:verification}).

\begin{definition}\label{def:d2-psi}
For a label $p$ with $J_T(p)$ invertible put $\Psi_T(p)=Q_T(\tilde v)$ if $p$ is same-sign and
$\Psi_1(p)=Q_T(\tilde v)$ if $p$ is one-different.  The label is \emph{isolated at second order}
if $\Phi(p)=0$, $J_T(p)$ is invertible and $Q_T(\tilde v)\ne0$; by
Lemmas~\ref{lem:d2-isolation} and~\ref{lem:d2-second} it is then an isolated point of $\mathcal X_m$, and
the omitted rows $F^Z_2,F^Z_5$ vanish to order exactly two along the curve $C$.
\end{definition}

At a label with $\Phi(p)=0$ the vector $\tilde v$ is determined by the triangle and the signs
$s_B$ (for an odd body $k$, $m_k\tilde v_k=-\sum_{B\ne k}m_B\tilde v_B$), so $\Psi_T$ and
$\Psi_1$ do not involve $m_1,m_2$.  For a one-different label with odd body $5$ the condition
$\Phi=0$ requires $|m_3\tilde v_3+m_4\tilde v_4|=m_5|v_5|$, i.e.\ the vanishing of the
\emph{sheet function} $G=|m_3\tilde v_3+m_4\tilde v_4|^2-m_5^2$ (recall $W_5=1$), which again
depends on the triangle only; then $\tilde v_5=\rho s_5v_5$ fixes $\rho$ up to sign, hence
$m_1/m_2$ up to exchanging $m_1$ and $m_2$.

\subsection{The case analysis}\label{sec:d2-cases}

By Proposition~\ref{prop:d2-first} and Lemma~\ref{lem:d2-isolation}(a) it remains to show that
every label $p$ with $\Phi(p)=0$ at positive masses is excluded.  Table~\ref{tab:d2-cases} lists
the cases; the computer-assisted statements follow.  Ball arithmetic is IEEE arithmetic with
proved rounding bounds (Appendix~\ref{app:verification}).

\begin{table}[ht]
\centering
\footnotesize
\setlength{\tabcolsep}{3pt}
\begin{tabular}{p{2.5cm}p{2.5cm}p{3.3cm}p{3.4cm}}
\hline
Labels & Mass locus & Method & Computation \\
\hline
$\Phi\ne0$ (any) & all $m$ & first order, $q(0)=\Phi$ & exact identities \\
same-sign, $S_4$, or $S_5$ with $L_5'=0$ & $L_4=0$, $L_5'=0$ & no point in $\Rp^3$ & exact \\
same-sign, $S_5$ & $L_5=0$: $m_3=m_4=m_5$ & order three; corank-two blow-up & exact jets \\
same-sign, $S_1$--$S_3$, $S_6$--$S_8$ & rays of $L_1$--$L_3$, $L_6$--$L_8$ & second order,
$|\Psi_T|>0.3$ & Bernstein subdivision, Krawczyk method; interval arithmetic \\
\hline
one-different, $m_1=m_2$, $S_3,S_6$--$S_8$ & $K_3,K_6,K_7,K_8=0$ & no point in $\Rp^3$ &
Bernstein coefficients \\
one-different, $m_1=m_2$, $S_1,S_2$ & $K_1=0$, $K_2=0$ & second order & exact; interval cover \\
one-different, $m_1=m_2$, $S_4$ & $K_4=0$: four arcs & second order & exact \\
one-different, $m_1=m_2$, $S_5$ & $K_5=0$: $\Rp(1,1,\rho_*)$ & orders one to four &
exact; interval arithmetic \\
\hline
one-different, $m_1\ne m_2$, not fold & $E_j(u)=O_j(u)=0$ & second order, $\Psi_1\ne0$: ball
covers, exact lemmas & ball and interval arithmetic; exact \\
one-different, $m_1\ne m_2$, fold & $E_5(u)=O_5(u)=0$ & corank one; $Q|_C$ has a simple zero &
exact \\
\hline
\end{tabular}
\caption{The Diagram~2 cases.  A label with $\Phi=0$ lies on $S_j$ for the $j$ given in the
first column (Lemma~\ref{lem:d2-fibre}); one-different labels are taken with odd body~$5$.}
\label{tab:d2-cases}
\end{table}

\begin{lemma}[computer-assisted]\label{lem:d2-empty}
\textup{(a)} The polynomials $L_4$, $L_5'$, $K_3$, $K_6$, $K_7$ and $K_8$ have no zero in $\Rp^3$.
More precisely, at positive masses every $S_4$ triangle is real \textup{(}the centre-of-mass
section of $S_4$ is a cubic with positive discriminant\textup{)} and violates \textup{(25)}.
\textup{(b)} If $m_3=m_4=m_5$, no $S_j$ with $j\ne5$ has a root satisfying \textup{(25)}.
\end{lemma}
The proof is an exact computation over $\mathbb Q$.  For $L_4$ we use a rational parametrization
of $S_4$; on a real $S_4$ triangle at positive masses the positions $W_B$ do not all have the same
sign, and then $\Lambda$ is a sum of nonnegative terms, one of them positive.  $L_5'$ has
positive coefficients, and $K_3$, $K_6$, $K_7$, $K_8$ have nonnegative Bernstein coefficients on
twelve projective charts, three for each of them, each split into the two sectors of a blow-up
at a corner of the cone.  For (b) we check that $S_j^\flat$ and $\Lambda^\flat$ are coprime at
$m_3=m_4=m_5$ for $j\ne5$ (Appendix~\ref{app:verification}).

\begin{lemma}[computer-assisted]\label{lem:d2-L5}
Let $m\in\Rp^5$ with $m_3=m_4=m_5$.  Every label at $m$ whose triangle lies on $S_5$ \textup{(}an
all-equal Lagrange class or a fold triangle\textup{)} and which satisfies $\Phi=0$ or is a fold
label is an isolated point of $\mathcal X_m$.  At each such $m$ these are $72$ labels of the all-equal
classes and $1152$ fold labels, $1224$ in all.
\end{lemma}

\emph{Masses.}  Common mass scaling is an exact isomorphism of the gauged chart, so we may take
$m_3=m_4=m_5=1$; then $\delta_*:=M_T^{-1/3}=3^{-1/3}$.  Every positive pair $(m_1,m_2)$ is, for
exactly one $(\nu,M)\in(0,1)\times\Rp$,
\[
  m_1=\frac{4M\nu^2}{(1+\nu^2)^2},\qquad m_2=\frac{M(1-\nu^2)^2}{(1+\nu^2)^2},
\]
namely $M=m_1+m_2$ and $\nu=1/(\chi+\sqrt{1+\chi^2})$, where
$\chi=\sqrt{m_2/m_1}=(1-\nu^2)/(2\nu)$.

\emph{The all-equal labels.}  The $72$ labels of the all-equal classes have $\Phi=0$, an
invertible $J_T(p)$ and $q(0)=q'(0)=0$, and along $C$ we have $F^Z_5=\kappa_3t^3+O(t^4)$, so
$q''(0)=-2\kappa_3$.  For the nine labels of one sign class with $D_{BB'}=\omega^k\delta_*$,
$(W_3,W_4,W_5)=(\omega,\omega^2,1)$ and $D_{12}=\omega^jM^{-1/3}$,
\[
  \kappa_3=\omega^jM^{5/3}R_1(1-i\chi)+3\omega^k\delta_*MR_2(1-i\chi^5),\qquad
  R_1=\frac{18\nu^2(1-\nu^2)}{(1+\nu^2)^3},\quad R_2=\frac{6\nu^2(1+\nu^2)}{(1-\nu^2)^3}.
\]
These nine are the images of the case $j=k=0$ under the automorphisms
$\delta_*\mapsto\omega^k\delta_*$, $M^{1/3}\mapsto\omega^{-j}M^{1/3}$ of the coefficient field,
which fix the masses and the polynomials; the other $63$ labels are images of these nine under
the exchange of bodies $3$ and $4$, conjugation and the reversal of $t$.  Here $R_1,R_2>0$, and
$(1-i\chi)/(1-i\chi^5)$ is a positive multiple of $1+\chi^6+i(\chi^5-\chi)$, where
$|\chi^5-\chi|\le\chi+\chi^5\le1+\chi^6$ because $(1+\chi^6)-(\chi+\chi^5)=(1-\chi)(1-\chi^5)\ge0$.
So the quotient of the two terms of $\kappa_3$ is $\omega^{j-k}$ times a number of argument at
most $\pi/4$ in absolute value, hence is not $-1$.  Therefore $\kappa_3\ne0$, and each of these
labels is isolated by Lemma~\ref{lem:d2-isolation}(b).

\emph{Fold labels of rank $20$.}  At $1080$ fold labels the Jacobian of the twenty-two
polynomials has rank $20$.  Twenty rows together with $h=W_4-W_4(p)$ have an invertible
derivative: the rows of $F_{\mathrm{sel}}$ at $864$ labels, and all rows other than $F^Z_2$ and
$F^Z_3$ at the other $216$.  An omitted row has a nonzero coefficient of $h^3$ along the curve
these rows define, so each of these labels is isolated by Lemma~\ref{lem:d2-isolation}(c).

\emph{Fold labels of rank $19$.}  The other $72$ fold labels form one orbit under the maps
listed below.  Its seed $p_*=p_*(\nu,M)$ has the pair part of Lemma~\ref{lem:d2-fibre}(a) with
$D_{12}=M^{-1/3}$ and the fold triangle
\[
  (W_3,W_4,W_5)=(\omega,\omega^2,1),\qquad Z_B=-\delta_*W_B,\qquad
  (D_{34},D_{35},D_{45})=\delta_*(1,\omega,\omega^2).
\]
Apply Lemma~\ref{lem:d2-corank-two} with $F_{R'}=(F^W_2,F^W_4,F^W_5)$, $F_R$ the other nineteen
rows and $y=Z_3$.  The kernel coordinates are then the chart parameter $t$ and
$a=Z_3-Z_3(p_*)$.  At $t=0$ the derivative of $(F_R,t,Z_3)$ is block lower triangular, as in the
proof of Lemma~\ref{lem:d2-second}.  Its diagonal blocks are $(1)$ for $dt$; the triangle block
(rows $F^Z_3,F^W_3,F^Z_4,F^Z_5,T_{34},T_{35},T_{45},G_2,dZ_3$, columns $Z_B,W_B,D_{BB'}$), with
determinant $18\delta_*$; the pair block (rows $F^Z_1,F^W_1,F^Z_2,T_{12},G_1$, columns
$Z_A,W_A,D_{12}$), with determinant $-3M^{-1/3}$; and the diagonal crossed block $(2/D_{AB})$,
with determinant $-512i\,\nu^3(1-\nu^2)^3/(1+\nu^2)^6$.  So the invertibility hypothesis holds
for all $(\nu,M)$.  The differentials of the three omitted rows vanish on the two-dimensional
kernel of $dF_R(p_*)$, so (a) holds.  Write $g=(\psi_2,\psi_4,\psi_5)$ with
$\psi_k=F^W_k\circ\phi$.

At $M=1$ the jet of $g$ through degree four was computed exactly over
$\mathbb Q(\nu)(\sqrt3,i,\delta_*)$, and the nineteen rows of $F_R$ vanish through degree four on
the computed jet of $\phi$.  In the quadratic part, $\psi_2^{(2)}=0$, there is no $a^2$ term,
and the $3\times3$ coefficient matrix has rank one.  So $g^{(2)}=\psi_4^{(2)}\,n$ for a fixed
vector $n=(0,1,n_5)$, where $\psi_4^{(2)}=t(\gamma_{11}a+\gamma_{02}t)$ and
\[
  \gamma_{11}=\frac{9i\delta_*(\nu-i)(\nu+i)^3\,\Upsilon(\nu)}{4\nu(\nu^2-1)(\nu^2+1)^2},\qquad
  \Upsilon(\nu)=\nu^4-2i\nu^3-6\nu^2+2i\nu+1=\tfrac14\bigl(3(\nu-i)^4+(\nu+i)^4\bigr).
\]
For real $\nu$, $|3(\nu-i)^4|=3(\nu^2+1)^2>|(\nu+i)^4|$, so $\Upsilon(\nu)\ne0$ and
$\gamma_{11}\ne0$ for $0<\nu<1$.  Hence the common zeros of $g^{(2)}$ are two lines: the
\emph{boundary line} $\xi_\partial=(1,0)$, tangent to $\{t=0\}$, and the \emph{slanted line}
$\xi_s=(-\gamma_{02},\gamma_{11})$.  Moreover $\Pi_\xi=n\,d\psi_4^{(2)}(\xi)$, with
$d\psi_4^{(2)}(\xi_\partial)=(0,\gamma_{11})$ and
$d\psi_4^{(2)}(\xi_s)=\gamma_{11}(\gamma_{11},\gamma_{02})$, both nonzero.  So (b) holds.
On the slanted line take $N=3$: the matrix $(\Pi_{\xi_s}\,|\,r_3(\xi_s))$ has the $2\times2$
minor
\[
  \mu_3=-\frac{3^{10}(\nu-i)^4(\nu+i)^{16}\,\Upsilon(\nu)^5}
             {2^{12}\delta_*\,\nu^6(\nu^2-1)^6(\nu^2+1)^8}\ne0 .
\]
On the boundary line take $e=(0,1)$, so that $\varpi=\gamma_{11}n$, and $N=4$.  Identically in
$\nu$, $r_3(\xi_\partial)=g^{(3)}(1,0)=-b_2\varpi$, and the matrix
$(\Pi_{\xi_\partial}\,|\,r_4(\xi_\partial))$ has the minor
\[
  \mu_4=-\frac{27i\delta_*(\nu-i)^3(\nu+i)^5\,\Upsilon(\nu)}{64\,\nu^3(\nu^2-1)^3}\ne0 .
\]
So (c) holds.  Every coefficient used is regular for $0<\nu<1$: the jets are polynomials in the
coordinates of $p_*$, the masses and the entries of the inverses of the three diagonal blocks,
and the only further division is by $\gamma_{11}$.  Hence the identities above hold at every
$\nu\in(0,1)$.

For general $M$, $p_*(\nu,M)$ is obtained from $p_*(\nu,1)$ by multiplying $W_1,W_2$ by $M^{2/3}$
and $D_{12}$ by $M^{-1/3}$.  The map
$\tilde g(a,\tilde t)=\bigl(M^{-2/3}\psi_2,\psi_4,\psi_5\bigr)(a,\tilde t/M)$ was computed in
exact Laurent arithmetic in $M^{1/3}$, and the selected rows vanish along both lines through the
orders used.  The quadratic part of $\tilde g$ does not depend on $M$, the compatibility
$r_3(\xi_\partial)\in\C\varpi$ persists, and the two minors become $M^{-3}\mu_3(\nu)$ and
$M^{-5/3}\mu_4(\nu)$.  By Lemma~\ref{lem:d2-corank-two}, applied to $\tilde g$, $p_*(\nu,M)$ is
an isolated point of $\mathcal X_m$ for every $0<\nu<1$ and $M>0$.

\emph{Transport.}  Two involutions preserve the twenty-two polynomials, since these have real
coefficients and $m_3=m_4$: the reversal $(t,D_{AB})\mapsto(-t,-D_{AB})$, which is biholomorphic,
and complex conjugation followed by the exchange of bodies $3$ and $4$, which is
antiholomorphic.  They commute, fix $(\nu,M)$, and carry $p_*$ to the labels with the other three
signs of the orbit.  Multiplying $t$, $Z_B$, $W_A$, $D_{12}$ and $D_{BB'}$ by $\omega^j$
multiplies each polynomial by a cube root of unity and fixes the gauges.  A permutation of the
bodies $3,4,5$, followed by the element of the chart torus that restores the gauges, is a local
biholomorphism of $\mathcal X_m$ by the implicit function theorem, because the transversality
determinant is $2(Z_2-Z_1)W_5=2$.  These $3\cdot6$ maps give eighteen labels for each of the four
signs.  The resulting $72$ labels are exactly one orbit of the discrete label action and are
disjoint from the $1080$ labels above.  All these maps preserve isolatedness.

The computations in the proof of Lemma~\ref{lem:d2-L5}, namely the coefficients $\kappa_3$, the
ranks and the coefficients of $h^3$ at the $1080$ fold labels, the determinants, jets and minors
at $p_*$, and the orbit of $p_*$, are exact, over number fields containing $\omega$, $i$ and
$\delta_*$ and over rational function fields over them in $\nu$ and $M^{1/3}$
(Appendix~\ref{app:verification}).

\begin{lemma}[computer-assisted]\label{lem:d2-rays}
\textup{(a)} The zero set of $L_8$ in $\Rp^3$ is $\Rp(p_*,q_*,1)\cup\Rp(q_*,p_*,1)$ with
$p_*=0.67987758\ldots$ and $q_*=0.82343785\ldots$.
\textup{(b)} The zero set of $L_1$ in $\Rp^3$ is $\Rp(a_*,b_*,1)\cup\Rp(a_*,1,b_*)$ with
$a_*=1.02227400\ldots$ and $b_*=0.90559635\ldots$; at $(a_*,b_*,1)$ the zeros of $S_1$ that
satisfy \textup{(25)} form exactly one pair of complex conjugate non-real shapes.
\textup{(c)} Consequently the zero sets of $L_1L_2L_3$ and of $L_6L_7L_8$ in $\Rp^3$ are the rays
through the six permutations of $(a_*,b_*,1)$ and of $(1,p_*,q_*)$, respectively.
\end{lemma}
For (a) we combine exact Bernstein subdivision of the mass simplex with the Krawczyk method in
rational interval arithmetic.  For (b) we eliminate the shape by subresultants, locate the zeros
by Bernstein subdivision and the signs of resultants, enclose them by the Krawczyk method, and
show that the quadratic subresultant has negative discriminant there
(Appendix~\ref{app:verification}).  Part (c) follows from (a), (b) and the permutation statement
after Definition~\ref{def:d2-mass}.

\begin{lemma}[computer-assisted]\label{lem:d2-Lray}
Let $m\in\Rp^5$ with $m'$ on one of the rays of Lemma~\ref{lem:d2-rays}(c), and let $p$ be a
same-sign label at $m$ with $\Phi(p)=0$.  Then $J_T(p)$ is invertible and $\Psi_T(p)\ne0$, so $p$
is isolated at second order.
\end{lemma}

At each permutation of $(1,p_*,q_*)$ the $48$ Lagrange triangles, and at each permutation of
$(a_*,b_*,1)$ the six Euler triangles on the conjugate shape pair, are enumerated; there only the
ray factor itself vanishes among $L_1,\dots,L_8,L_5'$, so by Proposition~\ref{prop:d2-first}(a)
no other triangle carries a same-sign label with $\Phi=0$.  For $36+36$ triangles a sign vector
with $\sum_Bm_Bs_Bv_B=0$ exists, unique up to sign; $J_T$ is invertible, and complex interval
arithmetic gives $|\Psi_T|>0.36$ ($L_6$--$L_8$) and $|\Psi_T|>0.3$ ($L_1$--$L_3$).  Common mass
scaling covers the rays, and $\Psi_T$ does not involve $m_1,m_2$.
The shapes and the square and cube roots are enclosed by the Krawczyk method, and the solutions
of the linear systems are enclosed rigorously (Appendix~\ref{app:verification}).

\begin{lemma}[computer-assisted]\label{lem:d2-K}
Let $m\in\Rp^5$ with $m_1=m_2$, let $p$ be a one-different label at $m$ with odd body $5$ and
$\Phi(p)=0$, and let $S_j$ carry its triangle, so that $K_j(m')=0$.  Then $p$ is an isolated
point of $\mathcal X_m$.  More precisely:
\begin{alphalist}
  \item $j\notin\{3,6,7,8\}$, by Lemma~\ref{lem:d2-empty};
  \item if $j\in\{1,2\}$: at every zero of $K_1$ in $\Rp^3$ the $S_1$ shapes compatible with
  \textup{(28)} at $y=0$ reduce to one simple common root, which is real and lies on one of two
  explicit real branches; $J_T(p)$ is invertible and $q'(0)\ne0$.  The case $j=2$ is the image of
  $j=1$ under the exchange of bodies $3$ and $4$;
  \item if $j=4$: $m'$ lies on one of four arcs, $k\mapsto\mathrm{II}(k),\mathrm{III}(k)$ for
  $5^{-1/4}<k<1$ and their images under $m_3\leftrightarrow m_4$, where
  $\alpha=4k(1+k^4)$, $\beta=(1-k^2)(1+6k^2+k^4)$, $\gamma=(1+k^2)(6k^2-1-k^4)$,
  $\mathrm{II}=((\gamma^3-\beta^3)\alpha,(\alpha^3-\gamma^3)\beta,(\alpha^3-\beta^3)\gamma)$ and
  $\mathrm{III}=((\gamma^3-\beta^3)\alpha,(\gamma^3+\alpha^3)\beta,(\alpha^3+\beta^3)\gamma)$;
  $J_T(p)$ is invertible and $q'(0)\ne0$;
  \item if $j=5$: $m'\in\Rp(1,1,\rho_*)$, where $\rho_*\in(1.5736409,1.573641)$ is the unique
  positive root of $\rho^8+2\rho^7+3\rho^6-4\rho^5-5\rho^4-6\rho^3-3\rho^2-12\rho-12$, and every
  label at $m$ with a Lagrange triangle is isolated, for all $27$ classes and all $64$
  crossed-sign assignments.
\end{alphalist}
\end{lemma}

In (b) and (c), $60$ of the $64$ crossed-sign assignments have $\Phi\ne0$, and for the other four
(one orbit of conjugation and $t\mapsto-t$) $q'(0)\ne0$.  In (d), the two-equal classes have
$\Phi\ne0$; for the all-equal classes $60$ assignments have $\Phi\ne0$ and $4$ have $q'(0)\ne0$;
for the fold classes an augmented system in $h=W_4-W_4(p)$ gives a nonzero coefficient of order
three ($60$ assignments) or four ($4$).  Exact grading makes all items uniform in $m_1=m_2$.
In (b) the common root is obtained from a linear subresultant over $\mathbb Q$, using the
smoothness of the zero set of $K_1$ in $\Rp^3$, which is proved with Bernstein coefficients; the
two real branches are found from a rational parametrization of $S_1$ and Sturm sequences, and
$q'(0)\ne0$ follows from an exact reduction and an interval cover of the branch parameter.  In (c)
the computation is exact, on the rational parametrizations of the four arcs, and classifies the
real $S_4$ triangles; by Lemma~\ref{lem:d2-empty}(a) every $S_4$ triangle at positive masses is
real.  In (d) the zero set of $K_5$ in $\Rp^3$ is determined by exact resultants and a Sturm
count, and the labels are classified in complex rational interval arithmetic, at one triangle in
each orbit of conjugation and the exchange of bodies $3$ and $4$
(Appendix~\ref{app:verification}).  Renumbering the bodies as at the beginning of this section
gives the same conclusions for all pair splits and odd-body placements.

\begin{lemma}[computer-assisted]\label{lem:d2-P}
Let $m\in\Rp^5$ with $m_1\ne m_2$, and let $p$ be a one-different label at $m$ with odd body $5$,
$\Phi(p)=0$, whose triangle is not a fold triangle.  Then $J_T(p)$ is invertible and
$\Psi_1(p)\ne0$, so $p$ is isolated at second order.
\end{lemma}

We cover the simplex $m_3+m_4+m_5=1$ for $S_1,S_3,S_4$ (all Euler shapes, enclosed
by Krawczyk or Pellet discs) and $S_5,S_6,S_8$ (all Lagrange shapes, up to conjugation and the
cube-root action), for both relative signs $s_3s_4$, by boxes on each of which $G$ has a strict
sign, or $\Psi_1$ is enclosed away from $0$ in ball arithmetic, or, on boxes meeting $m_3=0$ or
$m_4=0$ (where $\Psi_1$ vanishes), its transversal derivative is.  The sets left out of the cover
are treated separately: the vertex corners ($G\ne0$); the ends $(m_3,m_4)=(0,\frac12)$ of $S_1$
and $S_4$ and $(\frac12,0)$ of $S_4$, where a majority body becomes massless (an exact edge
identity and interval margins $0.787$, $0.588$, $0.551$); the corner $(\frac12,\frac12,0)$ of
$S_3$ (an exact weighted blow-up of weight two in $m_5$ of the Euler shape parameter, after which
the determinant and the numerator of the response are $m_5^{13}$ and $m_5^6$ times factors bounded
below in modulus by $6.96\cdot10^8$ and $1.63\cdot10^6$ in interval arithmetic); six real points
of $S_4$ where the Euler label degenerates, and the corner $(\frac12,\frac12,0)$ of $S_4$ (an
exact closed form of $\Psi_1$ on the Euler shapes, in which the degenerating factor cancels); and
$S_5$ by exact algebra (Appendix~\ref{app:verification}).  On $S_5$ let $g$ be the
$W_4$-component of the response of the triangle block (row $F^Z_5$ omitted, $dW_5=0$) to a unit
forcing of the row $F^Z_4$, and $X=\tilde v_4/W_4-\tilde v_3/W_3$.  Then
$\Psi_1=-\frac12m_4gX^2$; the numerator of $g$ is a nonzero constant times
$m_3(m_3^2+m_3m_4+m_4^2)^2(P_0+\omega P_1)$ with real polynomials $P_0,P_1$ whose only common zero
in the simplex is $m_3=m_4=\frac13$, and $X\ne0$ except at the same point.  There $m_3=m_4=m_5$,
and the label with $G=0$ has $\rho=-1$, which is impossible.  $S_2$ and $S_7$ follow from $S_1$
and $S_6$ by exchanging bodies $3$ and $4$.

\begin{lemma}[computer-assisted]\label{lem:d2-fold}
Let $m\in\Rp^5$ and let $p$ be a one-different label at $m$ with odd body $5$, $\Phi(p)=0$ and a
fold triangle.  Then $dF_{\mathrm{sel}}(p)$ has rank $20$ and $Q|_C$ has a simple zero at $p$; in
particular $p$ is excluded.
\end{lemma}

$J_T(p)$ has corank exactly one, with an explicit kernel vector and left null vector $\ell$.
Since the $t$-column of the triangle rows is $-\varphi_B/W_B$ on the $Z$-rows, $\ell$ applied to it
is $c_\sigma X$ with $X$ as in Lemma~\ref{lem:d2-P}, and $X\ne0$ at every label considered (the
exceptional point $m_3=m_4=m_5$ carries no one-different label with $\Phi=0$ at positive masses).
So $C$ exists, is parametrized by $h=W_4-W_4(p)$, has $dt/dh=0$ at $p$, and
$dQ/dh=-\frac12c_\sigma m_4X\ne0$ at $p$.
The proof is an exact computation over $\mathbb Q(m_3,m_4)(\omega)$, with $m_3,m_4$ as
indeterminates: it gives the kernel vector, the left null vector, a minor of order eight of
$J_T(p)$ that does not vanish on the simplex, and the zero set of $X$
(Appendix~\ref{app:verification}).

\begin{theorem}\label{thm:d2}
No Diagram~2 germ exists at positive masses: for $m\in\Rp^5$ no curve in $A_m$ has a Diagram~2
germ.  In particular Theorem~\ref{thm:exclusion}\textup{(a)} holds for $r=2$.
\end{theorem}

\begin{proof}
Let $\gamma$ be a Diagram~2 germ on a curve in $A_m$ and $p$ its limit label
(Proposition~\ref{prop:d2-cover}).  For each $B$, $\sigma_Bx$ is the limit of
$m_2\delta_{2B}/(m_1\delta_{1B})$ along $\gamma$, which does not depend on the chart; so the
crossed pattern of $p$ is determined by $\gamma$, and renumbering the bodies $3,4,5$ before
applying Proposition~\ref{prop:d2-cover} permutes it.  We may therefore assume that $p$ is
same-sign or has odd body $5$.
By Lemma~\ref{lem:d2-isolation}(a) it suffices to show that $p$ is excluded.  By
Proposition~\ref{prop:d2-first}, $\Phi(p)=0$; let $S_j$ carry the triangle of $p$
(Lemma~\ref{lem:d2-fibre}).

\emph{Same-sign.}  By Proposition~\ref{prop:d2-first}(a), $L_j(m')=0$, or $L_5L_5'=0$ if
$j=5$.  Lemma~\ref{lem:d2-empty} excludes $j=4$ and $L_5'=0$.  If $L_5=0$, then $m_3=m_4=m_5$,
$j=5$ by Lemma~\ref{lem:d2-empty}(b), and Lemma~\ref{lem:d2-L5} applies.  Otherwise
$j\in\{1,2,3,6,7,8\}$, $m'$ lies on a ray of Lemma~\ref{lem:d2-rays}(c), and
Lemma~\ref{lem:d2-Lray} applies.

\emph{One-different.}  If $m_1=m_2$, Lemma~\ref{lem:d2-K} applies.  If $m_1\ne m_2$,
Lemma~\ref{lem:d2-fold} applies when the triangle is a fold triangle and Lemma~\ref{lem:d2-P}
otherwise.

In every case $p$ is excluded (by Definition~\ref{def:d2-psi} for the labels isolated at second
order), which contradicts Lemma~\ref{lem:d2-isolation}(a).
\end{proof}

\begin{remark}\label{rem:d2-scope}
(i) Theorem~\ref{thm:d2} is stronger than Theorem~\ref{thm:exclusion}(a) requires: it excludes
every Diagram~2 germ, with or without an unbounded four-product, and uses no reality.  From
\cite{AK} it uses only Fact~\ref{fact:rules} (in Proposition~\ref{prop:d2-cover}), besides
Lemma~\ref{lem:germ}; the mass conditions (25)--(28) of \cite[\S8.2]{AK} are derived
(Proposition~\ref{prop:d2-first}), not assumed.

(ii) One can show that the layer $u=\infty$ of the $P_j$ misses the positive cone, and that no
non-degenerate triple of positions with positive centre-of-mass masses lies on two curves
$S_i\ne S_j$ (Appendix~\ref{app:verification}).  We do not use this.
\end{remark}



\section{Diagram 3}\label{sec:d3}

Diagram~3 of \cite[Fig.~11]{AK} is the diagram of Roberts' continuum at infinity
\cite[\S7.2]{AK}.  With the numbering of \cite[\S5.2 and \S8.3]{AK}, where the clustering scheme
is $z:12\ldots34$, $w:14\ldots23$, it is the quadrilateral $1234$ with $z$-edges $14$, $23$ and
$w$-edges $12$, $34$; the diagonals and body $5$ carry no stroke.  Exchanging $z$ and $w$ gives
the same diagram with the bodies $2$ and $4$ exchanged, so after renumbering a Diagram-3 germ has
exactly this diagram.  Unlike \cite[proof of Thm.~6(iii)]{AK} we take no second singular limit.

\subsection{The crossed-cluster chart}
The equations of \eqref{eq:system} other than the gauge are invariant under
$(z,w,\delta)\mapsto(\mu z,\mu^{-1}w,\delta)$, $\mu\in\C^*$, and here we use the gauge
$z_{13}=w_{13}$, as \cite[p.~543]{AK} does for Roberts' continuum.  The chart has the parameter
$t$ and the coordinates $A,B,x,y,\hat z_5,C,D,u,v,\hat w_5$ and $D_{kl}$:
\begin{equation}\label{eq:d3-chart}
\begin{gathered}
z_1=\tfrac{A}{t},\quad z_2=\tfrac{A}{t}+tx,\quad z_3=\tfrac{B}{t},\quad z_4=\tfrac{B}{t}+ty,
\quad z_5=t^2\hat z_5,\\
w_1=\tfrac{C}{t},\quad w_2=\tfrac{D}{t},\quad w_3=\tfrac{D}{t}+tv,\quad w_4=\tfrac{C}{t}+tu,
\quad w_5=t^2\hat w_5,\\
\delta_{kl}=D_{kl}\ \ (kl=12,14,23,34),\qquad
\delta_{kl}=t\,D_{kl}\ \ (kl=13,24,15,25,35,45).
\end{gathered}
\end{equation}
Multiplying the $z$- and $w$-equations of the bodies $1,\dots,4$ and the gauge equation by $t$,
and the equations of body $5$ by $t^{-2}$, gives $21$ polynomial equations
$F=F^0+t^2F^2+\dots+t^6F^6=0$, with no term linear in $t$.  The \emph{boundary equations}
$F^0=0$ say that the following expressions vanish:
\begin{equation}\label{eq:d3-boundary}
\begin{gathered}
R^z_1=A-m_4D_{14}^3(A-B),\qquad R^z_4=B-m_1D_{14}^3(B-A),\\
R^z_2=A-m_3D_{23}^3(A-B),\qquad R^z_3=B-m_2D_{23}^3(B-A),\\
R^w_1=C-m_2D_{12}^3(C-D),\qquad R^w_2=D-m_1D_{12}^3(D-C),\\
R^w_3=D-m_4D_{34}^3(D-C),\qquad R^w_4=C-m_3D_{34}^3(C-D),\\
\hat z_5+(m_1D_{15}^3+m_2D_{25}^3)A+(m_3D_{35}^3+m_4D_{45}^3)B,\\
\hat w_5+(m_1D_{15}^3+m_4D_{45}^3)C+(m_2D_{25}^3+m_3D_{35}^3)D,\\
D_{12}^2x(D-C)-1,\quad D_{34}^2y(C-D)-1,\quad D_{14}^2u(B-A)-1,\\
D_{23}^2v(B-A)-1,\quad D_{13}^2(B-A)(D-C)-1,\quad D_{24}^2(B-A)(C-D)-1,\\
D_{15}^2AC-1,\quad D_{25}^2AD-1,\quad D_{35}^2BD-1,\quad D_{45}^2BC-1,\quad (B-A)-(D-C);
\end{gathered}
\end{equation}
$R^z_k$, $R^w_k$ come from the equations of body $k$.  We work on the slice $B=A+1$; for a germ
we take $t=1/z_{13}$ (Proposition~\ref{prop:d3-cover}).

\subsection{Coverage}
We write $s$ for the parameter of a germ.

\begin{proposition}[Coverage]\label{prop:d3-cover}
Let $\gamma$ be a Diagram-3 germ, renormalized as in Lemma~\ref{lem:germ}.
\begin{romanlist}
  \item Bodies $1,\dots,4$ are circled in both colours, body $5$ in neither.
  \item The sides are maximal edges: $r_{12},r_{14},r_{23},r_{34}\approx1$,
  $z_{14},z_{23},w_{12},w_{34}\approx\epsilon^{-2}$ and $w_{14},w_{23},z_{12},z_{34}\approx\epsilon^2$.
  \item $z_{13},w_{13},z_{24},w_{24},z_{k5},w_{k5}\approx\epsilon^{-2}$ for $k\le4$, and
  $z_5,w_5\preccurlyeq\epsilon^4$.
  \item Renormalize instead by $\lambda=(w_{13}/z_{13})^{1/2}$, so that $z_{13}=w_{13}$, and put
  $t=1/z_{13}\approx\epsilon^2$.  Then $B-A=1$ and every chart coordinate converges as $s\to0$; the
  limit $X_0$ satisfies \eqref{eq:d3-boundary}, and $A,B,C,D,x,y,u,v$ and all $D_{kl}$ are
  nonzero at $X_0$.
\end{romanlist}
\end{proposition}

\begin{proof}
By Lemma~\ref{lem:germ} the relations of \cite[\S3]{AK} hold between orders along $\gamma$, and
Fact~\ref{fact:rules} applies.  (i) Each side is the only stroke of its colour at its two ends, so
Rule~1a \cite[pp.~545--546]{AK} circles $1,\dots,4$ in both colours; a circle needs a stroke of
its colour, so $5$ has none.
(ii) The isolated component $\{1,4\}$ of the $z$-diagram has $z$-circles, so its stroke $14$ is
maximal by Rule~1e \cite[p.~546]{AK}; likewise $23$, and $12$, $34$ in the $w$-diagram.  The right
equality case of Estimate~2 \cite[p.~548]{AK} gives the orders.
(iii) By (ii), $z_{13}=z_{14}-z_{34}$, $w_{13}=w_{12}+w_{23}$, $z_{24}=z_{23}+z_{34}$ and
$w_{24}=w_{14}-w_{12}$ are $\approx\epsilon^{-2}$; by (i), $z_k\approx Z=\epsilon^{-2}\succ z_5$ for
$k\le4$, and likewise for $w$.  The equation
$z_5(1-\sum_{k\le4}m_k\delta_{k5}^3)=-\sum_{k\le4}m_k\delta_{k5}^3z_k$ of body $5$, with
$\delta_{k5}\approx\epsilon^2$, gives $z_5\preccurlyeq\epsilon^4$; likewise for $w_5$.
(iv) By (iii), $\lambda$ differs from the factor of Lemma~\ref{lem:germ} by a Puiseux series of
order $0$, so (i)--(iii) persist.  Now $B-A=tz_{13}=1$, and $A=z_1/z_{13}$, $B=z_3/z_{13}$,
$C=w_1/w_{13}$, $D=w_2/w_{13}$, $x=z_{12}z_{13}$, $y=z_{34}z_{13}$, $u=w_{14}w_{13}$,
$v=w_{23}w_{13}$ and the $D_{kl}$ ($=\delta_{kl}$ for the sides, $\delta_{kl}z_{13}$ otherwise) are
$\approx1$, while $\hat z_5=z_5z_{13}^2$, $\hat w_5=w_5w_{13}^2\preccurlyeq1$.  So all chart
coordinates converge, and $F=0$ passes to the limit.
\end{proof}

By Proposition~\ref{prop:d3-cover}(ii) and~(iii), every Diagram-3 germ has an unbounded
four-product: $p_{12}=r_{12}^2r_{34}^2r_{35}^2r_{45}^2\approx\epsilon^{-8}$.

\subsection{The first condition}

\begin{proposition}[First condition]\label{prop:d3-first}
At every point of $\{t=0\}$ satisfying \eqref{eq:d3-boundary}, $m_1m_3=m_2m_4$.  Hence every
Diagram-3 germ satisfies (29) of \cite[\S8.3]{AK}.
\end{proposition}

\begin{proof}
Since $m_1R^z_1+m_4R^z_4=m_1A+m_4B$ and $m_2R^z_2+m_3R^z_3=m_2A+m_3B$ identically,
\begin{equation}\label{eq:d3-syzygy}
(B-A)(m_1m_3-m_2m_4)=(m_1+m_4)(m_2R^z_2+m_3R^z_3)-(m_2+m_3)(m_1R^z_1+m_4R^z_4),
\end{equation}
and likewise
$(D-C)(m_1m_3-m_2m_4)=(m_1+m_2)(m_3R^w_3+m_4R^w_4)-(m_3+m_4)(m_1R^w_1+m_2R^w_2)$.  With
$T_{13}=D_{13}^2(B-A)(D-C)-1$ (the distance equation of $13$) and $g=(B-A)-(D-C)$ (the gauge),
$D_{13}^2(B-A)^2-1=T_{13}+D_{13}^2(B-A)g$.  At a boundary point $T_{13}=g=0$, so
$D_{13}^2(B-A)^2=1$, $B-A\ne0$, and \eqref{eq:d3-syzygy} gives $m_1m_3=m_2m_4$.  The second
statement follows from Proposition~\ref{prop:d3-cover}(iv).
\end{proof}

These identities are the chart form of \cite[\S5.2]{AK}, where Rule~1c makes $m_1z_1+m_4z_4$ and
$m_2z_2+m_3z_3$ non-maximal while $z_1\sim z_2$ and $z_3\sim z_4$.

\subsection{Labels and the analytic graph}
From now on $m\in\Rp^5$ satisfies (29), $M=m_1+m_2+m_3+m_4$, and $X$ denotes the $19$ chart
coordinates other than $B=A+1$.  Let $F_{z_k},F_{w_k}$ be the rescaled equations of body $k$, let
$\Lambda$ be the $19\times21$ matrix that forms from $F$ the combinations $m_1F_{z_1}+m_4F_{z_4}$,
$F_{z_4}-F_{z_1}$, $F_{z_3}-F_{z_2}$, $m_1F_{w_1}+m_2F_{w_2}$, $F_{w_2}-F_{w_1}$,
$F_{w_3}-F_{w_4}$ and the other $13$ equations, and put
$\ell_zF=F_{z_1}-F_{z_2}-\frac{m_3}{m_2}F_{z_3}+\frac{m_4}{m_1}F_{z_4}$,
$\ell_wF=F_{w_1}+\frac{m_2}{m_1}F_{w_2}-\frac{m_2}{m_1}F_{w_3}-F_{w_4}$.

\begin{lemma}[computer-assisted]\label{lem:d3-fibre}
\begin{alphalist}
  \item On the slice, \eqref{eq:d3-boundary} has exactly $3^4\cdot2^6=5184$ solutions, the
  \emph{labels}: $D_{12},D_{34},D_{14},D_{23}$ are arbitrary cube roots of $\frac1{m_1+m_2}$,
  $\frac1{m_3+m_4}$, $\frac1{m_1+m_4}$, $\frac1{m_2+m_3}$; $D_{13},D_{24},D_{15},D_{25},D_{35},
  D_{45}$ are arbitrary square roots of $1$, $-1$, $\frac M{m_3}$, $-\frac M{m_4}$, $\frac M{m_1}$,
  $-\frac M{m_2}$; $A=-\frac{m_4}{m_1+m_4}$, $C=-\frac{m_2}{m_1+m_2}$, $D=\frac{m_1}{m_1+m_2}$,
  $x=(m_1+m_2)D_{12}$, $y=-(m_3+m_4)D_{34}$, $u=(m_1+m_4)D_{14}$, $v=(m_2+m_3)D_{23}$, and
  $\hat z_5,\hat w_5$ follow from the equations of body $5$.  The mass $m_5$ does not enter.
  \item At every label
  \[
  \det\partial_X(\Lambda F^0)=-\frac{5184\,m_1^2m_2^3m_4^3\,(D_{12}D_{14}D_{23}D_{34})^4\,
  D_{13}D_{24}D_{15}D_{25}D_{35}D_{45}}{(m_1+m_2)(m_1+m_4)}\neq0 .
  \]
\end{alphalist}
\end{lemma}

The proof is an exact computation over the field $\mathbb Q(m_1,m_2,m_4,m_5)$, with
$m_3=m_2m_4/m_1$.  On the slice, the $19$ entries of $\Lambda F^0$ and the relations listed in (a)
are combinations of each other with coefficients that have no poles at positive masses, and the
determinant in (b) is computed by exact elimination (Appendix~\ref{app:verification}).
Identically in $X$, $\ell_zF^0=-B(m_1m_3-m_2m_4)/(m_1m_2)$ and
$\ell_wF^0=-(D-C)D_{34}^3(m_1m_3-m_2m_4)/m_1$, so $\ell_z,\ell_w$ annihilate
$J_0=\partial_XF^0(X_0)$; as $\Lambda J_0$ is invertible, $J_0$ has rank $19$, $\ell_z,\ell_w$ span
its left kernel, and the $21$ rows of $\Lambda,\ell_z,\ell_w$ are independent.

\begin{lemma}\label{lem:d3-graph}
Let $X_0$ be a label.  There is a unique analytic $t\mapsto X(t)$ near $0$ with $X(0)=X_0$ and
$\Lambda F(t,X(t))=0$; put $G(t)=(\ell_zF(t,X(t)),\ell_wF(t,X(t)))=\sum_{k\ge1}G_kt^k$.  If the
chart coordinates of a Diagram-3 germ converge to $X_0$, then along the germ they are $X(t)$, and
$G\equiv0$.  Every analytic $Y(t)$ with $Y(0)=X_0$ and $F(t,Y(t))\equiv0$ equals $X(t)$.
\end{lemma}

\begin{proof}
The first statement is the analytic implicit function theorem (Lemma~\ref{lem:d3-fibre}(b)).
Along the germ, after a reparametrization $s\mapsto s^N$, the chart coordinates $Y(s)$ and
$t(s)=1/z_{13}(s)$ are convergent power series (Lemma~\ref{lem:germ},
Proposition~\ref{prop:d3-cover}) with $Y(0)=X_0$, $t(0)=0$, $t\not\equiv0$.  Since
$\Lambda F(t(s),Y(s))=0$, uniqueness gives $Y(s)=X(t(s))$.  So $F(t,X(t))$ vanishes on the image
of $t(s)$, a neighbourhood of $0$; hence $F(t,X(t))\equiv0$ and $G\equiv0$.  The last statement is
again uniqueness.
\end{proof}

So a label carries a Diagram-3 germ, ramified or not, only if $G\equiv0$ there.  As $F$ has no
term linear in $t$, $X(t)=X_0+X_2t^2+X_3t^3+\cdots$, and as $\ell J_0=0$ ($\ell=(\ell_z,\ell_w)$),
$G_k$ does not involve $X_k$: at $X_0$, $G_2=\ell F^2$, $G_3=\ell F^3$,
$G_4=\ell(F^4+\partial F^2\cdot X_2+\tfrac12\partial^2F^0[X_2,X_2])$ and
$G_5=\ell(F^5+\partial F^3\cdot X_2+\partial F^2\cdot X_3+\partial^2F^0[X_2,X_3])$.

\subsection{Orders three and four}

\begin{lemma}[computer-assisted]\label{lem:d3-order3}
At every label $G_2=0$ and
$G_3=\bigl(\frac{m_1+m_2}{m_1^2m_2}\Omega,\frac{m_1+m_4}{m_1^2m_4}\Omega\bigr)$, where
$\Omega=\Omega_{\mathrm{re}}+\Omega_{\mathrm{im}}$,
\[
\begin{aligned}
\Omega_{\mathrm{re}}&=m_1m_2m_4D_{13}+m_5m_1^2D_{15}+m_5m_2m_4D_{35},\\
\Omega_{\mathrm{im}}&=m_1m_2m_4D_{24}+m_5m_1m_2D_{25}+m_5m_1m_4D_{45}.
\end{aligned}
\]
\end{lemma}

The computation is exact, over the field $\mathbb Q(m_1,m_2,m_4,m_5)$ with $m_3=m_2m_4/m_1$:
$\ell F^2$ and $\ell F^3$ are evaluated at the labels and reduced modulo the cubic and quadratic
equations satisfied by the $D_{kl}$ (Appendix~\ref{app:verification}).

\begin{corollary}\label{cor:d3-order3}
If $G_3=0$ at some label, then $N_{\mathrm{re}}(m)=N_{\mathrm{im}}(m)=0$, where $Q=m_5^2M$ and
\begin{align*}
N_{\mathrm{re}}&=\bigl((m_1m_2m_4)^3-Q(m_1^6+m_2^3m_4^3)\bigr)^2-4Q^2m_1^6m_2^3m_4^3,\\
N_{\mathrm{im}}&=\bigl((m_1m_2m_4)^3-Qm_1^3(m_2^3+m_4^3)\bigr)^2-4Q^2m_1^6m_2^3m_4^3 .
\end{align*}
As $N_{\mathrm{re}}-N_{\mathrm{im}}+QS=0$ with
$S=(m_1^3-m_2^3)(m_1^3-m_4^3)\bigl(2(m_1m_2m_4)^3-Q(m_1^3+m_2^3)(m_1^3+m_4^3)\bigr)$, also
$S(m)=0$.
\end{corollary}

\begin{proof}
At positive masses $D_{13},D_{15},D_{35}$ are real and $D_{24},D_{25},D_{45}$ purely imaginary
(Lemma~\ref{lem:d3-fibre}), so $\Omega_{\mathrm{re}}=\Omega_{\mathrm{im}}=0$.  If
$\alpha+\beta+\gamma=0$ then $(\alpha^2-\beta^2-\gamma^2)^2-4\beta^2\gamma^2=\prod_\pm(\alpha+\beta
\pm\gamma)\prod_\pm(\alpha-\beta\pm\gamma)=0$.  Inserting the squares $(m_1m_2m_4)^2$,
$m_5^2m_1^4M/m_3$, $m_5^2m_2^2m_4^2M/m_1$ of the terms of $\Omega_{\mathrm{re}}$, using
$m_3=m_2m_4/m_1$ and multiplying by $(m_1m_2m_4)^2$ gives $N_{\mathrm{re}}=0$; the terms of
$\Omega_{\mathrm{im}}$ are $i$ times reals with squares $(m_1m_2m_4)^2$, $m_5^2m_1^2m_2^2M/m_4$,
$m_5^2m_1^2m_4^2M/m_2$, which gives $N_{\mathrm{im}}=0$.  Finally $Q>0$.
\end{proof}

\begin{lemma}[computer-assisted]\label{lem:d3-order4}
At every label $G_4=-\frac{\Pi}{m_1^2}\bigl(m_4(m_1+m_2),m_2(m_1+m_4)\bigr)$ with
$\Pi=(m_1D_{12}-m_4D_{34})(m_1D_{14}-m_2D_{23})$.
\end{lemma}

As for Lemma~\ref{lem:d3-order3}, the computation is exact and reduces modulo the equations
satisfied by the $D_{kl}$; it solves for $X_2$ and evaluates the formula for $G_4$ above
(Appendix~\ref{app:verification}).  At every label
$(m_1D_{12})^3-(m_4D_{34})^3=m_1(m_1+m_4)(m_1-m_4)/(m_1+m_2)$ and
$(m_1D_{14})^3-(m_2D_{23})^3=m_1(m_1+m_2)(m_1-m_2)/(m_1+m_4)$, so $\Pi=0$ forces $m_1=m_4$ or
$m_1=m_2$; with (29), $m=(p,q,q,p,m_5)$ or $(p,p,q,q,m_5)$.

\begin{lemma}[computer-assisted]\label{lem:d3-curves}
Let $m\in\Rp^5$ satisfy (29), $N_{\mathrm{re}}(m)=N_{\mathrm{im}}(m)=0$ and $m_1=m_2$ (resp.\
$m_1=m_4$).  Write $m_1=a^2$, $m_3=b^2$ (resp.\ $m_2=b^2$) with $a,b>0$, and
$c=(2(a^2+b^2))^{1/2}$.  Then
\[
m_5=\frac{a^3b^3}{c\,|a^3+b^3|}\qquad\text{or}\qquad m_5=\frac{a^3b^3}{c\,|a^3-b^3|}\quad(a\ne b).
\]
Conversely $N_{\mathrm{re}}$, $N_{\mathrm{im}}$ and $S$ vanish on these four families, called
$E^\pm_{12}$ (if $m_1=m_2$) and $E^\pm_{14}$ (if $m_1=m_4$).  They meet only at $a=b$, i.e.\ at
$m$ proportional to $(1,1,1,1,\frac14)$, which lies on $E^+_{12}\cap E^+_{14}$.
\end{lemma}

On these mass patterns $N_{\mathrm{re}}=a^{12}\prod_\pm\bigl(a^6b^6-m_5^2c^2(a^3\pm b^3)^2\bigr)$,
a polynomial identity in $a,b,m_5$.  The converse is checked by exact substitution, with $a,b,c$
given by a rational parametrization of the conic $c^2=2(a^2+b^2)$
(Appendix~\ref{app:verification}).

\begin{lemma}[computer-assisted]\label{lem:d3-census}
Let $m$ lie on one of the four families.
\begin{alphalist}
  \item If $a\ne b$, exactly $108$ labels have $G_3=G_4=0$: two of the eight sign choices for
  $(D_{13},D_{15},D_{35})$ give $\Omega_{\mathrm{re}}=0$, two of those for $(D_{24},D_{25},D_{45})$
  give $\Omega_{\mathrm{im}}=0$, and $\Pi=0$ exactly when $D_{14}=D_{23}$ (on $E^\pm_{12}$) or
  $D_{12}=D_{34}$ (on $E^\pm_{14}$), for $27$ of the $81$ choices of cube roots.
  \item If $a=b$, exactly $180$ labels have $G_3=G_4=0$: those with $D_{15}=D_{35}=-2D_{13}$,
  $D_{25}=D_{45}=-2D_{24}$, and $D_{12}=D_{34}$ or $D_{14}=D_{23}$ ($45$ choices of cube roots).
\end{alphalist}
\end{lemma}

An exact computation, with $a,b,c$ as indeterminates, finds the sign choices for which the
relevant linear forms in $a^3,b^3$ vanish identically (Appendix~\ref{app:verification}).  The
count holds at every point: on $E^\pm_{12}$, with $D_{13}=\sigma$, $D_{15}=s_1c/b$,
$D_{35}=s_2c/a$ and $\nu=\pm1$ the sign of $a^3\pm b^3$, $\Omega_{\mathrm{re}}$ is
$a^4b^2/|a^3\pm b^3|$ times $(\sigma\nu+s_1)a^3+(\pm\sigma\nu+s_2)b^3$, whose coefficients lie in
$\{0,2\sigma\nu\}$
and $\{0,\pm2\sigma\nu\}$; so for $a,b>0$ (and $a\ne b$ on $E^-_{12}$) it vanishes only if both
coefficients do.  The other cases are alike.

\subsection{Orders five and six}

\begin{lemma}[computer-assisted]\label{lem:d3-order5}
Let $m$ lie on one of the four families with $a\ne b$.  Then $G_5\neq0$ at each of the $108$
labels of Lemma~\ref{lem:d3-census}(a).  At $a=b$, $G_5=0$ at all $180$ labels of
Lemma~\ref{lem:d3-census}(b).
\end{lemma}

We compute $G_5$ in exact arithmetic over the field $\mathbb Q(i,a,b,c)$, $c^2=2(a^2+b^2)$, modulo
the cubic equations satisfied by $D_{12},D_{34},D_{14}$, and the product over the nine pairs below
as an iterated resultant (Appendix~\ref{app:verification}).  On $E^+_{12}$ the labels have
$D_{14}=D_{23}$, and the $z$-component of $G_5$ is a nonzero factor
times $(a-b)\,L(D_{12}/D_{14},D_{34}/D_{14})$, where
$L(\xi_1,\xi_2)=a^2F_4\xi_1-a^{-2}b^4F_4\xi_2-b^2(a^2-b^2)F_2$,
$F_4=a^4+a^3b+a^2b^2+ab^3+b^4$ and $F_2=a^2+ab+b^2$.  Up to a nonzero monomial factor and sign,
the product of $L$ over the nine pairs with $\xi_1^3=(a^2+b^2)/(2a^2)$, $\xi_2^3=(a^2+b^2)/(2b^2)$
is $(a-b)^3(a+b)^3P_{66}(a,b)$, where $P_{66}$ is homogeneous of degree $66$ with $67$ positive
coefficients; so $G_5\ne0$ if $a\ne b$.
On $E^-_{12}$ the same computation gives $P_{66}(a,-b)$, and an exact Sturm count shows that
$P_{66}(1,-\rho)$ has no root $\rho>0$.  Exchanging the bodies $2$ and $4$ together with $z$ and
$w$ maps the chart at $m$ onto the chart at $(m_1,m_4,m_3,m_2,m_5)$, $E^\pm_{12}$ onto
$E^\pm_{14}$ and the $108$ labels onto the $108$ labels; the relabelling $(13)(24)$ exchanges $a$
and $b$ and covers the minus families with $a<b$.  At $a=b$ both components of $G_5$ vanish at
the labels of $E^+_{12}$, hence, by the exchange, at all $180$ labels.

\begin{lemma}[computer-assisted]\label{lem:d3-order6}
Let $m=(\frac12,\frac12,\frac12,\frac12,\frac18)$.  Of the $180$ labels of
Lemma~\ref{lem:d3-census}(b), the $144$ labels with $D_{12}\ne D_{34}$ or $D_{14}\ne D_{23}$ have
$G_6\neq0$.
\end{lemma}

The proof is a computation of the jets modulo two primes (Appendix~\ref{app:verification}).  All
data lie in $R=\mathbb Z[\frac12,i,\omega]$, $\omega=e^{2\pi i/3}$; let $\mathfrak p$ be a prime of
$R$ above $p\in\{13,37\}$, so $R/\mathfrak p=\mathbb F_p$ as $p\equiv1\pmod{12}$.  If a germ passed
through a label, then $F(t,X(t))\equiv0$ (Lemma~\ref{lem:d3-graph}); each $X_k$ would be the
unique solution of the order-$k$ equations, given by any $19$ rows with invertible minor, and all
$21$ residuals would vanish.  We solve with $19$ rows whose minor is a unit modulo $\mathfrak p$
and find a nonzero residual at order $6$ for the $144$ labels.  Another choice of $\mathfrak p$
acts on the labels by $\mathrm{Gal}(\mathbb Q(i,\omega)/\mathbb Q)$ ($\omega\mapsto\omega^2$,
$i\mapsto-i$), which preserves the set of the $144$ labels; both primes give the same set.

\begin{lemma}[computer-assisted]\label{lem:d3-twisted}
Let $m=(\frac12,\frac12,\frac12,\frac12,\frac18)$, $\alpha,\beta\in\{1,\omega,\omega^2\}$ and
$\varepsilon,\eta\in\{\pm1\}$.  Let $(P,Q,Y)$ be the analytic solution near $\tau=0$ of
\[
(1-P)(1+Q)=\alpha\tau^2,\qquad (1+P)(1-Q)=\beta\tau^2,\qquad 4PQY^2+1=0
\]
with $(P,Q,Y)(0)=(1,1,\eta i/2)$; the Jacobian there is $\pm16i$.  Then
\[
z=\tfrac1\tau(1,P,-1,-P,0),\qquad w=\tfrac1\tau(1,-Q,-1,Q,0),
\]
\[
\delta_{12}=\delta_{34}=\alpha,\ \ \delta_{14}=\delta_{23}=\beta,\ \
\delta_{13}=\tfrac{\varepsilon\tau}2,\ \ \delta_{15}=\delta_{35}=-\varepsilon\tau,\ \
\delta_{24}=\tau Y,\ \ \delta_{25}=\delta_{45}=-2\tau Y
\]
solves \eqref{eq:system} with the gauge $z_{13}=w_{13}$.  In the chart, with $t=-\tau/2$, it is an
analytic solution of $F=0$ through the label with $D_{12}=D_{34}=\alpha$, $D_{14}=D_{23}=\beta$,
$D_{13}=-\varepsilon$ and $D_{24}=-\eta i$.  These $36$ labels are the labels of
Lemma~\ref{lem:d3-census}(b) with $D_{12}=D_{34}$ and $D_{14}=D_{23}$.  Along these germs
\begin{equation}\label{eq:d3-identities}
\delta_{15}+2\delta_{13}=\delta_{35}+2\delta_{13}=\delta_{25}+2\delta_{24}=\delta_{45}+2\delta_{24}=0 .
\end{equation}
\end{lemma}

We substitute into \eqref{eq:system} in exact arithmetic, with $\alpha,\beta,\varepsilon,\tau,P,Q,Y$
as indeterminates: the coordinate equations vanish identically once $\alpha^3=\beta^3=1$, the
distance equations reduce to the three displayed ones, and \eqref{eq:d3-identities} holds without
using them (Appendix~\ref{app:verification}).  The chart coordinates are analytic in $t$ because
$P,Q,Y$ are analytic in $\tau^2$ with $P(0)=Q(0)=1$.

\subsection{The exclusion}

\begin{theorem}\label{thm:d3}
Let $m\in\Rp^5$, let $\Gamma\subset A_m$ be an irreducible curve, and let $\gamma$ be a Diagram-3
germ on $\Gamma$, with the bodies numbered as in Diagram~3.  Then $m_1=m_2=m_3=m_4=4m_5$, and
\eqref{eq:d3-identities} holds identically on $\Gamma$.  In particular $\Gamma$ contains no
positive point, and Theorem~\ref{thm:exclusion}(c) holds.
\end{theorem}

\begin{proof}
Exchanging $z$ and $w$ and renumbering if necessary, we may assume that $\gamma$ has the diagram
described at the beginning of this section; the conclusions are invariant under these
operations.  By Proposition~\ref{prop:d3-cover} the chart coordinates of $\gamma$ converge to a
solution $X_0$ of \eqref{eq:d3-boundary} with $B-A=1$.  By Proposition~\ref{prop:d3-first},
(29) holds, so $X_0$ is a label (Lemma~\ref{lem:d3-fibre}), and $G\equiv0$ at $X_0$
(Lemma~\ref{lem:d3-graph}).  By Lemma~\ref{lem:d3-order3} and Corollary~\ref{cor:d3-order3},
$N_{\mathrm{re}}(m)=N_{\mathrm{im}}(m)=0$.  By Lemma~\ref{lem:d3-order4}, $m_1=m_4$ or $m_1=m_2$,
so $m$ lies on one of the families of Lemma~\ref{lem:d3-curves}.  By
Lemmas~\ref{lem:d3-census}(a) and~\ref{lem:d3-order5}, $a=b$, that is $m_1=m_2=m_3=m_4=4m_5$.

The system \eqref{eq:system} is covariant under
$(m,z,w,\delta)\mapsto(\kappa^3m,\kappa z,\kappa w,\kappa^{-1}\delta)$, $\kappa>0$, which maps
Diagram-3 germs to Diagram-3 germs and preserves \eqref{eq:d3-identities}.  So we may take
$m=(\frac12,\frac12,\frac12,\frac12,\frac18)$.
By Lemmas~\ref{lem:d3-census}(b) and~\ref{lem:d3-order6}, $D_{12}=D_{34}$ and $D_{14}=D_{23}$ at
$X_0$.  By Lemmas~\ref{lem:d3-twisted} and~\ref{lem:d3-graph}, the chart coordinates of $\gamma$
are those of the twisted Roberts germ through $X_0$.  The renormalizations do not change the
$\delta_{kl}$, so \eqref{eq:d3-identities} holds along $\gamma$.  Each expression in
\eqref{eq:d3-identities} is a regular function on the irreducible curve $\Gamma$ that vanishes at
the infinitely many points of $\gamma$, so it vanishes on $\Gamma$.  At a positive point
$\delta_{15}+2\delta_{13}>0$.
\end{proof}

So every Diagram-3 germ has the masses $(4,4,4,4,1)$ up to scale and is, up to the symmetries
used above, one of the $36$ twisted Roberts germs, which do exist.  The next remark shows that
the theorem cannot be weakened to curves through a real point.

\begin{remark}[The modified Roberts continuum]\label{rem:d3-roberts}
Let $m=(\frac12,\frac12,\frac12,\frac12,\frac18)$, proportional to $(4,4,4,4,1)$
\cite[Rem.~8, p.~583]{AK}.  The curve
\[
z=(a,b,-a,-b,0),\quad w=(a,-b,-a,b,0),\quad b^2=a^2-1,
\]
\[
\delta_{12}=\delta_{23}=\delta_{34}=\delta_{14}=1,\ \ \delta_{13}=\tfrac1{2a},\ \
\delta_{24}=\tfrac{i}{2b},\ \ \delta_{15}=\delta_{35}=-\tfrac1a,\ \ \delta_{25}=\delta_{45}=-\tfrac ib
\]
is the modified Roberts continuum \cite[p.~543]{AK}.  It solves \eqref{eq:system} with the gauge
$z_{13}=w_{13}$, $U=I=1$ on it, and \eqref{eq:d3-identities} holds, so it has no positive point.
Since $w_{12}/z_{12}=(a+b)^2$ on it, the rotation by $\mu=a+b$ gives an irreducible curve
$\Gamma_R\subset A_m$.  For $0<a<1$ and $b\in i\R$, $|\mu|=1$ and the points are real, with
$\delta_{13}>0>\delta_{15}=\delta_{35}$.
As $a\to\infty$ with $b\sim a$, the curve is the twisted Roberts germ of
Lemma~\ref{lem:d3-twisted} with $(\alpha,\beta,\varepsilon,\eta)=(1,1,1,1)$, $\tau=1/a$ and
$P=Q=b/a$.  Its strokes are $z_{14}$, $z_{23}$, $w_{12}$, $w_{34}$, so it is a Diagram-3 germ
(Roberts' continuum at infinity \cite[\S7.2]{AK}); the sides keep length $1$,
$r_{13},r_{24},r_{k5}\to\infty$, $p_{12}=a^2(1-a^2)$ is unbounded, and (29) holds.
So $\Gamma_R$ passes through real points, carries the positive constant $U=I=1$, and has a
Diagram-3 germ with an unbounded four-product.  An exclusion at the level of curves through a
real point is therefore false for Diagram~3.  Theorem~\ref{thm:d3} uses the positivity of points
only through the identity $\delta_{15}+2\delta_{13}=0$ on all of $\Gamma$, as
Remark~\ref{rem:levels} requires.
\end{remark}



\section{Diagrams 4 and 5: the shared triangle}\label{sec:d45}

Diagrams 4 and 5 contain the isolated uncircled $zw$-triangle $345$.

\begin{proposition}[Relation (23)]\label{prop:d45-first}
Let $m\in\Rp^5$ and let $\gamma$ be a Diagram 4 or Diagram 5 germ, renormalized as in
Lemma~\ref{lem:germ} so that $\mathcal Z$ and $\mathcal W$ have the same least order $-2\nu<0$.
Let $z_c$ and $w_c$ be the centres of mass of $T=\{3,4,5\}$.  Then the limits
\[
  \zeta_k=\lim_{t\to0}t^{-\nu}(z_k-z_c),\qquad \eta_k=\lim_{t\to0}t^{-\nu}(w_k-w_c),\qquad
  D_{kl}=\lim_{t\to0}t^{\nu}\delta_{kl}\qquad(k,l\in T)
\]
exist and solve the zero-multiplier triangle system of Section~\ref{sec:d1}, and
$\sum_km_k\zeta_k=\sum_km_k\eta_k=0$.  Hence $(\zeta,\eta,D)$ is as in
Lemma~\ref{lem:d1-shape}, and $Q(m_3,m_4,m_5)=0$, that is, the masses satisfy (23).
\end{proposition}

\begin{proof}
The sides of $T$ are $zw$-edges, so by Estimate~1 and the left equality case of Estimate~2
(Fact~\ref{fact:rules}, read through Lemma~\ref{lem:germ}) $\ord z_{kl}=\ord w_{kl}=\nu$ and
$\ord\delta_{kl}=-\nu$ for $k,l\in T$.
As $z_k-z_c$ and $w_k-w_c$ are combinations of these differences, the limits exist, and the sums
$\sum_km_k\zeta_k$, $\sum_km_k\eta_k$ vanish by the choice of the centres.  Let $k\in T$.  In the
$z$-row $z_k=\sum_{l\ne k}m_l\delta_{kl}^3(z_k-z_l)$ the terms with $l\in T$ have order
$-2\nu$.  The term $z_k$ has larger order, because $k$ is not circled, and so have the terms
with $l\notin T$, because $T$ is an isolated component of the diagram and no $z$-stroke leaves
it.  Multiplying by $t^{2\nu}$ and letting $t\to0$ gives $\phi_k=0$; likewise $\psi_k=0$, and
the ties of the pairs in $T$ give $\chi_{kl}=0$.  Lemma~\ref{lem:d1-shape} now applies, and for
positive masses $Q=0$ is equivalent to (23).
\end{proof}

\subsection{Diagram 4}\label{sec:d4}

Diagram~4 consists of the isolated $zw$-triangle $345$, whose vertices are not circled, and the
isolated $z$-stroke $12$ with $z$-circles at $1$ and $2$; there are no other strokes or circles
\cite[\S7.4, p.~567, and Fig.~11]{AK}.
For this diagram \cite[\S8.4]{AK} record only relation (23).  We use the notation of
Section~\ref{sec:d1}: $P=\{1,2\}$, $T=\{3,4,5\}$, $i\in P$, $k,l\in T$, $M_P$, $M_T$, $Z_P$, $Z_T$,
$W_P$, $W_T$, the equations $E^z_j,E^w_j$, the triangle functions $\phi_k,\psi_k,\chi_{kl}$, the
torque $T_P$ of Lemma~\ref{lem:d1-syzygies} and the orbit $O$ of Lemma~\ref{lem:d1-mb}.
Renumbering inside $P$ or inside $T$ does not change the diagram.

\begin{theorem}\label{thm:d4}
Let $m\in\Rp^5$.  No curve in $A_m$ carries a Diagram~4 germ.
\end{theorem}

This is part~(a) of Theorem~\ref{thm:exclusion} for $r=4$, without any assumption on the
four-products.  A torque argument excludes a triangle body at the $w$-centre of the triangle
(Lemma~\ref{lem:d4-central}); then every germ converges in a ramified chart to a point of a
smooth boundary fibre $\mathcal B_4$ (Proposition~\ref{prop:d4-cover} and
Lemma~\ref{lem:d4-fibre}).  A Lyapunov--Schmidt lemma (Lemma~\ref{lem:d4-ls}) turns the
existence of a germ through a boundary point into the vanishing of successive obstructions.
The first says that the triangle exerts equal leading $w$-forces on the two outer bodies, the
second vanishes identically, and the third says that it exerts equal leading $z$-forces
(Lemma~\ref{lem:d4-obstr}).  An exact census shows that the two balances never hold together
(Lemma~\ref{lem:d4-census}).

\subsubsection{The chart}\label{d4-chart}
The equations of \eqref{eq:system} other than the gauge are invariant under
$(z,w,\delta)\mapsto(cz,c^{-1}w,\delta)$, $c\in\C^*$; in this subsection we use the gauge
$z_{34}=w_{34}$ instead of $z_{12}=w_{12}$.  Let $t$ be a parameter, put
\[
  (\zeta_3,\zeta_4,\zeta_5)=(m_5u,\ m_5v,\ -m_3u-m_4v),\qquad
  (\eta_3,\eta_4,\eta_5)=(m_5p,\ m_5q,\ -m_3p-m_4q),
\]
so that $\sum_km_k\zeta_k=\sum_km_k\eta_k=0$, and
\begin{equation}\label{eq:d4-chart}
\begin{gathered}
  z_i=t^{-4}A_i,\qquad w_i=t^{4}B_i\qquad(i\in P),\\
  z_k=t^{-1}\hat Z+t^{2}\zeta_k,\qquad w_k=t^{5}\hat W+t^{2}\eta_k\qquad(k\in T),\\
  \delta_{12}=D_{12},\qquad\delta_{ik}=t\,D_{ik},\qquad\delta_{kl}=t^{-2}D_{kl}.
\end{gathered}
\end{equation}
The twenty chart coordinates $x$ are $A_1,A_2,B_1,B_2,\hat Z,\hat W,u,v,p,q$ and the ten $D_{jl}$;
for $t\ne0$ the chart is invertible.  Put
\[
\begin{gathered}
  \Sigma_i=\sum_km_kD_{ik}^3,\qquad f_i=\sum_km_kD_{ik}^3\eta_k,\\
  R^z_i=A_i-m_{i'}D_{12}^3(A_i-A_{i'}),\qquad R^w_i=B_i-m_{i'}D_{12}^3(B_i-B_{i'}),
\end{gathered}
\]
where $\{i,i'\}=P$.  Substituting \eqref{eq:d4-chart} gives
\begin{equation}\label{eq:d4-rows}
\begin{aligned}
  t^{4}E^z_i&=R^z_i-\textstyle\sum_km_kD_{ik}^3\,(t^3A_i-t^6\hat Z-t^9\zeta_k),\\
  t^{-4}E^w_i&=R^w_i+t\,f_i-\textstyle\sum_km_kD_{ik}^3\,(t^3B_i-t^4\hat W),\\
  t^{4}E^z_k&=-\phi_k+t^3\hat Z+t^6\zeta_k-\textstyle\sum_im_iD_{ik}^3\,(t^6\hat Z+t^9\zeta_k-t^3A_i),\\
  t^{4}E^w_k&=-\psi_k+t^6\eta_k+t^9\hat W-\textstyle\sum_im_iD_{ik}^3\,(t^9\eta_k+t^{12}\hat W
  -t^{11}B_i),
\end{aligned}
\end{equation}
while the ties become $D_{12}^2(A_1-A_2)(B_1-B_2)-1$, the \emph{crossed ties}
$D_{ik}^2(t^3\hat Z+t^6\zeta_k-A_i)(\eta_k-t^2B_i+t^3\hat W)-1$ and $\chi_{kl}$, and the gauge
becomes $t^{-2}(z_{34}-w_{34})=(\zeta_4-\zeta_3)-(\eta_4-\eta_3)$.  These $21$ polynomials form
$F(t,x)=\sum_{j\ge0}t^jF_j(x)$.  The \emph{boundary fibre} $\mathcal B_4=\mathcal B_4(m)$ is the
zero set of $F_0$.  By \eqref{eq:d4-rows}, $F_1$ vanishes except in the two rows
$t^{-4}E^w_i$, where it equals $f_i$, and $F_2$ vanishes except in the crossed ties, where it
equals $D_{ik}^2A_iB_i$.

\subsubsection{Coverage}\label{d4-cover}
Let $\gamma$ be a Diagram~4 germ at $m\in\Rp^5$, renormalized as in Lemma~\ref{lem:germ} so that
$Z\approx W\approx\epsilon^{-2}$.  After a base change of the parameter $\tau$ of the branch,
all coordinates are Laurent series in $\tau$ and $\ord_\tau Z=\ord_\tau W=-4q$ for an integer
$q\ge1$; we put $t=\tau^q$ and measure orders in units of $\ord t$, so that the maximal
components of $\mathcal Z$ and $\mathcal W$ have order $-4$ and $\ord\epsilon=2$.

\begin{lemma}\label{lem:d4-val}
Along $\gamma$, for $i\in P$ and $k\ne l$ in $T$:
\begin{romanlist}
  \item $\ord z_k>-4$ and $\ord z_{ik}=-4$;
  \item $\ord z_1=\ord z_2=\ord z_{12}=-4$, $\ord\delta_{12}=0$, $\ord w_{12}=4$,
  $\ord w_1,\ord w_2\ge4$, and hence $\ord W_T=\ord W_P\ge4$;
  \item $\ord z_{kl}=\ord w_{kl}=2$ and $\ord\delta_{kl}=-2$.
\end{romanlist}
\end{lemma}

\begin{proof}
(i) The bodies of $T$ are not $z$-circled, while $1$ and $2$ are.
(ii) Bodies $1,2$ are $z$-circled, so $\ord z_1=\ord z_2=-4$.  The isolated component $\{1,2\}$
of the $z$-diagram has circled vertices, so its stroke $12$ is maximal by Rule~1e
\cite[p.~546]{AK}, and the right equality case of Estimate~2 \cite[p.~548]{AK} gives
$r_{12}\approx1$, $z_{12}\approx\epsilon^{-2}$ and $w_{12}\approx\epsilon^2$.  The triangle edges are
not maximal, since by (iii), whose proof does not use (ii), $z_{kl}\approx\epsilon$ on them; so
$12$ is the only maximal $z$-edge and, by Estimate~1
\cite[p.~548]{AK} with $z$ and $w$ exchanged, $w_{12}\prec w_{ik}$ for $i\in P$, $k\in T$.  As in
the proof of \cite[Prop.~4, p.~564]{AK}, \cite[Prop.~1, p.~559]{AK} with $z$ and $w$ exchanged,
applied to bodies $1$ and $2$, puts the origin in the $w$-cluster of size $w_{12}$ formed by
$w_1$ and $w_2$, so $w_1,w_2\preccurlyeq w_{12}$.  Finally $M_PW_P+M_TW_T=0$.
(iii) The triangle edges are $zw$-edges, so $r_{kl}\approx z_{kl}\approx w_{kl}\approx\epsilon$ by
Estimate~1 and the left equality case of Estimate~2.
\end{proof}

\begin{lemma}[computer-assisted]\label{lem:d4-torque}
Let $0<e<2$.  In \eqref{eq:d4-chart} replace $\hat Z$ by $t^{-3e/2}\hat Z$, $\hat W$ by
$t^{-e/2}\hat W$, $(\eta_3,\eta_4,\eta_5)$ by $(t^e\omega,\ m_5q,\ -m_3m_5^{-1}t^e\omega-m_4q)$ and
$D_{i3}$ by $t^{-e/2}d_i$ $(i\in P)$.  Then
\[
  T_P=t^{1-e/2}\Big(m_3\sum_{i\in P}m_id_i\,\big(d_i^2A_i\omega\big)+\sum_{g\in G}c_g\,t^g\Big),
\]
where the $c_g$ are polynomials in the chart coordinates and
$G=\{\tfrac12e,\ \tfrac32e,\ 3,\ 3-\tfrac32e,\ 5-\tfrac52e,\ 5-e,\ 8-e,\ 8+\tfrac12e\}$; all
elements of $G$ are positive.
\end{lemma}
We obtain this expansion by an exact computation over $\mathbb Q$, with the masses and the chart
coordinates as indeterminates and with exponents of $t$ that are affine functions of $e$
(Appendix~\ref{app:verification}).

\begin{lemma}\label{lem:d4-central}
For $k\in T$ let $o_k$ be the coefficient of $t^2$ in $w_k-W_T$ along $\gamma$.  Then
$o_3,o_4,o_5\ne0$.
\end{lemma}

\begin{proof}
By Lemma~\ref{lem:d4-val}(iii), $o_l-o_k$ is the leading coefficient of $w_{kl}$, so at most one
$o_k$ vanishes.  Renumbering $T$, suppose that $o_3=0$ and let $\ord(w_3-W_T)=2+e$ with $e>0$.
By Lemma~\ref{lem:d4-val}(ii), $\ord(W_T-w_i)\ge4$.

If $e\ge2$, then $\ord w_{i3}\ge4$, so $\ord\delta_{i3}=-\frac12(\ord z_{i3}+\ord w_{i3})\le0$ and
$\ord Z_{i3}=\ord z_{i3}+3\ord\delta_{i3}\le-4$.  Then $Z_{i3}$ is a maximal component of
$\mathcal Z$, that is, a $z$-stroke between $i$ and $3$, which Diagram~4 does not have.

Let $0<e<2$.  Then $\ord w_{i3}=2+e$ and $\ord\delta_{i3}=1-\frac e2$, while $\ord w_{ik}=2$ and
$\ord\delta_{ik}=1$ for $k=4,5$.  The first two identities of Lemma~\ref{lem:d1-syzygies} express
$M_PW_P$ and $M_TZ_T$ as sums of crossed terms $m_im_k\delta_{ik}^3(w_i-w_k)$ and
$m_im_k\delta_{ik}^3(z_k-z_i)$, of orders at least $5-\frac e2$ and $-1-\frac32e$; with
$M_PW_P+M_TW_T=0$ this gives $\ord W_T\ge5-\frac e2$ and $\ord Z_T\ge-1-\frac32e$.  So $\gamma$ is
described by the chart of Lemma~\ref{lem:d4-torque}, with all its coordinates of nonnegative
order and $A_i,\omega,d_i$ of order $0$.  At $t=0$ the ties $\delta_{i3}^2z_{i3}w_{i3}=1$ give
$d_i^2A_i\omega=-1$, and $m_1A_1+m_2A_2=t^4(m_1z_1+m_2z_2)=-t^4M_TZ_T$ vanishes.  By
Lemma~\ref{lem:d4-torque}, $t^{e/2-1}T_P$ tends to $-m_3(m_1d_1+m_2d_2)$, which is nonzero because
\[
  (m_1d_1)^2-(m_2d_2)^2=-\frac{m_1^2}{A_1\omega}+\frac{m_2^2}{A_2\omega}
  =-\frac{m_1^3+m_2^3}{m_1A_1\omega}\ne0 .
\]
But $T_P=0$ on $A_m$ by Lemma~\ref{lem:d1-syzygies}.
\end{proof}

\begin{proposition}[Coverage]\label{prop:d4-cover}
Let $m\in\Rp^5$ and let $\gamma$ be a Diagram~4 germ.  Then $m$ satisfies (23), and
$\ord Z_T\ge-1$, $\ord W_T\ge5$ along $\gamma$.  Renormalize by $(z,w)\mapsto(cz,c^{-1}w)$ with
$c=(w_{34}/z_{34})^{1/2}$, so that $z_{34}=w_{34}$.  After a base change of $\tau$, every chart
coordinate of \eqref{eq:d4-chart} is a power series in $\tau$; the value $y$ of $x$ at $\tau=0$
lies in $\mathcal B_4(m)$, and $\eta_k(y)\ne0$ for $k\in T$.
\end{proposition}

\begin{proof}
Relation (23) is Proposition~\ref{prop:d45-first}.  By Lemmas~\ref{lem:d4-val}
and~\ref{lem:d4-central}, $\ord w_{ik}=2$ for all $i\in P$, $k\in T$, hence $\ord\delta_{ik}=1$.
The first two identities of Lemma~\ref{lem:d1-syzygies} then give $\ord Z_T\ge3-4=-1$ and
$\ord W_P\ge3+2=5$, and $\ord W_T=\ord W_P$.  The factor $c$ has order $0$ by
Lemma~\ref{lem:d4-val}(iii), so all orders persist after the renormalization.  Inverting
\eqref{eq:d4-chart}, $A_i=t^4z_i$, $B_i=t^{-4}w_i$, $\hat Z=tZ_T$, $\hat W=t^{-5}W_T$,
$\zeta_k=t^{-2}(z_k-Z_T)$, $\eta_k=t^{-2}(w_k-W_T)$, $D_{12}=\delta_{12}$, $D_{ik}=t^{-1}\delta_{ik}$
and $D_{kl}=t^2\delta_{kl}$; by Lemma~\ref{lem:d4-val} all have nonnegative order, and
$\eta_k\to o_k\ne0$.  The rows of $F$ are monomial multiples of the equations, so they vanish
along $\gamma$, and $F_0(y)=0$.
\end{proof}

\subsubsection{The boundary fibre}\label{d4-fibre}
On $\C^{20}$ consider the scalings
\[
\begin{gathered}
  \alpha\cdot x:\ (A,B,D_{ik})\mapsto(\alpha^2A,\ \alpha^{-2}B,\ \alpha^{-1}D_{ik}),\\
  \beta\cdot x:\ (\zeta,\eta,D_{kl},D_{ik})\mapsto(\beta^2\zeta,\ \beta^2\eta,\ \beta^{-2}D_{kl},\
  \beta^{-1}D_{ik}),
\end{gathered}
\]
$\alpha,\beta\in\C^*$, the other coordinates being fixed, and the translations of $(\hat Z,\hat W)$.
They preserve $F_0$ up to nonzero factors on each row, hence preserve $\mathcal B_4(m)$; they
are not symmetries of $F$.

\begin{lemma}\label{lem:d4-fibre}
Let $m\in\Rp^5$ satisfy (23) and let $y\in\mathcal B_4(m)$.
\begin{alphalist}
  \item $D_{12}^3=M_P^{-1}$, $A_2=\rho A_1$ and $B_2=\rho B_1$ with $\rho=-m_1/m_2$,
  $D_{12}^2(A_1-A_2)(B_1-B_2)=1$ and $D_{ik}^2A_i\eta_k=-1$; in particular $A_i,B_i,\eta_k\ne0$.
  The point $(\zeta_3,\zeta_4,\eta_3,\eta_4,D_{34},D_{35},D_{45})$ solves the zero-multiplier
  triangle system, so, after renumbering $T$, it lies on the orbit $O$ of Lemma~\ref{lem:d1-mb}.
  \item $dF_0(y)$ has rank $16$.  Its kernel is spanned by $\partial_{\hat Z}$, $\partial_{\hat W}$
  and the generators of the two scalings.  Its cokernel is spanned by the linear form
  \[
    v_P(r)=(B_1-B_2)\,(r^z_1-r^z_2)-(A_1-A_2)\,(r^w_1-r^w_2),
  \]
  where $r^z_i,r^w_i$ are the components of $r\in\C^{21}$ in the rows $t^4E^z_i$, $t^{-4}E^w_i$, and
  by four linear forms that involve only the rows $t^4E^z_k$, $t^4E^w_k$, $\chi_{kl}$ and the gauge.
  \item Near $y$, $\mathcal B_4(m)$ coincides with the orbit of $y$ under the scalings and
  the translations; it is a complex submanifold of dimension $4$ with tangent space
  $\ker dF_0(y)$ at $y$.
\end{alphalist}
\end{lemma}

\begin{proof}
(a) Identically $R^z_1-R^z_2=(A_1-A_2)(1-M_PD_{12}^3)$ and $m_1R^z_1+m_2R^z_2=m_1A_1+m_2A_2$, and
likewise for $R^w$.  The tie of $12$ gives $A_1\ne A_2$ and $B_1\ne B_2$; so $F_0(y)=0$ gives
$D_{12}^3=M_P^{-1}$, $m_1A_1+m_2A_2=0$ and $m_1B_1+m_2B_2=0$.  The crossed ties are
$D_{ik}^2A_i\eta_k=-1$.  The triangle rows of $F_0$ are $\phi_k,\psi_k,\chi_{kl}$ and the gauge;
the last statement is the proof of Proposition~\ref{prop:d1-first}.
(b) Order the coordinates as $(D_{ik};\,A,B,D_{12};\,u,v,p,q,D_{kl};\,\hat Z,\hat W)$ and the rows
as (crossed ties; pair rows $R^z_i,R^w_i$ and the tie of $12$; triangle rows).  The pair rows
involve only the pair coordinates and the triangle rows only the triangle coordinates, and
$D_{ik}$ occurs only in its own crossed tie, with $\partial/\partial D_{ik}=-2D_{ik}A_i\eta_k\ne0$.
So $dF_0(y)$ is block lower triangular, and its rank is $6+\operatorname{rank}J_P+
\operatorname{rank}J_T$, with $J_P$, $J_T$ the pair and triangle blocks.  At $y$, the
differentials of $R^z_1-R^z_2$ and $R^w_1-R^w_2$ are $-3M_PD_{12}^2$ times
$(A_1-A_2)\,dD_{12}$ and $(B_1-B_2)\,dD_{12}$, those of $m_1R^z_1+m_2R^z_2$ and $m_1R^w_1+m_2R^w_2$
are $m_1dA_1+m_2dA_2$ and $m_1dB_1+m_2dB_2$, and the tie of $12$ adds $D_{12}^2$ times
$(B_1-B_2)(dA_1-dA_2)+(A_1-A_2)(dB_1-dB_2)$ modulo $dD_{12}$; so $\operatorname{rank}J_P=4$, the
kernel of $J_P$ is spanned by $(A,-B,0)$, and $v_P$, the derivative at $y$ of the identity
$(B_1-B_2)(R^z_1-R^z_2)-(A_1-A_2)(R^w_1-R^w_2)=0$, spans its cokernel.  By
Lemma~\ref{lem:d1-mb} the rows $\phi_3,\phi_4,\psi_3,\psi_4,\chi_{kl}$ have rank $5$ at $y$, with
kernel spanned by the generators $(\zeta,0,-\frac12D)$ and $(0,\eta,-\frac12D)$ of $\mathbb T$;
$\phi_5$ and $\psi_5$ add nothing because $\sum_km_k\phi_k=\sum_km_k\psi_k=0$; and the gauge row
takes the value $\zeta_4-\zeta_3$, which is nonzero by $\chi_{34}=0$, on the first generator and
$0$ on $(\zeta,\eta,-D)$.  So
$\operatorname{rank}J_T=6$, with kernel spanned by $(\zeta,\eta,-D)$.  Hence the rank is $16$, and
the four listed vectors, completed in the $D_{ik}$ by the crossed ties, span the kernel.  A
linear form vanishing on the image vanishes on the crossed ties (columns $D_{ik}$), so it is the
sum of a form in the cokernel of $J_P$ and one in the cokernel of $J_T$, which has dimension
$10-6=4$.
(c) The orbit map has differential onto $\ker dF_0(y)$, so the orbit of $y$ is a submanifold of
dimension $4$ near $y$, contained in $\mathcal B_4(m)$.  By (b) and the rank theorem,
$\mathcal B_4(m)$ is contained near $y$ in a submanifold of dimension $20-16=4$.
\end{proof}

\subsubsection{A Lyapunov--Schmidt lemma}\label{d4-ls}
For a polynomial curve $x(t)$ we write $[F]_j$ for the coefficient of $t^j$ in $F(t,x(t))$.

\begin{lemma}[Lyapunov--Schmidt]\label{lem:d4-ls}
Let $F(t,x)=\sum_{j\ge0}t^jF_j(x)$ be a polynomial map $\C\times\C^n\to\C^p$, let
$Z=F_0^{-1}(0)$, and let $y\in Z$ be a point near which $Z$ is a complex submanifold with
$T_{y'}Z=\ker dF_0(y')$ for all $y'\in Z$.  For $y'\in Z$ let $Q_{y'}$ be the cokernel of
$dF_0(y')$.  Let $N\ge1$ and suppose that every $y'\in Z$ near $y$ is the value at $t=0$ of a
polynomial curve $x(t)$ with $[F]_j=0$ for $j<N$.
\begin{alphalist}
  \item The class $\Omega_N(y')\in Q_{y'}$ of $[F]_N$ does not depend on the choice of such a curve.
  \item If $\Omega_N(y)\ne0$, no holomorphic germ $c\colon(\C,0)\to(F^{-1}(0),(0,y))$ has
  $t\circ c\not\equiv0$.
\end{alphalist}
For $N=1$ the hypothesis is empty and $\Omega_1(y')$ is the class of $F_1(y')$.
\end{lemma}

\begin{proof}
Let $K=\ker dF_0(y)$ and $R=dF_0(y)(\C^n)$, choose complements $\C^n=K\oplus K'$ and
$\C^p=R\oplus R'$, and let $\pi$ be the projection onto $R$ along $R'$.  Since $\pi\,dF_0(y)$ maps
$K'$ isomorphically onto $R$, the implicit function theorem gives a holomorphic $h(t,k)\in K'$,
defined near $(0,0)\in\C\times K$ with $h(0,0)=0$, such that near $(0,y)$ the equation
$\pi F(t,y+k+h')=0$ with small $h'\in K'$ holds exactly for $h'=h(t,k)$; uniqueness also holds for
formal power series in $t$.  Put $G(t,k)=F(t,y+k+h(t,k))\in R'$.  Near $(0,y)$,
$F^{-1}(0)=\{(t,y+k+h(t,k)):G(t,k)=0\}$.  The map $k\mapsto y+k+h(0,k)$ parametrizes the
submanifold $\{\pi F_0=0\}$, of dimension $\dim K=\dim Z$, which contains $Z$; so it
parametrizes $Z$ near $y$, and $G(0,k)=0$.  Write $G=\sum_{j\ge1}t^jG_j(k)$.  For $y'\in Z$ near
$y$ the space $R'$ is also a complement of $dF_0(y')(\C^n)$, so $R'\to Q_{y'}$ is an isomorphism.

Let $j\le N$, let $y'=y+k'+h(0,k')$, suppose that $G_1,\dots,G_{j-1}$ vanish near $k'$, and let
$x(t)=y+\kappa(t)+\eta(t)$, with $\kappa(t)\in K$, $\eta(t)\in K'$, be a polynomial curve with
$x(0)=y'$ and $[F]_i=0$ for $i<j$.  Since $\pi F(t,x(t))=O(t^j)$, formal uniqueness gives
$\eta(t)-h(t,\kappa(t))=e\,t^j+O(t^{j+1})$ for some $e$.  By Taylor's formula
$F(t,x(t))=F(t,y+\kappa+h(t,\kappa))+dF_0(y')(e)\,t^j+O(t^{j+1})$, and
$F(t,y+\kappa+h(t,\kappa))=G(t,\kappa(t))=t^jG_j(k')+O(t^{j+1})$.  So $[F]_j\equiv G_j(k')$
modulo $dF_0(y')(\C^n)$.  For $j<N$ the curves of the hypothesis give $G_j(k')\in R'\cap
dF_0(y')(\C^n)=0$ for all $k'$ near $0$; by induction $G_1,\dots,G_{N-1}$ vanish near $0$, and the
case $j=N$ gives (a), $\Omega_N(y')$ being the image of $G_N(k')$.  If $\Omega_N(y)\ne0$, then
$G=t^NH$ with $H$ holomorphic and $H(0,0)=G_N(0)\ne0$, so near $(0,0)$ the equation $G=0$
forces $t=0$.  A germ $c$ as in (b) has the form $(t,y+k+h(t,k))$ with $G(t,k)=0$, so $t\circ
c\equiv0$.
\end{proof}

\subsubsection{The first three obstructions}\label{d4-obstructions}
For $y\in\mathcal B_4(m)$ put $S_i=\sum_km_kD_{ik}$ and
\[
  \Theta=\sum_km_k\,\frac{D_{1k}-D_{2k}}{\eta_k}=A_2\Sigma_2-A_1\Sigma_1 ;
\]
the second expression and $f_i=-S_i/A_i$ follow from $1/\eta_k=-A_iD_{ik}^2$.  Under the
scalings, $S_i/A_i$ is multiplied by $\alpha^{-3}$ and by $\beta^{-1}$, and $\Theta$ by
$\alpha^{-1}$ and by $\beta^{-3}$; neither involves $\hat Z$, $\hat W$, $D_{12}$ or $m_1,m_2$.
Since the crossed forces on body $i$ are
$\sum_km_k\delta_{ik}^3(z_i-z_k)=t^{-1}A_i\Sigma_i+O(t^2)$ and
$\sum_km_k\delta_{ik}^3(w_i-w_k)=-t^5f_i+O(t^7)$, the conditions $f_1=f_2$ and $\Theta=0$ say
that the triangle pulls the two outer bodies equally, at leading order, in the $w$- and in the
$z$-coordinate.

\begin{lemma}\label{lem:d4-obstr}
Let $m\in\Rp^5$ satisfy (23), let $y\in\mathcal B_4(m)$, and let $\Omega_N$ be the classes of
Lemma~\ref{lem:d4-ls} for $F$ and $Z=\mathcal B_4(m)$.
\begin{alphalist}
  \item $\Omega_1(y)=0$ if and only if $f_1=f_2$, that is, $S_1/A_1=S_2/A_2$.
  \item If $f_1=f_2$ at $y$, then the hypothesis of Lemma~\ref{lem:d4-ls} holds at $y$ for
  $N=3$, and
  \[
    v_P\big(\Omega_3(y)\big)=\tfrac32\,(A_1-A_2)\,\frac{B_1}{A_1}\,\Theta .
  \]
\end{alphalist}
\end{lemma}

\begin{proof}
(a) By Lemma~\ref{lem:d4-fibre}(b), $\Omega_1(y)$, the class of $F_1(y)$, vanishes if and only if
$F_1(y)$ is killed by $v_P$ and by the four triangle forms.  $F_1$ vanishes on the triangle
rows, and $v_P(F_1(y))=-(A_1-A_2)(f_1-f_2)$.

(b) By homogeneity $f_1=f_2$ on the orbit of $y$, which by Lemma~\ref{lem:d4-fibre}(c) is a
neighbourhood of $y$ in $\mathcal B_4(m)$.  Let $y'$ be a point of it, $N\in\{2,3\}$, and $x(t)$
a polynomial curve with $x(0)=y'$ and $F(t,x(t))=O(t^N)$; such a curve exists for $N=2$ by (a).
We write $A_i(t),\dots$ for its components, $A_i,\dots$ for their values at $y'$, and
$\Phi(t)=1-M_PD_{12}(t)^3$.

\emph{Pair.}  By \eqref{eq:d4-rows}, $t^4E^z_i=R^z_i+O(t^3)$.  As $R^z_1-R^z_2=(A_1(t)-A_2(t))\Phi(t)$
and $m_1R^z_1+m_2R^z_2=m_1A_1(t)+m_2A_2(t)$, we get $\Phi=O(t^N)$ and $A_i(t)=A_i(1+a(t))+O(t^N)$
with one series $a(t)=O(t)$.

\emph{Triangle.}  The triangle rows of $F$ agree with those of $F_0$ up to $O(t^3)$, so the
triangle components of $x(t)$ solve the triangle block of $F_0$ modulo $t^N$.  By the proof of
Lemma~\ref{lem:d4-fibre}(b),(c), the zero set of that block is, near $y'$, the curve
$\beta\mapsto\beta\cdot y'$, of dimension $\dim\ker J_T$; so the triangle components of $x(t)$
agree modulo $t^N$ with a curve on it, and $\eta_k(t)=\theta(t)\eta_k+O(t^N)$ with $\theta(0)=1$.
For $N=2$ the triangle components of $[F]_2$ therefore lie in the image of $J_T$.

\emph{Crossed ties.}  They give $D_{ik}(t)^2=-1/\big(A_i(t)(\eta_k(t)-t^2B_i)\big)+O(t^N)$, hence
\[
  D_{ik}(t)^3\eta_k(t)=-\frac{D_{ik}(t)}{A_i(t)}\Big(1+t^2\frac{B_i}{\eta_k}\Big)+O(t^N),\qquad
  D_{ik}(t)=\frac{D_{ik}}{\sqrt{(1+a)\theta}}\Big(1+\frac{t^2B_i}{2\eta_k}\Big)+O(t^N),
\]
and, with $f_i(t)=\sum_km_kD_{ik}(t)^3\eta_k(t)$,
\[
  f_i(t)=-\frac{S_i/A_i}{(1+a)^{3/2}\theta^{1/2}}-\frac32\,t^2\,\frac{B_i}{A_i}\sum_km_k\frac{D_{ik}}{\eta_k}
  +O(t^N).
\]
Since $S_1/A_1=S_2/A_2$ and $B_1/A_1=B_2/A_2$, $f_1(t)-f_2(t)=-\frac32\,t^2\,\frac{B_1}{A_1}\,\Theta+O(t^N)$.

\emph{The form $v_P$.}  By \eqref{eq:d4-rows}, $v_P$ applied to the rows of $F(t,x(t))$, with the
coefficients of $v_P$ taken at $y'$, equals
\[
  \Phi(t)\big[(B_1-B_2)(A_1(t)-A_2(t))-(A_1-A_2)(B_1(t)-B_2(t))\big]
  -(A_1-A_2)\,t\,\big(f_1(t)-f_2(t)\big)-t^3\,\Delta+O(t^4),
\]
with $\Delta=(B_1-B_2)(A_1\Sigma_1-A_2\Sigma_2)-(A_1-A_2)(B_1\Sigma_1-B_2\Sigma_2)$, which vanishes
because $A_2=\rho A_1$ and $B_2=\rho B_1$.  The bracket vanishes at $t=0$ and $\Phi=O(t^N)$.

For $N=2$ this gives $v_P([F]_2)=0$.  The crossed ties are killed by every form of the
cokernel, so with the triangle part $\Omega_2(y')=0$: the curve can be corrected at order $2$ to
a curve with $F(t,x(t))=O(t^3)$.  This proves the hypothesis for $N=3$.  For $N=3$ the
displayed expression gives $v_P([F]_3)=-(A_1-A_2)\cdot\big(-\frac32\frac{B_1}{A_1}\Theta\big)$.
\end{proof}

\subsubsection{The census}\label{d4-census}
Let $m$ satisfy (23).  After renumbering $T$,
$(\sqrt{m_3},\sqrt{m_4},\sqrt{m_5})=\lambda\,(bc,\ b(b+c),\ c(b+c))$ with $\lambda,b,c>0$, as after
Lemma~\ref{lem:d1-mb}; exchanging $4$ and $5$ exchanges $b$ and $c$, and $b=c$ is the
central-body slice.  By Lemma~\ref{lem:d1-shape} and the proof of
Proposition~\ref{prop:d1-first}, the triangle offsets $\eta_k$ of a point of $\mathcal B_4(m)$
are proportional to
\[
  (o_3,o_4,o_5)=\Big(-\frac{b^2-c^2}{bc},\ \frac{c(2b+c)}{b(b+c)},\ -\frac{b(b+2c)}{c(b+c)}\Big).
\]
The rational curve
\[
  b=2(n^4+1),\qquad c=-n^4+6n^2-1,\qquad 1<n<1+\sqrt2,
\]
has $b-c=3(n^2-1)^2$, $2b+c=3(n^2+1)^2$ and $b+2c=12n^2$, and $b/c$ takes every value in
$(1,\infty)$ on it.  Put
\[
  g_3=1,\qquad g_4=\frac{i(b+c)(n^2-1)}{c(n^2+1)},\qquad g_5=\frac{(b+c)(n^2-1)}{2nb},
\]
so that $g_k^2=o_3/o_k$, and $(m_3,m_4,m_5)=\big((bc)^2,(b(b+c))^2,(c(b+c))^2\big)$.  For a sign
vector $s=(s_{13},s_{14},s_{15},s_{23},s_{24},s_{25})\in\{\pm1\}^6$ and $j\in\{1,3\}$ put
\[
  X_j=\sum_ks_{1k}\,m_kg_k^j,\qquad Y_j=\sum_ks_{2k}\,m_kg_k^j .
\]

\begin{lemma}[computer-assisted]\label{lem:d4-census}
Let $n\in(1,1+\sqrt2)$ and $\nu>0$.
\begin{alphalist}
  \item For $48$ of the $64$ sign vectors, $\nu^3X_1+iY_1\ne0$ for all such $(n,\nu)$.  For
  each of the other $16$, namely those with $s_{13}s_{24}=1$ and $s_{14}s_{23}=-1$, there is
  exactly one such pair $(n,\nu)$ with $\nu^3X_1+iY_1=0$.
  \item At each of these $16$ pairs, $\nu X_3\ne iY_3$.
\end{alphalist}
\end{lemma}
The proof of (a) is an exact computation in three steps.  For each of the $64$ sign vectors,
multiplication by $\pm4n(n^2+1)(n^4+1)(n^2-2n-1)(n^2+2n-1)$, which has no zero on
$(1,1+\sqrt2)$, turns $\nu^3X_1+iY_1$ into a polynomial in $n$ and $\nu$, the census numerator.
Exact real-root counting splits the census numerators into $48$ classes without and $16$ with a
zero; each zero class has exactly one zero, by Sturm counts on the eliminant
$\operatorname{Im}(-\bar X_1\cdot iY_1)$ with the sign condition $\nu^3>0$.
For (b), if $\nu X_3=iY_3$ and $\nu^3X_1+iY_1=0$ then $Y_1X_3^3=X_1Y_3^3$; the numerator of
$Y_1X_3^3-X_1Y_3^3$ is coprime over $\mathbb Q(i)$ to the minimal polynomial of the zero $n$ in
all $16$ cases (Appendix~\ref{app:verification}).

\subsubsection{Proof of Theorem~\ref{thm:d4}}
Suppose that $\gamma$ is a Diagram~4 germ at $m\in\Rp^5$.  By Proposition~\ref{prop:d4-cover},
$m$ satisfies (23), and $\tau\mapsto(t(\tau),x(\tau))$, with $t=\tau^q$, is a holomorphic germ in
$F^{-1}(0)$ with $t\not\equiv0$ and $x(0)=y\in\mathcal B_4(m)$.  By Lemma~\ref{lem:d4-fibre}(c),
Lemma~\ref{lem:d4-ls} applies at $y$ with $Z=\mathcal B_4(m)$.  For $N=1$ it gives
$\Omega_1(y)=0$, that is, $f_1=f_2$ at $y$ (Lemma~\ref{lem:d4-obstr}(a)).  Then
Lemma~\ref{lem:d4-obstr}(b) and Lemma~\ref{lem:d4-ls} for $N=3$ give $\Omega_3(y)=0$; as
$(A_1-A_2)B_1/A_1\ne0$, $\Theta=0$ at $y$.

Both conditions are symmetric in the bodies of $T$ and homogeneous under the scalings, and they
involve $m_3,m_4,m_5$ linearly.  So we may renumber $T$ as in \S\ref{d4-census} with $b>c$
($b\ne c$ because $\eta_k(y)\ne0$ by Proposition~\ref{prop:d4-cover} and $o_3=0$ for $b=c$),
choose $n\in(1,1+\sqrt2)$ with $b/c$ equal to the ratio of the curve, replace $(m_3,m_4,m_5)$ by
the values on the curve, and apply the scalings so that $A_1=1$ and $\eta_3=-1$.  Then
$\eta_k=-o_k/o_3$ and $D_{1k}^2=-1/\eta_k=g_k^2$; with $\nu=(m_1/m_2)^{1/2}$ we have $A_2=-\nu^2$ and
$D_{2k}^2=-D_{1k}^2/\nu^2$.  Hence $D_{1k}=s_{1k}g_k$ and $D_{2k}=s_{2k}\,ig_k/\nu$ for a sign
vector $s$, and $S_1=X_1$, $S_2=iY_1/\nu$, $\Sigma_1=X_3$, $\Sigma_2=-iY_3/\nu^3$.  The condition
$S_1/A_1=S_2/A_2$ becomes $\nu^3X_1+iY_1=0$, and $\Theta=A_2\Sigma_2-A_1\Sigma_1=0$ becomes
$\nu X_3=iY_3$.  This contradicts Lemma~\ref{lem:d4-census}. \qed

\begin{remark}\label{rem:d4-witnesses}
(a) Positivity of the masses enters through (23) and the parametrization of the triangle,
through $m_1^3+m_2^3\ne0$, and through $\nu>0$; no reality of the germ is used, so
Theorem~\ref{thm:d4} excludes all complex germs.
(b) The two obstructions involve the outer masses only through $\rho=-m_1/m_2$, and not through
the ratio $M_P/M_T$.
\end{remark}


\subsection{Diagram 5}\label{sec:d5}

Diagram~5 is the isolated $zw$-triangle $345$ alone: its only strokes are the six strokes of
the three triangle edges, and no vertex is circled \cite[\S7.4 and Fig.~11]{AK}.
We write $T=\{3,4,5\}$ and call $1$ and $2$ the \emph{outer bodies}.  For this diagram
\cite[\S\S8.4--8.5]{AK} record only relation (23), and \cite[\S9]{AK} avoid estimating the
distances that carry no edge.
We build no chart.  We classify the orders of the outer bodies relative to the triangle and
decide by exact linear programming which are compatible with the equations on a curve through
a real point, where $U=I$ has a nonzero constant value.

\begin{theorem}\label{thm:d5}
Let $m\in\Rp^5$.  No Diagram 5 germ with an unbounded four-product lies on a curve through a
real point: if an irreducible curve $\Gamma\subset A_m$ carries such a germ, then $\Gamma$
contains no real point.
\end{theorem}

This is part~(b) of Theorem~\ref{thm:exclusion} for $r=5$.  Germs with all four-products bounded
are not excluded, and need not be (Proposition~\ref{prop:reduction}); they do occur on curves
through real points, e.g.\ at the collision end of the modified Roberts continuum at the masses
$(4,4,4,4,1)$ \cite[p.~543]{AK}, which is real but not positive.

\subsubsection{The frame}\label{d5-frame}
Let $\gamma$ be a Diagram 5 germ, numbered as in \cite[Fig.~11]{AK} and renormalized as in
Lemma~\ref{lem:germ}, with orders measured in units of $\ord\epsilon$, so that $\mathcal Z$ and
$\mathcal W$ have least order $-2$.  After a constant reparametrization the Puiseux parameter
$t$ has $\epsilon\approx|t|$, and we write orders in $t$.
By Fact~\ref{fact:rules}: $-2\le\ord r_{kl}\le1$, with $\ord r_{kl}=1$ exactly for $zw$-edges, and
$-2\le\ord z_{kl},\ord w_{kl}\le2$ (Estimate~1); $\ord z_{kl}=\ord w_{kl}=1$ and
$\ord\delta_{kl}=-1$ on $T$ (Estimate~2); $\ord z_k,\ord w_k>-2$ for all $k$, as nothing is circled
(Rule~1a at the outer vertices; on $T$ the strokes of the $zw$-edges are not maximal by
Estimates~1 and~2, so Rule~1e applies); and $\ord Z_{kl},\ord W_{kl}>-2$
unless $k,l\in T$.

Put $M_T=m_3+m_4+m_5$, $M=\sum_km_k$, $z_c=\sum_{k\in T}m_kz_k/M_T$, $w_c$ likewise,
$\xi_k=z_k-z_c$ and $\eta_k=w_k-w_c$ ($k\in T$), and let
$F_j=\sum_{k\in T}m_k\delta_{jk}^3(z_j-z_k)$ and $G_j=\sum_{k\in T}m_k\delta_{jk}^3(w_j-w_k)$ be the
forces of the triangle on the outer body $j$.  By Proposition~\ref{prop:d45-first}, (23) holds
and $(\xi/t,\eta/t)$ tends to a zero-multiplier triangle.  Renumbering $T$,
$(\sqrt{m_3},\sqrt{m_4},\sqrt{m_5})=\lambda(b,b(b+1),b+1)$ with $b,\lambda>0$.
This triangle has real, mass-centred, pairwise distinct offsets $s_3,s_4,s_5$ with distinct
cubes; $s_4,s_5\ne0$, and $s_3=0$ exactly for $b=1$, the \emph{central-body slice} $m_4=m_5=4m_3$.
Its $z$- and $w$-shapes agree, and it is unique up to translations, the scalings
$(\xi,\eta)\mapsto(a\xi,a'\eta)$ and the sign of its inverse distances.  A constant
reparametrization and a constant renormalization $(z,w)\mapsto(az,a^{-1}w)$, which preserves
rows, ties, orders, the $r_{kl}^2$, $U$ and $I$, give
\begin{equation}\label{d5-eq-offsets}
  \xi_k=t\,(s_k+\sigma_k),\qquad \eta_k=t\,(s_k+\tau_k),\qquad \sigma_k,\tau_k\to0.
\end{equation}

For $j\in\{1,2\}$ let $\alpha_{jc},\beta_{jc}>-2$ be the orders and $\hat z_j,\hat w_j$ the leading
coefficients of $z_j-z_c$, $w_j-w_c$; for $P\in\{12,1c,2c\}$ let $\alpha_P,\beta_P$ be the orders of
the $z$- and $w$-differences, $r_P^2$ their product and $\rho_P=\frac12(\alpha_P+\beta_P)=\ord r_P$.
The \emph{pattern} of $\gamma$ is the word $X_{12}X_{1c}X_{2c}$, where $X_P$ is C, R or E as
$\rho_P>0$, $=0$ or $<0$.

\subsubsection{The reality input}\label{d5-reality}
\begin{lemma}\label{d5-lem-reality}
Let $\Gamma\subset A_m$ be an irreducible curve containing a real point $q$.  Then $U=I=I_0$ on
$\Gamma$, where $I_0=\sum_km_k|z_k(q)|^2>0$, and along every Diagram 5 germ of $\Gamma$, in the
frame above,
\begin{equation}\label{d5-eq-inertia}
  I_T+I_3=I_0,\qquad I_T=\sum_{k\in T}m_k\xi_k\eta_k,\qquad
  I_3=\frac{m_1m_2r_{12}^2+m_1M_Tr_{1c}^2+m_2M_Tr_{2c}^2}{M},
\end{equation}
with $\ord I_T=2$ and hence $\ord I_3=0$.
\end{lemma}

\begin{proof}
By Lemma~\ref{lem:U-constant}, $U$ is constant on $\Gamma$, and $U=I$ on $A_m$.  At $q$,
$w_k=\bar z_k$ and $z_1(q)\ne z_2(q)$ (tie of $12$), so $I(q)>0$.  The renormalization preserves
$U$ and $I$.  As $\sum_km_kz_k=\sum_km_kw_k=0$ on $A_m$, $I-I_T-I_3$ vanishes identically, and
$I_T=t^2(\sum_km_ks_k^2+o(1))$ with $\sum_km_ks_k^2>0$ by \eqref{d5-eq-offsets}.
\end{proof}

This is the only input from reality.  It enters through the conditions (A) and (B) of
\S\ref{d5-census}, that two quantities have order exactly $0$; so only $I_0\ne0$ is used, not its
sign, nor $\sigma(\Gamma)=\Gamma$.

\subsubsection{Regimes of an outer body}\label{d5-regimes}
\begin{definition}\label{d5-def-regimes}
Let $j\in\{1,2\}$, $\alpha=\alpha_{jc}$ and $\beta=\beta_{jc}$.  The outer body $j$ is
\begin{alphalist}
\item \emph{far} if $\alpha,\beta<1$.  Then $\delta_{jk}/\delta_{jc}\to\varepsilon_{jk}=\pm1$ for a
fixed branch $\delta_{jc}$ of $((z_j-z_c)(w_j-w_c))^{-1/2}$.  The body is a \emph{monopole} if
$\mu_j=\sum_km_k\varepsilon_{jk}\ne0$.  Otherwise it is \emph{dipole-$z$} if $\alpha>\beta$,
\emph{dipole-$w$} if $\beta>\alpha$, and, if $\alpha=\beta$, \emph{dipole-equal-$F$},
\emph{dipole-equal-$G$} or \emph{dipole-equal} as $\hat w_j+3\hat z_j$, $3\hat w_j+\hat z_j$ or
neither vanishes, the leading coefficients of $F_j$, $G_j$ at the monopole orders raised by
$1-\alpha$ being nonzero multiples of these.
\item \emph{near in $z$} if $\alpha\ge1>\beta$.  For $\alpha=1$ let
$\theta_j=\lim t^{-1}(z_j-z_c)$.  The body is \emph{$z$-dominant-coincidence} if $\alpha=1$ and
$\theta_j=s_k$, its \emph{special body} $k$, so that $\ord(z_j-z_k)=1+e_j$ with $e_j>0$;
\emph{$z$-coincident} if $\alpha=1$ and $\theta_j\notin\{s_3,s_4,s_5\}$, with the variants
\emph{$z$-coincident-$\Phi$}, \emph{-$\Psi$}, \emph{-both} if $\Phi_j=0\ne\Psi_j$,
$\Psi_j=0\ne\Phi_j$, $\Phi_j=\Psi_j=0$, where $\Phi_j=\sum_km_ky_{jk}$, $\Psi_j=\sum_km_ky_{jk}^3$,
$\delta_{jk}=y_{jk}\nu_j(1+o(1))$, $y_{jk}^{-2}=\theta_j-s_k$ and $\nu_j$ is a branch of
$(t(w_j-w_c))^{-1/2}$; \emph{$z$-inner} if $\alpha>1$ off the slice; and
\emph{$z$-dominant-central} if $\alpha>1$ on the slice, with special body $3$,
$1+e_j=\ord(z_j-z_3)$ and $g=\ord\xi_3$.
\item \emph{near in $w$} if $\beta\ge1>\alpha$, with the seven mirror regimes.
\end{alphalist}
\end{definition}

\begin{lemma}\label{d5-lem-complete}
Let $m\in\Rp^5$.  Each outer body lies in exactly one of these $20$ regimes, and
$z$-coincident-both and $w$-coincident-both are empty.  The dipoles need $m_a=m_b+m_c$ for the
minority body $a$, which fails on the slice; and for a $z$-dominant-central body the least of
$\alpha_{jc}$, $1+e_j$, $g$ is attained at least twice.
\end{lemma}

\begin{proof}
For $k\in T$, $z_j-z_k=(z_j-z_c)-\xi_k$ and $w_j-w_k=(w_j-w_c)-\eta_k$, with
$\ord\xi_k,\ord\eta_k\ge1$.  If $\alpha,\beta\ge1$, then $\ord r_{jk}\ge1$ and $jk$ would be a
$zw$-edge.  If $\alpha,\beta<1$, then $\delta_{jk}^2/\delta_{jc}^2\to1$; $\mu_j=0$ needs a single
minority body $a$ with $m_a=m_b+m_c$, false for the slice masses $(1,4,4)$; and
$\hat w_j+3\hat z_j=3\hat w_j+\hat z_j=0$ would give $\hat z_j=\hat w_j=0$.  Let $\alpha=1>\beta$:
at most one $s_k$ equals $\theta_j\ne0$, and otherwise every $z_j-z_k$ has order $1$; at positive
masses on (23), $\Phi_j=\Psi_j=0$ forces $y=0$, i.e.\ $\theta_j=\infty$.  Let $\alpha>1>\beta$: then
$z_j-z_k=-ts_k(1+o(1))$ when $s_k\ne0$; on the slice $\ord(z_j-z_3)\ge\min(\alpha,g)>1$, and
$z_j-z_c=(z_j-z_3)+\xi_3$ gives the tropical condition, including the deep case
$\alpha=g<1+e_j$.  The rest is symmetric.
\end{proof}

If $z_j\equiv z_c$ or $\xi_3\equiv0$, the order $\alpha_{jc}$ or $g$ is infinite.  Such an order
enters the constraints below only through conditions that remain true when it is replaced by any
sufficiently large finite value, and the census, whose variables are free, includes such values.
The order
$\alpha_{12}$ is always finite, as the tie of $12$ gives $z_1\ne z_2$.  Up to the exchanges
$1\leftrightarrow2$ and $z\leftrightarrow w$, which permute patterns and four-products, there are
$111$ pairs of regimes.

\subsubsection{Identities and orders}\label{d5-census}
Let $\Pi=\delta_{12}^3(z_1-z_2)$, $U_T=\sum_{k<l\in T}m_km_l\delta_{kl}$ and, summing over
$j\in\{1,2\}$, $k\in T$, $V_1=\sum m_jm_k\delta_{jk}^3\xi_k(w_k-w_j)$,
$V_2=\sum m_jm_k\delta_{jk}^3\eta_k(z_k-z_j)$.  The census uses the conditions below; the identities
in (R), (T), the virial identity $I_T=U_T+\frac12(V_1+V_2)$ and the pair-centre identity of
Proposition~\ref{d5-prop-over} are exact combinations of rows and ties, as a symbolic computation
shows (Appendix~\ref{app:verification}).
\begin{itemize}
\item[(R)] $z_1-z_2=(m_1+m_2)\Pi+F_1-F_2$, $z_1-z_c=m_2\Pi+(1+\frac{m_1}{M_T})F_1+\frac{m_2}{M_T}F_2$
and $z_2-z_c=-m_1\Pi+\frac{m_1}{M_T}F_1+(1+\frac{m_2}{M_T})F_2$, with their $w$-analogues; the three
forms in $(\Pi,F_1,F_2)$ are pairwise independent on each coordinate plane, with common kernel
$(1,-m_2,m_1)$;
\item[(C)] $z_1-z_2=(z_1-z_c)-(z_2-z_c)$, and likewise in $w$;
\item[(T)] $V_1=V_2$: the outer bodies exert no torque on the triangle about its centre;
\item[(A)] $\ord I_3=0$ (Lemma~\ref{d5-lem-reality});
\item[(B)] $\sum_{j,k}m_jm_k\delta_{jk}\bigl(1+\frac12\bigl(\frac{\xi_k}{z_j-z_k}+\frac{\eta_k}{w_j-w_k}\bigr)\bigr)
+m_1m_2\delta_{12}=I_0-I_T$, of order exactly $0$, from $U=I_0$ and the virial identity;
\item[(V)] $V_1=V_2=I_T-U_T$, where $\ord I_T=2$ and $\ord U_T\ge\min(5,-1+2\kappa)$ by
\eqref{d5-eq-UT};
\item[(M)] for two outer bodies near in the same colour, the relations of
Lemma~\ref{d5-lem-consequences}.
\end{itemize}

Each regime fixes the orders of $F_j$, $G_j$, of the shares of body $j$ in the torque
$V_1-V_2$, in (B), $V_1$ and $V_2$, of its forcings of the triangle bodies and of the $r_{jk}^2$:
exactly with nonzero leading coefficient, by lower bounds, by the vanishing of the generic
coefficient (of $F_j$ in dipole-equal-$F$ and $z$-coincident-$\Phi$, of $G_j$ in dipole-equal-$G$
and $z$-coincident-$\Psi$), or not at all.  Nonvanishing comes from $\mu_j\ne0$, from the dipole
moment $\sum_km_k\varepsilon_{jk}\xi_k=-2m_a\xi_a$ ($a$ the body of minority sign; of order $1$ unless
$a$ is central), from $\sum_km_ks_k^2>0$, from the definitions, from $s_k\ne0$ at a dominant special
body, and, for $z$-inner, from the following observation: off the slice a single $s_k$ has the
minority sign, and as $y_{jk}^{-2}=-s_k$, its term alone forms the real or the imaginary part of
$\Phi_j$ and of $\Psi_j$.  A second computation re-derives these orders in the $12$ regimes of the $z$-side ($w$ by
symmetry) from exact first-order expansions with remainder bounds, and checks them, every census
constraint and the formulas for distances and four-products by exact linear programming over each
regime; it finds one exception (\S\ref{d5-slice}).  Valid constraints that the census omits, such
as $\alpha_{jc},\beta_{jc}>-2$, only weaken it.

A vanishing sum of Puiseux series attains its least order twice; a sum of order exactly $0$ has
least order $\le0$, attained twice if negative; and in (R), at the least order of $\Pi,F_1,F_2$, at
most one equation cancels if two of them attain it, and none, one or all three if all do.  So
each identity becomes a finite disjunction of rational linear systems, with strict and non-strict
inequalities, in the orders.  The census explores the disjuncts depth first and decides each
system exactly over $\mathbb Q$ by the simplex method, with free variables and a common slack
$0<s\le1$ for the strict inequalities.

\subsubsection{The Morse--Bott lemma}\label{d5-mb}
In reduced coordinates $\hat x_k=\xi_k/t$, $\hat y_k=\eta_k/t$, $D_{kl}=t\delta_{kl}$ ($k,l\in T$)
we have $U_T=t^{-1}\hat U$, $\hat U=\sum_{k<l}m_km_lD_{kl}$, with
$D_{kl}=((\hat x_k-\hat x_l)(\hat y_k-\hat y_l))^{-1/2}$ on the branch of the limit.  For
$\psi_k=\sum_lm_lD_{kl}^3(\hat y_k-\hat y_l)$ and $\phi_k=\sum_lm_lD_{kl}^3(\hat x_k-\hat x_l)$,
$\partial\hat U/\partial\hat x_k=-\frac{m_k}2\psi_k$, $\partial\hat U/\partial\hat y_k=-\frac{m_k}2\phi_k$
and, by homogeneity, $\hat U=\sum_km_k\hat x_k\psi_k=\sum_km_k\hat y_k\phi_k$.  The triangle rows
read exactly
\begin{equation}\label{d5-eq-rows}
  \psi_k=t^3\hat y_k-\hat g_k,\qquad \phi_k=t^3\hat x_k-\hat g^z_k,
\end{equation}
where $g_k=t^2\sum_jm_j\delta_{jk}^3(w_k-w_j)$ is the crossed forcing of the $w$-row of $k$,
$g^z_k$ its $z$-analogue, and the hat is the deviation from the $m$-weighted mean over $T$.  Let
$O=\{(as+u\mathbf 1,a's+v\mathbf 1)\}$ be the orbit of the zero-multiplier triangle under
translations and scalings; the reduced triangle tends to $x_0=(s,s)\in O$.

\begin{lemma}[computer-assisted]\label{d5-lem-mb}
At each point $(as,a's)$ of $O$, $\hat U=0$, $d\hat U=0$, the translations and scalings lie in
the Hessian kernel, and on $(\hat x_3,\hat y_3)$ the Hessian $H$ has entries $3X/(\varrho a^2)$,
$X/(\varrho aa')$, $3X/(\varrho a'^2)$, where $\varrho^2=aa'$ and $X=b^4(b+1)^6/4$ for masses
normalized to $b^2,b^2(b+1)^2,(b+1)^2$; its determinant is $8X^2/\varrho^6$, and
$(H^{-1})_{\hat x\hat x}=3a^2\varrho/(8X)\ne0$.
\end{lemma}
The proof is an exact computation in the field $\mathbb Q(b,a,\varrho)$ of rational functions
(Appendix~\ref{app:verification}).

\begin{lemma}[Morse--Bott]\label{d5-lem-mb2}
If $x(t)\to x_0$ is a germ of reduced triangles, $\kappa_\psi=\ord\psi(x(t))$ and
$\kappa_\phi=\ord\phi(x(t))$, then $\ord\hat U(x(t))\ge2\min(\kappa_\psi,\kappa_\phi)$, and
$\kappa_\psi<\kappa_\phi$ implies $\kappa_\phi\le2\kappa_\psi$ (and symmetrically).
\end{lemma}

\begin{proof}
As $s_4\ne s_5$, the $\hat x_3$- and $\hat y_3$-directions span a complement of the tangent space of
$O$, which by Lemma~\ref{d5-lem-mb} is thus the Hessian kernel; write points near $x_0$ as
$o(u)+v_1e^x_3+v_2e^y_3$, $o(u)\in O$.  Taylor's formula gives $\hat U=\frac12v^TH(u)v+O(|v|^3)$ and
$\partial_v\hat U=H(u)v+O(|v|^2)=-\frac{m_3}2(\psi_3,\phi_3)$ with $H(u)$ invertible, so
$\ord\hat U\ge2\ord v=2\ord\partial_v\hat U\ge2\min(\kappa_\psi,\kappa_\phi)$.  If
$\kappa_\psi<\kappa_\phi$, the leading coefficient of $\psi$ satisfies $\sum_km_k\psi_k=0$ and, as
$\hat U=\sum_km_k\hat x_k\psi_k$ has order $>\kappa_\psi$, $\sum_km_ks_k\psi_k=0$; as $s_4\ne s_5$,
its third component is nonzero, so $\partial_v\hat U=t^{\kappa_\psi}(\chi,0)+o(t^{\kappa_\psi})$ with
$\chi\ne0$ and $\hat U=\frac12(\partial_v\hat U)^TH^{-1}\partial_v\hat U+O(|v|^3)
=\frac12(H^{-1})_{\hat x\hat x}\chi^2t^{2\kappa_\psi}+o(t^{2\kappa_\psi})$.  But
$\hat U=\sum_km_k\hat y_k\phi_k$ has order $\ge\kappa_\phi$.
\end{proof}

By Lemma~\ref{d5-lem-mb2}, as the self-terms in \eqref{d5-eq-rows} have order $3$, along every
Diagram~5 germ
\begin{equation}\label{d5-eq-UT}
  \ord U_T\ge\min(5,-1+2\kappa),\qquad \kappa=\min(\ord\hat g,\ord\hat g^z).
\end{equation}
\begin{corollary}[uniform forcing]\label{d5-cor-uniform}
Let both outer bodies be near in $z$, and let $\kappa_*$ be the least order of the terms
$t^2m_j\delta_{jk}^3(w_k-w_j)$ ($j\in\{1,2\}$, $k\in T$).  Then the coefficient of $t^{\kappa_*}$ in
$g_k$ is the same for $k=3,4,5$; likewise with $z$ and $w$ exchanged.
\end{corollary}

\begin{proof}
Let $\ord(z_j-z_k)=1+e_{jk}$, $e_{jk}\ge0$.  The term of body $j$ in $g_k$ has order
$\kappa_{jk}=\frac12(1-\beta_{jc}-3e_{jk})$, in $g^z_k$ the order
$\frac12(3-3\beta_{jc}-e_{jk})=2\kappa_{jk}+\frac12(1-\beta_{jc}+5e_{jk})$.  As $\beta_{jc}>-2$ and
$e_{jk}=0$ for some $k$, $\kappa_*<\frac32$.  Unequal coefficients would give $\kappa_\psi=\kappa_*$
(as $3>\kappa_*$), while $\hat g^z$ has order $>2\kappa_*$ (as $\beta_{jc}<1$) and $3>2\kappa_*$; so
$\kappa_\phi>2\kappa_\psi$, with $\kappa_\psi>0$ as $\psi\to0$, against Lemma~\ref{d5-lem-mb2}.
\end{proof}

\begin{lemma}[computer-assisted]\label{d5-lem-consequences}
Let both outer bodies be near in $z$ (or, mirrored, in $w$).  (a) A forcing from one source,
$k\mapsto c(\theta-s_k)^{-3/2}$ with $c\ne0$, or supported on fewer than three triangle bodies, is
not uniform.  (b) Two $z$-dominant-coincidence bodies share their special body; so
$\alpha_{12}>1$.  (c) A $z$-coincident-$\Phi$ or $z$-coincident-$\Psi$ body does not occur with a
$z$-dominant-coincidence body.  (d) A $z$-inner body does not occur with a $z$-dominant-central
body.  (e) Two $z$-dominant-central bodies have the same $g$; their forcings of body $3$, of orders
$\varphi_j=\ord(\delta_{j3}^3(w_3-w_j))=-\frac32(1+e_j)-\frac12\beta_{jc}$, tie and cancel;
$g>3+\varphi_1$; and $1+e_j=g<\alpha_{jc}$ with $\beta_{jc}<-g$ for both $j$ is impossible.
\end{lemma}
The lemma follows from Corollary~\ref{d5-cor-uniform} by exact computations:
(a) through the resultant $9(s_4-s_5)^2\cdot\frac12\sum_{k<l}(s_k-s_l)^2\ne0$ of
$(\theta-s_3)^3-(\theta-s_4)^3$ and $(\theta-s_3)^3-(\theta-s_5)^3$, (c) through twelve gcd
computations over $\mathbb Q(\omega)$ at the points $(s_k-\omega s_{k'})/(1-\omega)$,
$\omega^3=1\ne\omega$, followed by exact real-root counting in $b$ (Appendix~\ref{app:verification});
(d) is the definition, and the last case of (e) falls as in Proposition~\ref{d5-prop-over}.

\begin{proof}[Proof of the tie and of $g>3+\varphi_1$ in (e)]
The term of body $j$ in $g_k$ has order $2+\varphi_j$ for $k=3$, and order
$\frac12-\frac12\beta_{jc}>2+\varphi_j$ for $k=4,5$, where $\ord(z_j-z_k)=1$ and $e_j>0$.  So the
leading crossed forcing is supported on body $3$, and by Corollary~\ref{d5-cor-uniform} and (a)
the two forcings of body $3$ have one order $\varphi:=\varphi_1=\varphi_2$ and cancel; hence
$\ord\hat g_3>2+\varphi$.  The $z$-forcing of body $3$ by body $j$ has order
$2+\varphi+(1+e_j)-\beta_{jc}>2+\varphi$, as $\beta_{jc}<1$, and the self-terms in
\eqref{d5-eq-rows} have order $\ge3>2+\varphi$, as $\varphi<-\frac12$ by $\beta_{jc}>-2$.  So
$\ord\psi_3,\ord\phi_3>2+\varphi$.  In the proof of Lemma~\ref{d5-lem-mb2},
$\partial_v\hat U=H(u)v+O(|v|^2)=-\frac{m_3}2(\psi_3,\phi_3)$ with $H(u)$ invertible and $v\to0$,
so $\ord v\ge\min(\ord\psi_3,\ord\phi_3)>2+\varphi$.  On the slice $s_3=0$, so the mass-centred
point $o(u)$ of the orbit has third $\hat x$-component $0$, and mass-centring gives
$\hat x_3=(1-m_3/M_T)v_1$.  Therefore $g=\ord\xi_3=1+\ord\hat x_3\ge1+\ord v>3+\varphi$.
\end{proof}

\subsubsection{The census and over-rotation}\label{d5-result}

\begin{lemma}[computer-assisted]\label{d5-lem-census}
For each of the $111$ pairs of regimes and each of the $27$ patterns, the conditions of the frame
and of the two regimes together with (R), (C), (T), (A), (B), (V), (M) are infeasible, except for
the pairs and patterns of Table~\ref{d5-tab-survivors} and for pairs containing an empty regime.
For the five monopole pairs of \S\ref{d5-slice} this holds with the rows corrected there.
\end{lemma}
The proof is a finite computation: for each of the $111$ pairs of regimes and each of the $27$
patterns, the finite disjunction of rational linear systems obtained above is decided by exact
linear programming over $\mathbb Q$, with the orders checked by the second computation and, for
the five monopole pairs, with the corrected rows (Appendix~\ref{app:verification}).

\begin{table}[ht]
\centering
\small
\begin{tabular}{p{8.9cm}p{3.3cm}}
\hline
pair of regimes (up to the exchanges) & surviving patterns \\
\hline
(monopole, monopole); (monopole, $X$) for $X$ = dipole-$z$, dipole-equal, dipole-equal-$F$;
(dipole-$z$, dipole-$z$); (dipole-equal, $X$) for $X$ = dipole-equal, dipole-equal-$F$;
(dipole-equal-$F$, $X$) for $X$ = dipole-equal-$F$, dipole-equal-$G$ & RRR \\
($X$, $Y$) for $X,Y\in\{z$-coincident, $z$-coincident-$\Phi$, $z$-coincident-$\Psi\}$ & RRR \\
($z$-inner, $z$-inner) & RRR \\
($z$-inner, $X$) for $X\in\{z$-coincident, $z$-coincident-$\Phi$, $z$-coincident-$\Psi\}$ & RCR \\
($z$-dominant-coincidence, $z$-dominant-coincidence) & CEE, CRR \\
($z$-dominant-central, $z$-dominant-central) & RRR, RRC, RCR, CEE, CRR \\
every other pair of nonempty regimes & none \\
\hline
\end{tabular}
\caption{The census of Lemma~\ref{d5-lem-census}.  A pattern lists the signs of $\rho_{12}$,
$\rho_{1c}$, $\rho_{2c}$: C positive, R zero, E negative.}\label{d5-tab-survivors}
\end{table}

Each four-product contains a triangle edge, of order $2$, and its order is the sum of its edge
orders: $\alpha_{12}+\beta_{12}$ for $r_{12}^2$; for $r_{jk}^2$, $\alpha_{jc}+\beta_{jc}$ if $j$ is far
and $1+\beta_{jc}$ (or $1+e_j+\beta_{jc}$ at the special body) if $j$ is near in $z$.

\begin{lemma}[computer-assisted]\label{d5-lem-four}
For every pair and pattern of Table~\ref{d5-tab-survivors} other than ($z$-inner, $z$-inner) with
RRR, adding $\ord p_{ab}<0$ for any one of the ten four-products makes the constraints infeasible.
For ($z$-inner, $z$-inner) with RRR exactly the six four-products $p_{jk}$, $j\in\{1,2\}$, $k\in T$,
can have negative order.
\end{lemma}
The proof is a computation of the same kind: for each surviving pair and pattern and each of the
ten four-products, $260$ cases in all, we add $\ord p_{ab}<0$, with $\ord p_{ab}$ expressed through
the edge orders as above, and decide the resulting systems by exact linear programming over
$\mathbb Q$ (Appendix~\ref{app:verification}).

By Lemmas~\ref{d5-lem-census} and~\ref{d5-lem-four} (and \S\ref{d5-slice} for pairs with a
monopole), a germ with an unbounded four-product would be \emph{over-rotated}: both outer bodies
inner in one colour, say $z$, and all $\rho_P=0$, so that $\beta_{jc}=-\alpha_{jc}<-1$ and the
crossed distances, of order $\frac12(1+\beta_{jc})<0$, are unbounded.

\begin{proposition}\label{d5-prop-over}
Let $m\in\Rp^5$.  No Diagram 5 germ has both outer bodies $z$-inner and pattern RRR, nor both
$w$-inner and pattern RRR.
\end{proposition}

\begin{proof}
Here $\beta_{jc}=-\alpha_{jc}<-1$ and every $z_j-z_k$ ($k\in T$) has order $1$; as
$\delta_{jk}^3(w_j-w_k)=\delta_{jk}/(z_j-z_k)$,
$\ord G_j\ge-\frac32-\frac12\beta_{jc}=\beta_{jc}+\frac32(\alpha_{jc}-1)>\beta_{jc}$.  By (R) in $w$, the
pair-centre identity $m_1(w_1-w_c)+m_2(w_2-w_c)=\frac M{M_T}(m_1G_1+m_2G_2)$ holds.  If
$\beta_{1c}<\beta_{2c}$ (the other case is symmetric), its left side has order $-\alpha_{1c}$ and its
right side order at least
$-\frac32+\frac12\alpha_{2c}$, larger by $(\alpha_{1c}-1)+\frac12(\alpha_{2c}-1)>0$.  So
$\beta_{1c}=\beta_{2c}=:\beta$ and $m_1\hat w_1+m_2\hat w_2=0$.  Next, $z_k-z_j=ts_k(1+o(1))$ and
$w_k-w_j=-\hat w_jt^\beta(1+o(1))$, so
$(t^2m_j\delta_{jk}^3(w_k-w_j))^2=-(m_j^2/\hat w_j)s_k^{-3}t^{1-\beta}(1+o(1))$, and the coefficient of
$t^{(1-\beta)/2}$ in $g_k$ is $(c_1\varepsilon'_{1k}+c_2\varepsilon'_{2k})s_k^{-3/2}$, with fixed
branches, $c_j^2=-m_j^2/\hat w_j$ and signs $\varepsilon'_{jk}$.  By Corollary~\ref{d5-cor-uniform}
it is one value $f$ for $k=3,4,5$.  If $f\ne0$, then $f^2s_k^3=(c_1\varepsilon'_{1k}+c_2\varepsilon'_{2k})^2$
depends only on $\varepsilon'_{1k}\varepsilon'_{2k}$, so two of the cubes $s_k^3$ coincide, which is
false.  Hence $f=0$ and $c_1^2=c_2^2$, that is, $m_1^2\hat w_2=m_2^2\hat w_1$.  The two linear
relations for $(\hat w_1,\hat w_2)$ have determinant $m_1^3+m_2^3>0$, so $\hat w_1=\hat w_2=0$,
which is absurd.
\end{proof}
No reality input is used here.

\subsubsection{The central-body slice}\label{d5-slice}
On the slice body $3$ sits at the centre of the zero-multiplier triangle; this gives the
dominant-central regimes, voids $z$-inner, and causes the one exception found by the second
computation: for a monopole whose minority crossed sign is on body $3$ the dipole moment
$-2m_3\xi_3$ has order $g>1$, so for its shares in $V_1$, $V_2$ the orders the census uses as exact
are only lower bounds (the coefficients there vanish).  Its other rows are those of every
monopole: the leading coefficients of $F_j$, $G_j$ and of its term in (B) are multiples of
$\mu_j\ne0$, and its torque row is a lower bound.  The correction can matter only for a monopole
facing a monopole, a $z$-coincident (or -$\Phi$, -$\Psi$) or a $z$-dominant-coincidence body
(elsewhere the masses avoid the slice, the regime is void on it, or there is no virial row);
recomputed with it, these five pairs keep their surviving patterns.

\begin{remark}\label{d5-lem-ablation}
Suppose that we drop (T), (B), (A), (V), and keep of (M) only its exact part ($z$-inner is void
on the slice; two $z$-dominant-central bodies share $g$ and satisfy the tropical condition of
$z_1-z_2=(z_1-z_3)-(z_2-z_3)$).  Then, among all pairs of nonempty regimes, patterns and
four-products, a four-product of negative order is feasible only for ($z$-inner, $z$-inner) with
RRR, ($z$-dominant-coincidence, $z$-dominant-coincidence) with EEE or REE, and
($z$-dominant-central, $z$-dominant-central) with RRR, RRC, RCR, CEE, CRR or CCC.  This shows
where (A), (B) and the Morse--Bott relations are needed.  None of these pairs contains a monopole,
so the slice exception cannot create an unbounded four-product even before the correction.  We do
not use this.
\end{remark}

\subsubsection{Proof of Theorem~\ref{thm:d5}}
Let $\gamma$ be a Diagram 5 germ with an unbounded four-product on an irreducible curve
$\Gamma\subset A_m$ containing a real point.  By Lemma~\ref{d5-lem-reality}, (A) and (B) hold along
$\gamma$; the other conditions of \S\ref{d5-census} hold along every germ, and all of them are
necessary conditions on the orders.  By Lemma~\ref{d5-lem-complete} and the exchanges, the outer
bodies form one of the $111$ pairs of regimes, both nonempty, and some $p_{ab}$ has negative order
(Definition~\ref{def:germ}).  If the pair contains a monopole, Lemma~\ref{d5-lem-census}, with the
rows of \S\ref{d5-slice} for the five slice-sensitive pairs, leaves only pairs of two far bodies
with pattern RRR (Table~\ref{d5-tab-survivors}).  Then $\rho_{12}=\rho_{1c}=\rho_{2c}=0$, so
$r_{12}^2$ and every $r_{jk}^2$ with $j\in\{1,2\}$, $k\in T$ have order $0$, and every four-product,
which contains a triangle edge of order $2$, has order at least $2$: a contradiction.  Otherwise
every order used by the census is confirmed by the second computation, and
Lemmas~\ref{d5-lem-census} and~\ref{d5-lem-four} leave only ($z$-inner, $z$-inner) with RRR or
its mirror, which Proposition~\ref{d5-prop-over} excludes.  \qed

\begin{remark}\label{d5-rem-h1}
We do not prove the \emph{outer-triple regularity} of Diagram~5 germs, that is, that their pattern
is RRR: the census leaves RCR, CEE, CRR and, on the slice, RRC, RCR, CEE, CRR.
Theorem~\ref{thm:d5} does not use it.
\end{remark}



\section{Diagram 6}\label{sec:d6}

\subsection{The diagram and the butterfly chart}

Diagram~6 is the \emph{butterfly} of \cite[\S7.4]{AK}, a $z$-triangle and a $w$-triangle (the
\emph{wings}) with one common vertex, with both wings circled.  In the numbering of
\cite[\S8.6]{AK} the $z$-wing $\{1,2,5\}$ has the $z$-strokes $12,15,25$ and $z$-circles at $1,2$;
the $w$-wing $\{3,4,5\}$ has the $w$-strokes $34,35,45$ and $w$-circles at $3,4$; vertex $5$ is
not circled, and no stroke joins $\{1,2\}$ to $\{3,4\}$.
The butterfly has three circlings, Diagrams~6--8, and $5$ is never circled
\cite[\S7.4]{AK} (Fact~\ref{fact:fig11}); in Diagram~6 both wings are circled, and a circled
$z$-wing has $z$-circles exactly at $1,2$ (Rules~1a, 1d), and likewise for $w$.

Permutations of the bodies (with the masses) and $z\leftrightarrow w$ are symmetries of
\eqref{eq:system}.  The thirty coloured labels of Diagram~6 (a common vertex and an ordered pair
of wings) are images of this numbering, and $z\leftrightarrow w$ composed with $(13)(24)$
preserves it.  The chart below uses rotation invariants only, so the gauge $z_{12}=w_{12}$ of
$A_m$ plays no role, and it suffices to treat this numbering at arbitrary positive masses.

\begin{definition}\label{d6-def-chart}
The \emph{butterfly chart} has the $21$ coordinates $t$, $Z_1,\dots,Z_5$, $W_1,\dots,W_5$,
$D_{kl}$, related to \eqref{eq:system} by $(z_k,w_k)=(t^{-1}Z_k,tW_k)$ for $k=1,2$,
$(z_k,w_k)=(tZ_k,t^{-1}W_k)$ for $k=3,4$, $(z_5,w_5)=(tZ_5,tW_5)$, $\delta_{kl}=tD_{kl}$ on the
\emph{crossed pairs} $13,14,23,24$ and $\delta_{kl}=D_{kl}$ on the \emph{wing pairs}
$12,15,25,34,35,45$.  Its \emph{rows} are the $z$- and $w$-equations
$z_k-\sum_{l\ne k}m_l\delta_{kl}^3(z_k-z_l)=0$ (likewise for $w$) of body $k$, multiplied by
$(t,t^{-1})$ for $k=1,2$, by $(t^{-1},t)$ for $k=3,4$ and by $(t,t)$ for $k=5$ and denoted
$R^z_k,R^w_k$; the ten distance equations $\delta_{kl}^2z_{kl}w_{kl}-1=0$; and the gauges
$Z_2-Z_1-1=0$, $W_4-W_3-1=0$.  We write $F=0$ for these $22$ rows and
$\mathcal B=F^{-1}(0)\cap\{t=0\}$ for the \emph{boundary fibre}.  ($Z_k,W_k$ are not the terms
$Z_{kl}=z_{kl}\delta_{kl}^3$ of \cite[\S3]{AK}.)
\end{definition}

\begin{lemma}[computer-assisted]\label{d6-lem-chart}
Let $m\in\Rp^5$.
\begin{romanlist}
  \item The rows are polynomials in the $21$ coordinates, of degree at most $5$ in $t$.
  \item Let $(z,w,\delta)\in A_m$ (so $z_{12}w_{34}\ne0$), choose $t$ with
  $t^2=(z_{12}w_{34})^{-1}$ and put $\nu=(tz_{12})^{-1}$.  The chart point of the rotated
  solution $(\nu z,\nu^{-1}w,\delta)$ has $Z_k=z_k/z_{12}$, $W_k=w_kz_{12}$ $(k=1,2)$,
  $Z_k=z_kw_{34}$, $W_k=w_k/w_{34}$ $(k=3,4)$, $Z_5=z_5w_{34}$, $W_5=w_5z_{12}$,
  $D_{kl}=\delta_{kl}/t$ on crossed pairs and $D_{kl}=\delta_{kl}$ on wing pairs, and satisfies
  $F=0$: each coordinate row is $t^{\pm1}\nu^{\pm1}$ times an equation of \eqref{eq:system}, the
  distance rows are those of \eqref{eq:system}, and both gauges hold.
  \item At $t=0$ the crossed distance rows read $D_{kl}^2Z_kW_l=-1$, and the wing distance rows
  read $D_{12}^2(Z_2-Z_1)(W_2-W_1)=1$, $-D_{15}^2Z_1(W_5-W_1)=1$, $-D_{25}^2Z_2(W_5-W_2)=1$,
  $D_{34}^2(Z_4-Z_3)(W_4-W_3)=1$, $-D_{35}^2(Z_5-Z_3)W_3=1$, $-D_{45}^2(Z_5-Z_4)W_4=1$.  So
  every $D_{kl}$ is nonzero on $\mathcal B$.
  \item Identically,
  $m_1R^z_1+m_2R^z_2+m_5R^z_5+t^2(m_3R^z_3+m_4R^z_4)=m_1Z_1+m_2Z_2+t^2(m_3Z_3+m_4Z_4+m_5Z_5)$,
  and likewise after $z\leftrightarrow w$, $1\leftrightarrow3$, $2\leftrightarrow4$.  Hence on
  $\mathcal B$: $Z_1=-m_2/(m_1+m_2)$, $Z_2=m_1/(m_1+m_2)$, $W_3=-m_4/(m_3+m_4)$,
  $W_4=m_3/(m_3+m_4)$.
  \item At $t=0$ the rows $R^z_1,R^z_2,R^z_5,R^w_1,R^w_2$, the distance rows of $12,15,25$ and
  the first gauge involve only the \emph{$z$-wing coordinates}
  $Z_1,Z_2,W_1,W_2,W_5,D_{12},D_{15},D_{25}$; the remaining rows except the crossed distance rows
  involve only the \emph{$w$-wing coordinates} $Z_3,Z_4,Z_5,W_3,W_4,D_{34},D_{35},D_{45}$; and
  $D_{13},D_{14},D_{23},D_{24}$ occur only in the crossed distance rows.
\end{romanlist}
\end{lemma}
We check (i)--(iv) by exact expansion over $\mathbb Q$, with the masses and the coordinates as
indeterminates; for (ii) we substitute the chart point of the rotated solution into the $22$
rows (Appendix~\ref{app:verification}).  Item (v) holds because every term coupling the two
wings carries a positive power of $t$.

\subsection{Coverage}

\begin{proposition}[Coverage]\label{prop:d6-cover}
Let $m\in\Rp^5$ and let $\gamma(s)$, $s\to0$, be a Diagram-6 germ, in the numbering above, on a
curve in $A_m$.  Along $\gamma$ define $t$ and $\nu$ as in Lemma~\ref{d6-lem-chart}(ii), and
let $\tilde\gamma(s)$ be the chart point of the rotated solution.  Then $t\to0$, every chart
coordinate is bounded along $\tilde\gamma$, and $\tilde\gamma(s)$ converges to a point
$p\in\mathcal B$ at which $Z_1,Z_2,W_3,W_4$ and all ten $D_{kl}$ are nonzero.
\end{proposition}

\begin{proof}
Normalize as in \cite[p.~547]{AK} so that the largest components of $\mathcal Z$ and $\mathcal W$
both have modulus $\epsilon^{-2}$; by Lemma~\ref{lem:germ} this is realized along $\gamma$, with
$\epsilon$ comparable to a positive power of $s$.  Let $\ord$ be the order in $\epsilon$
($\ord f>\alpha$ means $f\epsilon^{-\alpha}\to0$).  The chart coordinates are invariant under
$z\mapsto az$, $w\mapsto a^{-1}w$, so we compute in this frame.

By the circles, $\ord z_1=\ord z_2=\ord w_3=\ord w_4=-2$, and the uncircled
$z_3,z_4,z_5,w_1,w_2,w_5$ have order $>-2$; so $\ord z_{15}=\ord z_{25}=-2$.  A $z$-stroke $kl$
means $\ord(z_{kl}\delta_{kl}^3)=\ord(z_{kl}^{-1/2}w_{kl}^{-3/2})=-2$, that is,
$\ord z_{kl}+3\ord w_{kl}=4$.  So $\ord w_{15}=\ord w_{25}=2$, hence $\ord w_{12}\ge2$, and the
stroke $12$ gives $\ord z_{12}\le-2$; with $\ord z_{12}\ge-2$ (Estimate~1,
Fact~\ref{fact:rules}) we get $\ord z_{12}=-2$, $\ord w_{12}=2$.  Thus $r_{kl}^2=z_{kl}w_{kl}$ and
$\delta_{kl}$ have order $0$ on the three $z$-wing pairs (maximal edges, $r\approx1$ as in
Estimate~2), and symmetrically on the $w$-wing.  On a crossed pair $z_{kl}$ and $w_{kl}$ are
differences of a circled and an uncircled coordinate, so both have order $-2$ and
$\ord\delta_{kl}=2$.  In the $w$-equation $w_1=\sum_{l\ne1}m_l\delta_{1l}^3(w_1-w_l)$ the terms
have orders $2$ ($l=2,5$) and $4$ ($l=3,4$); so $\ord w_1\ge2$, likewise
$\ord w_2,\ord z_3,\ord z_4\ge2$, and $w_5=w_1+w_{15}$, $z_5=z_3+z_{35}$ have order $\ge2$.
(These are the estimates $A_w,B_z,m_1w_1+m_2w_2,w_5,B_w,A_z\preccurlyeq\epsilon^2$ of
\cite[\S8.6]{AK}.)

Finally $\ord t=-\frac12(\ord z_{12}+\ord w_{34})=2$ and $\ord\nu=-\ord t-\ord z_{12}=0$.  By
Lemma~\ref{d6-lem-chart}(ii), $Z_1,Z_2,W_3,W_4$ and all ten $D_{kl}$ have order exactly $0$
(crossed pairs: $2+\frac12(-2-2)=0$), and $W_1,W_2,W_5,Z_3,Z_4,Z_5$ have order $\ge0$.  Bounded
Puiseux series in $s$ converge, so $\tilde\gamma(s)\to p\in\mathcal B$, and the coordinates of
order exactly $0$ have nonzero limits.
\end{proof}
In particular $D_{12}=\delta_{12}$ and $D_{34}=\delta_{34}$ have finite nonzero limits along every
Diagram-6 germ, at every mass.

\subsection{The first condition}

Put $\Sigma=\sum_{k\in\{1,2\},\,l\in\{3,4\}}m_km_lD_{kl}^3Z_kW_l$.

\begin{lemma}[computer-assisted]\label{d6-lem-torque}
There is a polynomial $Q$ in the chart coordinates, with coefficients polynomial in $m$, such
that identically
\begin{equation}\label{d6-eq-torque}
  \sum_{k=1,2}m_k\bigl(W_kR^z_k-Z_kR^w_k\bigr)+m_5W_5R^z_5=-t\,\Sigma+t^2Q.
\end{equation}
\end{lemma}
Expanding the left-hand side of \eqref{d6-eq-torque} in exact arithmetic over $\mathbb Q$, with
the masses and the chart coordinates as indeterminates, we find that its coefficients of $t^0$
and $t^1$ are $0$ and $-\Sigma$ (Appendix~\ref{app:verification}).
With $\rho^z_k,\rho^w_k$ the left-hand sides of \eqref{eq:system}, the pair torque
$\sum_{k=1,2}m_k(w_k\rho^z_k-z_k\rho^w_k)$ equals
$\sum_{k=1,2}\sum_{l=3,4,5}m_km_l\delta_{kl}^3(w_kz_l-z_kw_l)$ (the terms with $l\in\{1,2\}$
cancel), which is $2i$ times (31) of \cite[\S8.6]{AK}.  In the chart it is the first sum in
\eqref{d6-eq-torque}; its terms with $l=3,4$ are $-t\Sigma+t^5\sum m_km_lD_{kl}^3W_kZ_l$, those with
$l=5$ are $\sum_{k=1,2}m_km_5D_{k5}^3(t^2W_kZ_5-Z_kW_5)$, and
$R^z_5-(m_1D_{15}^3Z_1+m_2D_{25}^3Z_2)$ is divisible by $t^2$.

\begin{proposition}[First condition]\label{prop:d6-first}
Let $m\in\Rp^5$, let $\gamma$ be a Diagram-6 germ in the numbering above, and let $p$ be its
limit (Proposition~\ref{prop:d6-cover}).  Then:
\begin{romanlist}
  \item $\Sigma(p)=0$, and $\Sigma(p)=-(m_1m_3D_{13}+m_1m_4D_{14}+m_2m_3D_{23}+m_2m_4D_{24})(p)$;
  \item $m_1=m_2$ and $m_3=m_4$; this is relation (32) of \cite[\S8.6]{AK};
  \item $(D_{13},D_{14},D_{23},D_{24})(p)=(2i\eta_{13},\,2\eta_{14},\,2\eta_{23},\,2i\eta_{24})$
  with signs $\eta_{kl}\in\{\pm1\}$ satisfying $\eta_{14}=-\eta_{23}$ and $\eta_{13}=-\eta_{24}$.
\end{romanlist}
In another numbering, wings $\{a,b,c\}$ and $\{d,e,c\}$ force $m_a=m_b$ and $m_d=m_e$.
\end{proposition}

\begin{proof}
All rows vanish along $\tilde\gamma$, so \eqref{d6-eq-torque} gives $t\Sigma=t^2Q$, and $t\ne0$
gives $\Sigma=tQ$.  By Proposition~\ref{prop:d6-cover}, $Q$ is bounded and $t\to0$, so
$\Sigma(p)=0$.  By Lemma~\ref{d6-lem-chart}(iii), $D_{kl}^2Z_kW_l=-1$ at $p$ for the crossed
pairs, so $\Sigma(p)=\sum m_km_lD_{kl}(D_{kl}^2Z_kW_l)=-\sum m_km_lD_{kl}$; this is (i).

By Lemma~\ref{d6-lem-chart}(iv), $Z_1W_3=m_2m_4/s_0^2$, $Z_1W_4=-m_2m_3/s_0^2$,
$Z_2W_3=-m_1m_4/s_0^2$ and $Z_2W_4=m_1m_3/s_0^2$ at $p$, where $s_0>0$,
$s_0^2=(m_1+m_2)(m_3+m_4)$.  Write $m_k=\mu_k^2$ with $\mu_k>0$.  Then $D_{kl}^2=-1/(Z_kW_l)$ gives
\[
  D_{13}=\frac{i\eta_{13}s_0}{\mu_2\mu_4},\quad D_{14}=\frac{\eta_{14}s_0}{\mu_2\mu_3},\quad
  D_{23}=\frac{\eta_{23}s_0}{\mu_1\mu_4},\quad D_{24}=\frac{i\eta_{24}s_0}{\mu_1\mu_3}
\]
with signs $\eta_{kl}\in\{\pm1\}$, and
\[
  \Sigma(p)=-\frac{s_0}{\mu_1\mu_2\mu_3\mu_4}\bigl(P_{\mathrm{re}}+iP_{\mathrm{im}}\bigr),\qquad
  P_{\mathrm{re}}=\eta_{14}\mu_1^3\mu_4^3+\eta_{23}\mu_2^3\mu_3^3,\quad
  P_{\mathrm{im}}=\eta_{13}\mu_1^3\mu_3^3+\eta_{24}\mu_2^3\mu_4^3.
\]
Up to the factor $\pm i$, $P_{\mathrm{re}}+iP_{\mathrm{im}}$ is the left-hand side of (12) of
\cite[\S5.1]{AK} written in the $\mu_k$, and each of its eight sign choices arises from exactly two
of the sixteen sign classes $\eta$.
Since the $\mu_k$ and $s_0$ are real, $\Sigma(p)=0$ forces $P_{\mathrm{re}}=P_{\mathrm{im}}=0$.  If
$\eta_{14}=\eta_{23}$, then $P_{\mathrm{re}}\ne0$, and if $\eta_{13}=\eta_{24}$, then
$P_{\mathrm{im}}\ne0$; this excludes twelve classes.  In the other four,
$P_{\mathrm{re}}=\eta_{14}(x-y)(x^2+xy+y^2)$ with $x=\mu_1\mu_4$, $y=\mu_2\mu_3$ and
$x^2+xy+y^2>0$, so $\mu_1\mu_4=\mu_2\mu_3$; likewise $P_{\mathrm{im}}=0$ gives
$\mu_1\mu_3=\mu_2\mu_4$.  Then
$\mu_3(\mu_1\mu_4-\mu_2\mu_3)-\mu_4(\mu_1\mu_3-\mu_2\mu_4)=\mu_2(\mu_4-\mu_3)(\mu_4+\mu_3)$
vanishes, so $\mu_3=\mu_4$ and then $\mu_1=\mu_2$: this is (ii).  With $m_1=m_2$ and $m_3=m_4$
the displayed values of $D_{kl}$ become those of (iii).
\end{proof}
At $m_1=m_2$, $m_3=m_4$ the expression $\Sigma(p)$ vanishes for exactly the four classes of (iii).
Albouy and Kaloshin \cite[\S8.6]{AK} obtain (32) ``as in 5.1'' from the estimate
$S\preccurlyeq\epsilon^4$ in (31); the proof above makes this exact in the chart.

\subsection{The finite wing atlas}

From now on $m=(a,a,b,b,c)$ with $a,b,c>0$.

\begin{lemma}[computer-assisted]\label{d6-lem-atlas}
The points of $\mathcal B$ are exactly the points with $Z_1=W_3=-\frac12$, $Z_2=W_4=\frac12$,
\begin{gather*}
  (D_{12},D_{15},D_{25})=(d,\lambda d,\omega\lambda d),\\
  W_1=-d(a+2c\lambda),\quad
  W_2=d(a+2c\omega\lambda),\quad W_5=ad\bigl(4\lambda^{-2}-1\bigr),\\
  (D_{34},D_{35},D_{45})=(e,\kappa e,\omega'\kappa e),\\
  Z_3=-e(b+2c\kappa),\quad
  Z_4=e(b+2c\omega'\kappa),\quad Z_5=be\bigl(4\kappa^{-2}-1\bigr),\\
  (D_{13},D_{14},D_{23},D_{24})=(2i\eta_{13},\,2\eta_{14},\,2\eta_{23},\,2i\eta_{24}),\quad
  \eta_{kl}\in\{\pm1\},
\end{gather*}
where
\begin{equation}\label{d6-eq-wing}
  \omega^3=1,\quad\lambda^2=2(1+\omega),\quad(2a+c\lambda^3)d^3=1,\qquad
  \omega'^3=1,\quad\kappa^2=2(1+\omega'),\quad(2b+c\kappa^3)e^3=1.
\end{equation}
For each $\omega$ both roots $\lambda$ are nonzero.  The scale denominator $2a+c\lambda^3$
vanishes only for $(\omega,\lambda)=(1,-2)$ and $a=4c$: for $\omega=1$, $\lambda=\pm2$; for
$\omega\ne1$, $\lambda^6=-8$ and the two denominators have product $4a^2+8c^2$.  Hence a wing
has $18$ classes $(\omega,\lambda,d)$ if $a\ne4c$ and $15$ if $a=4c$, and $\mathcal B$ has
$18\cdot18\cdot16=5184$ points if $a\ne4c\ne b$, $4320$ if exactly one of $a,b$ equals $4c$, and
$3600$ if $a=b=4c$.
\end{lemma}
The proof is an exact computation over the field $\mathbb Q(i)(a,b,c)$, with the rows reduced
modulo \eqref{d6-eq-wing} by a Gr\"obner basis, and follows the argument below
(Appendix~\ref{app:verification}).  By Lemma~\ref{d6-lem-chart}(iv),(v), $Z_1=-\frac12$,
$Z_2=\frac12$, and $\mathcal B$ is the product of the two wing fibres and the crossed roots.
On the $z$-wing put $\lambda=D_{15}/D_{12}$ and $\omega=D_{25}/D_{15}$.  At $t=0$,
$2(R^z_1+R^z_2)/c=D_{15}^3-D_{25}^3$ gives $\omega^3=1$; the wing distance rows give
$D_{12}^{-2}=W_2-W_1=2D_{15}^{-2}+2D_{25}^{-2}$,
that is, $\lambda^2=2(1+\omega)$; and $2R^z_1=(2a+c\lambda^3)d^3-1$.  The row $R^w_1$ and the
distance rows of $12,15$ are linear in $(W_1,W_2,W_5)$ with determinant $d^4\lambda^2/2\ne0$;
conversely the displayed point satisfies all eight $z$-wing rows modulo \eqref{d6-eq-wing}.  The
$w$-wing follows by the wing exchange below.

Germs at infinity are complex (Remark~\ref{rem:levels}), so all eighteen classes are needed, not
only the positive real class $(\omega,\lambda)=(1,2)$, $D_{15}=D_{25}=2D_{12}$.

\begin{definition}\label{d6-def-labels}
A \emph{label} is a point of $\mathcal B$, that is, a choice of $(\omega,\lambda,d)$,
$(\omega',\kappa,e)$ and $\eta$ in Lemma~\ref{d6-lem-atlas}.  The wing \emph{shapes}
$(\omega,\lambda)$ and $(\omega',\kappa)$ are \emph{Gaussian} if $\omega=1$ (so $\lambda=\pm2$) and
\emph{primitive} otherwise (so $\lambda^4-2\lambda^2+4=0$); each wing has two Gaussian and four
primitive shapes.  A label is \emph{corank-one} if $\eta_{14}=-\eta_{23}$ and
$\eta_{13}=-\eta_{24}$, and \emph{regular} otherwise.  For $a\ne4c\ne b$ there are
$12\cdot324=3888$ regular and $4\cdot324=1296$ corank-one labels.
\end{definition}

\emph{Symmetries.}  At the masses $(a,a,b,b,c)$, $F^{-1}(0)$ is preserved by the exchange of $1$
and $2$ combined with $t\mapsto it$, $(Z_1,Z_2,W_1,W_2,W_5)\mapsto(-Z_2,-Z_1,-W_2,-W_1,-W_5)$ and
$D_{kl}\mapsto-iD_{kl}$ on crossed pairs, by $(t,D_{13},D_{14},D_{23},D_{24})\mapsto
-(t,D_{13},D_{14},D_{23},D_{24})$, and by complex conjugation; the \emph{wing exchange}
($z\leftrightarrow w$, $1\leftrightarrow3$, $2\leftrightarrow4$) maps the chart at $(a,a,b,b,c)$ to
that at $(b,b,a,a,c)$.  On the $576$ \emph{states} ($z$-wing shape, $w$-wing shape, $\eta$) these
maps have $29$ orbits, $21$ regular and $8$ corank-one (Appendix~\ref{app:verification}).  The
computations below express the masses through the scales, $a=\frac12(d^{-3}-c\lambda^3)$,
$b=\frac12(e^{-3}-c\kappa^3)$, so one symbolic computation covers the nine choices of $(d,e)$ of a
state at every positive $(a,b,c)$.

\begin{remark}[Regular labels]\label{d6-rem-regular}
Regular labels are already excluded by Proposition~\ref{prop:d6-first}(iii).  Independently, $DF$
has rank $21$ at every regular label, which is therefore an isolated point of $F^{-1}(0)$: at one
state of each regular orbit the minor omitting $R^z_1$ is a unit in $\{\pm1,\pm i,\pm1\pm i\}$
times $-\frac{81}{128}a^2b\,d^2e^4H(\lambda)H(\kappa)(2a+4c\lambda)/(2a+c\lambda^3)$, where
$H(x)=x^3(x^2-2)^2((x^2-2)^2+4)$, and this is nonzero on finite labels.  The minors are computed
in exact arithmetic over the field $\mathbb Q(i,\sqrt2,\sqrt3)(c,d,e)$, and their nonvanishing is
checked with resultants and norms (Appendix~\ref{app:verification}).
\end{remark}

\begin{lemma}[computer-assisted]\label{d6-lem-corank}
Let $p$ be a corank-one label.
\begin{romanlist}
  \item The $20$ rows of $F$ other than $R^w_2$ and $R^w_4$, together with the coordinate $t$,
  have invertible derivative at $p$.  Hence near $p$ the common zeros of these $20$ rows form a
  smooth curve $C_p$, the graph of an analytic map of $t$.
  \item Along $C_p$, $R^w_2=c_kt^k+O(t^{k+1})$ with $c_k\ne0$, where $k$ and $c_k$ are as in
  Table~\ref{d6-tab-corank}.
  \item Consequently $p$ is an isolated point of $F^{-1}(0)$.
\end{romanlist}
\end{lemma}

\begin{table}[ht]
\centering
\small
\begin{tabular}{p{3.3cm}p{0.9cm}p{0.3cm}p{6.6cm}}
\hline
wing shapes ($z$, $w$) & labels & $k$ & $c_k$ is a nonzero multiple of \\
\hline
Gaussian, Gaussian & $144$ & $4$ &
$Q_{2,2}=-\frac{2a+2b+c}{60(a+4c)(b+4c)}$, $Q_{2,-2}=\frac{2a+2b+c}{4(a+4c)(b-4c)}$,
$Q_{-2,-2}=\frac{2a+2b+c}{68(a-4c)(b-4c)}$ for $(\lambda,\kappa)=(2,2),(2,-2),(-2,-2)$;
$(-2,2)$ by wing exchange \\
one Gaussian, one primitive & $576$ & $3$ &
$P=A\kappa^3-4A\kappa+B\kappa^3-B\kappa^2-4B\kappa+B$ for $\lambda=\pm2$, $\kappa$ primitive;
product over the primitive $\kappa$: $9(8A^2+16AB+9B^2)^2$ \\
primitive, primitive: orbit of $(\lambda_0,\lambda_0)$ & $288$ & $2$ &
$1-2\kappa^3$; product over the primitive $\kappa$: $1089$ \\
primitive, primitive: orbits of $(\lambda_0,-\lambda_0)$ & $288$ & $2$ &
$2B/(2B+\kappa^3)$; product over the primitive $\kappa$: $B^4/(B^2+2)^2$ \\
\hline
\end{tabular}
\caption{Order $k$ and leading coefficient of $R^w_2$ along $C_p$ at the corank-one labels
(Lemma~\ref{d6-lem-corank}), at one representative state of each orbit.  Counts are for
$a\ne4c\ne b$; $A=a/c$, $B=b/c$; $\lambda_0=\frac12(\sqrt6+i\sqrt2)$ is the primitive ratio for
$\omega=\frac12(-1+i\sqrt3)$.  In the first row $R^w_4$ has the same order.}\label{d6-tab-corank}
\end{table}

The proof of (i) and (ii) is an exact computation at one representative state of each of the
eight corank-one orbits, over the field $\mathbb Q(i)(c,d,e)$ if both wings are Gaussian and over
its extension by a root $\kappa$ of $\kappa^4-2\kappa^2+4$ otherwise
(Appendix~\ref{app:verification}).  In (i) the derivative is reduced to the Schur complement of
the diagonal block of the distance rows.  For two Gaussian wings its determinant (rows multiplied
by $d^3e^3$) is $82944\,d^{61}e^{59}(cd^3-\frac{\lambda}{16})(ce^3-\frac{\kappa}{16})$, with
$cd^3-\frac{\lambda}{16}=-\frac{\lambda a}{8(2a+c\lambda^3)}\ne0$.  Otherwise the complement, after
removal of the unit row of $t$, is block diagonal with one $5\times5$ block per wing
(Lemma~\ref{d6-lem-chart}(v)); we check that an off-diagonal block vanishes and evaluate both
diagonal blocks.  For one Gaussian and one primitive wing the two determinants are nonzero
multiples of $8cd^3-\frac{\lambda}{2}=-\lambda A/(2A+\lambda^3)$ and of a polynomial in $ce^3$ and
$\kappa$ whose product over the primitive $\kappa$ is
$(\frac9{16})^4\cdot64B^4(B^4-8B^2+64)/(B^2+2)^4$, where $B^4-8B^2+64=(B^2-4)^2+48$.  This
product and those in Table~\ref{d6-tab-corank} are obtained from resultants with
$\kappa^4-2\kappa^2+4$.  For (ii) the jets of $C_p$ are computed by solving the $20$ rows of (i)
order by order in $t$, and are checked by substitution into all $22$ rows.  Item (iii) follows
from (i) and (ii): near $p$, $F^{-1}(0)\subset C_p$, and on $C_p$ the row $R^w_2$ vanishes only at
$t=0$, that is, at $p$.  With Remark~\ref{d6-rem-regular}, every label is thus an isolated point
of $F^{-1}(0)$; for $a\ne4c\ne b$ these are $3888+1296=5184$ labels.

\subsection{Exclusion}

\begin{theorem}\label{thm:d6}
Let $m\in\Rp^5$.  No curve in $A_m$ has a Diagram-6 germ, for any numbering of the bodies and
either colouring.  In particular the case $r=6$ of Theorem~\ref{thm:exclusion}(a) holds, at the
complex level.
\end{theorem}

\begin{proof}
By the symmetries described at the beginning of this section we may use the numbering above.  Let
$\gamma(s)$ be a Diagram-6 germ on a curve in $A_m$.  By Proposition~\ref{prop:d6-cover} the
chart germ $\tilde\gamma(s)$ has $t(s)\to0$ with $t(s)\ne0$ for $s\ne0$, bounded coordinates,
and a limit $p\in\mathcal B$.  By Proposition~\ref{prop:d6-first}, $m=(a,a,b,b,c)$ and the
crossed signs of $p$ satisfy $\eta_{14}=-\eta_{23}$, $\eta_{13}=-\eta_{24}$; by
Lemma~\ref{d6-lem-atlas}, $p$ is a label, hence a corank-one label.  By
Lemma~\ref{d6-lem-corank}, $p$ is an isolated point of $F^{-1}(0)$.  The coordinates of
$\tilde\gamma$ are Puiseux series in $s$ without negative exponents, so after a substitution
$s=\sigma^N$ they are convergent power series in $\sigma$, and $\sigma\mapsto\tilde\gamma(\sigma^N)$
is a holomorphic map of a disc into $F^{-1}(0)$ with value $p$ at $\sigma=0$.  Since $p$ is
isolated, this map is constant near $\sigma=0$, contradicting $t\not\equiv0$.
\end{proof}

The argument covers every ramification, every crossed sign class and the faces $a=4c$, $b=4c$,
since Lemma~\ref{d6-lem-atlas} describes the whole fibre $\mathcal B$.

\begin{remark}\label{d6-rem-projective}
On the face $a=4c$ the wing class $(\omega,\lambda)=(1,-2)$ has no finite scale; it lies at
$D_{12}=\infty$.  The proof of Theorem~\ref{thm:d6} never meets this class: by
Proposition~\ref{prop:d6-cover}, $D_{12}$ and $D_{34}$ have finite nonzero limits along every
Diagram-6 germ.  In a one-wing projective chart with $h=1/D_{12}$, and on $a=b=4c$ in a
double-projective ratio chart, one can show that no analytic branch through a point of
$\{t=h=0\}$ leaves the divisor $h=0$, though not that these points are isolated
(Appendix~\ref{app:verification}).  We do not use this.
\end{remark}



\section{Diagram 7}\label{sec:d7}

Diagram~7 is the butterfly of \cite[\S7.4]{AK} without circles.  Albouy and Kaloshin derive for
it the mass relations (36), (39) and (40) in \cite[\S\S8.7.1--8.7.3]{AK}.  We use none of these
relations.

\begin{theorem}\label{thm:d7}
Let $m\in\Rp^5$.  No Diagram~7 germ with an unbounded four-product lies on a curve through a
real point: if an irreducible curve $\Gamma\subset A_m$ carries such a germ, then $\Gamma$
contains no real point.
\end{theorem}

This is part~(b) of Theorem~\ref{thm:exclusion} for $r=7$.  Germs with $a+b\ge0$ in
Proposition~\ref{d7-cover} have no unbounded four-product (Lemma~\ref{d7-bounded}); the others
lie in the open cone of a nested chart, where an identity valid modulo \eqref{eq:system} gives
$I\to0$ (Lemma~\ref{d7-vanish}), so $I=U\equiv0$ on $\Gamma$, which is impossible at a real point.

\subsection{The diagram and the nested chart}
In the numbering of \cite[\S8.7]{AK}, the \emph{upper wing} $\{1,2,5\}$ is an isolated triangle
of $z$-edges without $z$-circles, the \emph{lower wing} $\{3,4,5\}$ is an isolated triangle of
$w$-edges without $w$-circles, and the four \emph{crossed} pairs $kj$ ($k\in\{1,2\}$,
$j\in\{3,4\}$) are not edges.
Since Diagram~7 has no circles, Rule~1e shows that neither wing has a maximal stroke (see the
proof of Proposition~\ref{d7-cover}(i)).  By Definition~\ref{def:germ} and the convention after
Theorem~\ref{thm:exclusion}, we may use this numbering and colouring.

\emph{Orders.}  As in Lemma~\ref{lem:germ} we renormalize a germ so that $Z$ and $W$ have a
common order, and reparametrize it by a power $t$ of its Puiseux parameter so that this order is
$-2$; thus $t$ plays the role of $\epsilon$ in \cite[\S3]{AK}, and $\ord$ is the order in $t$.
Renormalizations $(z,w)\mapsto(\lambda z,\lambda^{-1}w)$ preserve the diagram, the ties, $I$, $U$
and every $r_{kl}^2=z_{kl}w_{kl}$.

\emph{The nested chart} with exponents $P,Q,R$ has the coordinates $A_1,A_2,C,u_3,u_4$,
$B_3,B_4,D,v_1,v_2$ and $D_{kl}$ ($k<l$):
\begin{gather*}
  z_1=A_1,\quad z_2=A_2,\qquad w_3=t^{-Q}B_3,\quad w_4=t^{-Q}B_4,\\
  z_j=C+t^Pm_5u_j\ \ (j=3,4),\qquad z_5=C-t^P(m_3u_3+m_4u_4),\\
  w_k=t^{-Q}D+t^Rm_5v_k\ \ (k=1,2),\qquad w_5=t^{-Q}D-t^R(m_1v_1+m_2v_2),
\end{gather*}
and $\delta_{kl}=t^{e_{kl}}D_{kl}$, where $e_{kl}=-R/2$ on the upper wing, $Q/2$ on the crossed
pairs and $(Q-P)/2$ on the lower wing ($D$ without indices is not an inverse distance).
So $C$ and $t^{-Q}D$ are the centres of mass of the $z$-cluster $\{3,4,5\}$ and of the
$w$-cluster $\{1,2,5\}$, $P$ and $R$ are their internal orders, and $-Q$ is the order of the
common $w$-scale.
The exponents $e_{kl}$ balance all ten ties $\delta_{kl}^2z_{kl}w_{kl}=1$; those of the pairs
$12$ and $34$ read exactly
\[
  m_5D_{12}^2(A_2-A_1)(v_2-v_1)=1,\qquad m_5D_{34}^2(u_4-u_3)(B_4-B_3)=1.
\]
The \emph{open cone} is $P>Q>0$, $R>0$.

\emph{Gauges.}  The chart does not impose the gauge $z_{12}=w_{12}$.  By the displayed ties,
$A_2-A_1$ and $B_4-B_3$ never vanish, so they could serve as gauges for the rotations and for
the rescalings of $t$.  We fix none, since $I$, $U$, the ties and the zero set of the rows of
\eqref{eq:system} are invariant under the rotations.

\subsection{Coverage}

\begin{proposition}[Coverage]\label{d7-cover}
Let $\gamma$ be a Diagram~7 germ.  There are $a,b$ with $-2<a<1$ and $-2<b<1$ such that the
following hold.
\begin{romanlist}
  \item On the upper wing $\ord z_{kl}=a$ and $\ord w_{kl}=(4-a)/3$; on the lower wing
  $\ord w_{kl}=b$ and $\ord z_{kl}=(4-b)/3$; on the crossed pairs $\ord z_{kl}=a$ and
  $\ord w_{kl}=b$.  Moreover $\ord z_k\ge a$ and $\ord w_k\ge b$ for every $k$.
  \item After the rotation $(z,w)\mapsto(t^{-a}z,t^{a}w)$, $\gamma$ is a germ of the nested
  chart with exponents
  \[
    P=\frac{4-b}{3}-a,\qquad Q=-(a+b),\qquad R=\frac{4+2a}{3};
  \]
  all chart coordinates have nonnegative order, so they converge, and every $D_{kl}$ has order
  zero.  Moreover $R=\frac23(a+2)>0$ and $P-Q=\frac23(b+2)>0$, so $(P,Q,R)$ lies in the open
  cone if and only if $a+b<0$.
\end{romanlist}
\end{proposition}

\begin{proof}
(i) We use Fact~\ref{fact:rules}.  By Rule~2e \cite[p.~552]{AK} the three $z$-separations of a
wing have one order, and so do its three $w$-separations; let $a$ be the $z$-order of the upper
wing and $b$ the $w$-order of the lower wing.  As $\delta_{kl}^2z_{kl}w_{kl}=1$,
$\ord Z_{kl}=-(\ord z_{kl}+3\ord w_{kl})/2$ and $\ord W_{kl}=-(3\ord z_{kl}+\ord w_{kl})/2$;
since the strokes have order $-2$, this gives the other two wing orders.  The upper wing is an
isolated component of the $z$-diagram without $z$-circles, so it has no maximal $z$-stroke
(Rule~1e, p.~546), and along a $z$-stroke Estimate~2 (p.~548) gives $-2\le\ord z_{kl}\le1$, with
$-2$ exactly for a maximal $z$-edge and $1$ exactly for a $zw$-edge.  Hence $-2<a<1$, and dually
$-2<b<1$.  As $z_{13}=z_{15}+z_{53}$
with $\ord z_{53}=(4-b)/3>1>a$, and $w_{13}=w_{15}+w_{53}$ with $\ord w_{15}=(4-a)/3>1>b$, the
crossed orders are $a$ and $b$; likewise for $14$, $23$, $24$.  Finally
$\sum_km_kz_k=\sum_km_kw_k=0$ on $A_m$ (proof of Lemma~\ref{lem:distances}), so each $z_k$ is a
combination of $z$-differences, all of order at least $a$; likewise $\ord w_k\ge b$.

(ii) After the rotation the $z$-differences have order $0$ on the upper wing and the crossed
pairs and $P$ on the lower wing, the $w$-differences have order $R$ on the upper wing and $-Q$
elsewhere, and $\ord z_k\ge0$, $\ord w_k\ge-Q$.  Inverting the chart map, $A_k=z_k$,
$B_j=t^Qw_j$, $C$ and $t^{-Q}D$ are the centres of the two clusters, $m_5t^Pu_j=z_j-C$ and
$m_5t^Rv_k=w_k-t^{-Q}D$ are combinations of differences inside them, and
$D_{kl}=t^{-e_{kl}}\delta_{kl}$, where $\ord\delta_{kl}=-(\ord z_{kl}+\ord w_{kl})/2=e_{kl}$.  So
every coordinate has nonnegative order and $\ord D_{kl}=0$.
\end{proof}

Thus a germ lies in the open cone exactly when $a+b<0$, and the germs with $a+b\ge0$ have no
unbounded four-product (Lemma~\ref{d7-bounded}).  In the terminology of \cite[\S8.7]{AK}, a germ
with a wing in the general case has $a+b<0$; we shall not need this.

\subsection{The inertia identity}
Write
\[
  F^z_k=z_k-\sum_{l\ne k}m_l\delta_{kl}^3(z_k-z_l),\qquad
  F^w_k=w_k-\sum_{l\ne k}m_l\delta_{kl}^3(w_k-w_l)
\]
for the coordinate rows of \eqref{eq:system} and $T_{kl}=\delta_{kl}^2z_{kl}w_{kl}-1$ for the
ties.  Then $I-U=\sum_km_kw_kF^z_k+\sum_{k<l}m_km_l\delta_{kl}T_{kl}$, so $U=I$ on $A_m$
\cite[p.~541]{AK}.

\begin{lemma}[Inertia identity]\label{d7-identity}
Let
\[
  S_5=\sum_{j=3,4}\sum_{l=1,2}m_jm_l\delta_{jl}^3
  \bigl((z_j-z_5)(w_l-w_5)-(w_j-w_5)(z_l-z_5)\bigr).
\]
Then, as polynomials in the masses and in $(z,w,\delta)$,
\begin{multline*}
  I=\sum_{k=1,2}m_kz_k(w_k-w_5)+\sum_{j=3,4}m_jw_j(z_j-z_5)+S_5\\
  -\sum_{j=3,4}m_j\bigl((z_j-z_5)F^w_j-(w_j-w_5)F^z_j\bigr)+w_5\sum_{k=1}^5m_kF^z_k.
\end{multline*}
In particular $I=\sum_{k=1,2}m_kz_k(w_k-w_5)+\sum_{j=3,4}m_jw_j(z_j-z_5)+S_5$ on $A_m$.
\end{lemma}

\begin{proof}
Put $\hat z_k=z_k-z_5$ and $\hat w_k=w_k-w_5$.  By antisymmetry of the force terms,
$\sum_km_kF^z_k=\sum_km_kz_k$.  Since $z_j\hat w_j-w_j\hat z_j=z_5w_j-w_5z_j$,
\begin{align*}
  I&=\sum_km_kz_k\hat w_k+w_5\sum_km_kz_k\\
   &=\sum_{k=1,2}m_kz_k\hat w_k+\sum_{j=3,4}m_j\bigl(w_j\hat z_j+z_5w_j-w_5z_j\bigr)
    +w_5\sum_km_kF^z_k.
\end{align*}
For $j\in\{3,4\}$, using $z_j-z_l=\hat z_j-\hat z_l$ and $w_j-w_l=\hat w_j-\hat w_l$,
\[
  \hat z_jF^w_j-\hat w_jF^z_j=(w_5z_j-z_5w_j)
  +\sum_{l\ne j}m_l\delta_{jl}^3(\hat z_j\hat w_l-\hat w_j\hat z_l).
\]
In $\sum_{j=3,4}m_j(\hat z_jF^w_j-\hat w_jF^z_j)$ the terms with $l=5$ vanish, those with
$l\in\{3,4\}$ cancel in pairs, and those with $l\in\{1,2\}$ add up to $S_5$.  Hence
$\sum_{j=3,4}m_j(z_5w_j-w_5z_j)=S_5-\sum_{j=3,4}m_j(\hat z_jF^w_j-\hat w_jF^z_j)$, and
substituting this into the display gives the identity.
\end{proof}

The term $S_5$ is the torque about body~$5$ of the crossed forces on the lower wing; compare
\cite[(31)]{AK}.

\begin{lemma}\label{d7-vanish}
Let $\gamma$ be a Diagram~7 germ on $A_m$ with $a+b<0$ (Proposition~\ref{d7-cover}).  Then
$I(\gamma(t))\to0$ as $t\to0$.
More precisely, in the nested chart of Proposition~\ref{d7-cover}(ii) the three terms
$\sum_{k=1,2}m_kz_k(w_k-w_5)$, $\sum_{j=3,4}m_jw_j(z_j-z_5)$ and $S_5$ are sums of powers $t^e$
times polynomials in the chart coordinates, with exponents
\[
  e=R;\qquad e=P-Q;\qquad e\in\bigl\{\tfrac Q2,\ P+\tfrac Q2,\ \tfrac{3Q}2+R,\ P+\tfrac{3Q}2+R\bigr\},
\]
respectively.
\end{lemma}

\begin{proof}
On $A_m$ the last two sums in Lemma~\ref{d7-identity} vanish, and the three remaining terms are
invariant under the rotations, so we may evaluate them on the rotated germ in the chart.  There
\begin{gather*}
  w_k-w_5=t^R\bigl((m_k+m_5)v_k+m_{k'}v_{k'}\bigr),\qquad
  z_j-z_5=t^P\bigl((m_j+m_5)u_j+m_{j'}u_{j'}\bigr),\\
  w_j-w_5=t^{-Q}(B_j-D)+t^R(m_1v_1+m_2v_2),\\
  z_l-z_5=A_l-C+t^P(m_3u_3+m_4u_4),
\end{gather*}
where $\{k,k'\}=\{1,2\}$, $\{j,j'\}=\{3,4\}$ and $l\in\{1,2\}$, and
$\delta_{jl}^3=t^{3Q/2}D_{jl}^3$.  Substituting, together with $z_k=A_k$ and $w_j=t^{-Q}B_j$,
gives the stated exponents.  Since $a+b<0$, Proposition~\ref{d7-cover}(ii) puts $(P,Q,R)$ in the
open cone, so every exponent is positive, and all chart coordinates are bounded.  Hence the three
terms, and with them $I$, tend to zero.
\end{proof}

\subsection{Germs outside the open cone}

\begin{lemma}\label{d7-bounded}
Let $\gamma$ be a Diagram~7 germ with $a+b\ge0$ (Proposition~\ref{d7-cover}).  Then every
four-product tends to zero along $\gamma$.  If $a+b=0$, then $-1<a<1$, and $p_{12}$, $p_{34}$
and each of the other eight four-products have the orders $(16-4a)/3$, $(16+4a)/3$ and $8/3$.
If $a+b>0$, every $r_{kl}$ tends to zero, and so does $I$.
\end{lemma}

\begin{proof}
By Proposition~\ref{d7-cover}(i), $r_{kl}^2=z_{kl}w_{kl}$ has order $(4+2a)/3>0$ on the upper
wing, $(4+2b)/3>0$ on the lower wing and $a+b\ge0$ on the crossed pairs.  The four-product
$p_{ij}$ contains the three squared sides of the triangle $\{k,l,m\}$ complementary to
$\{i,j\}$, and this triangle contains a wing pair: if $5\in\{k,l,m\}$, both pairs through~$5$
are wing pairs; otherwise $\{k,l,m\}\subset\{1,2,3,4\}$ contains $\{1,2\}$ or $\{3,4\}$ (the
crossed pairs form the bipartite graph $K_{2,2}$).  Hence
$\ord p_{ij}\ge\min\{(4+2a)/3,(4+2b)/3\}>0$.  If $a+b=0$, the bounds $-2<a,b<1$ give $-1<a<1$,
and adding orders gives the stated values.  If $a+b>0$, every $r_{kl}^2$ has positive order, and
the Lagrange identity $MI=(\sum_km_kz_k)(\sum_km_kw_k)+\sum_{k<l}m_km_lz_{kl}w_{kl}$, with
$\sum_km_kz_k=\sum_km_kw_k=0$, gives $I\to0$.
\end{proof}

\begin{remark}\label{d7-balanced}
The germs with $a+b=0$ are not excluded, and need not be: they carry no unbounded four-product,
so Proposition~\ref{prop:reduction} never produces them.
\end{remark}

\subsection{End of the proof}

\begin{proof}[Proof of Theorem~\ref{thm:d7}]
Let $\gamma$ be a Diagram~7 germ with an unbounded four-product on an irreducible curve
$\Gamma\subset A_m$, numbered as above, and let $a,b$ be as in Proposition~\ref{d7-cover}.  Some
four-product has negative order along $\gamma$ (Definition~\ref{def:germ}), so $a+b<0$ by
Lemma~\ref{d7-bounded}, and $I(\gamma(t))\to0$ by Lemma~\ref{d7-vanish}.  By
Lemma~\ref{lem:U-constant}, $I=U$ is constant on $\Gamma$, with value
$\lim_{t\to0}I(\gamma(t))=0$.  At a real point $w_k=\bar z_k$, so $I=\sum_km_k|z_k|^2$,
which is positive because $m_k>0$ and the tie of the pair $12$ forces $z_1\ne z_2$.  Hence
$\Gamma$ contains no real point.
\end{proof}

\begin{remark}\label{d7-masses}
The proof of Theorem~\ref{thm:d7} uses no relation among the masses: Lemma~\ref{d7-identity} is a
polynomial identity, Proposition~\ref{d7-cover} and Lemma~\ref{d7-bounded} use only Rules~1e,
2e and Estimate~2 of \cite[\S3]{AK}, Lemma~\ref{d7-vanish} only the signs of the exponents, and
the last step only Lemma~\ref{lem:U-constant} and $m_k>0$.  So the theorem does not depend on
whether the masses satisfy (36), (39) or (40), and no mass condition is derived for Diagram~7.
\end{remark}



\section{Diagram 8}\label{sec:d8}

Diagram~8 is the butterfly of \cite[\S7.4]{AK} with exactly one circled wing
(Fact~\ref{fact:fig11}).  We use the numbering of \cite[\S8.8]{AK}: the upper wing $\{1,2,5\}$
is an isolated triangle of $z$-edges without $z$-circles, the lower wing $\{3,4,5\}$ is a
triangle of $w$-edges with $w$-circles at $3$ and $4$ (vertex $5$ is never circled
\cite[\S7.4]{AK}, and Rules~1a, 1d place the circles of a circled wing at its other two
vertices), and the crossed pairs $13$, $14$, $23$, $24$ are not edges.
The other circling becomes this one after exchanging $z$ and $w$ and renumbering
$1\leftrightarrow3$, $2\leftrightarrow4$; both operations are symmetries of \eqref{eq:system}
(the second permutes the masses), so it suffices to treat this numbering for every
$m\in\Rp^5$.
The mass condition of the diagram is relation (40) of \cite[\S8.8]{AK},
\begin{equation}\label{eq:d8-40}
  m_5^2(m_1+m_2)^4=m_1^3m_2^3,
\end{equation}
which \cite[p.~578]{AK} obtain ``by the same arguments'' as for Diagram~7.  For the other
circling it reads $m_5^2(m_3+m_4)^4=m_3^3m_4^3$, which is their (39).
We show that every Diagram~8 germ is a ray germ of a chart with a finite lower wing
(Proposition~\ref{prop:d8-cover}), derive \eqref{eq:d8-40} from the boundary of the chart
(Proposition~\ref{prop:d8-first}), and exclude all germs, complex ones included, by an
order-six pairing, an order-nine pairing and a second-order elimination on one curve of
masses (Theorem~\ref{thm:d8}).

Two further symmetries are used tacitly.  Scaling the masses by $\lambda>0$ and $(z,w,\delta)$
by $(\lambda^{1/3},\lambda^{1/3},\lambda^{-1/3})$ maps $A_m$ onto $A_{\lambda m}$ and multiplies
every coordinate and every component of $\mathcal Z$, $\mathcal W$ by a nonzero constant, so it
preserves orders and hence diagrams (Lemma~\ref{lem:germ}).
The equations of \eqref{eq:system} other than $z_{12}=w_{12}$ are invariant under $z\mapsto cz$,
$w\mapsto c^{-1}w$; this shifts all $z$-orders by $\ord c$ and all $w$-orders by $-\ord c$, so it
preserves diagrams \cite[p.~547]{AK}.

\subsection{The ratio chart}
Let $\gamma$ be a Diagram~8 germ.  We rotate it by $z_{12}^{-1}$, so that $z_{12}=1$ along
$\gamma$.  The \emph{ratio chart} has two parameters $t,h$ and nineteen coordinates, the
positions and shapes $A_1,A_2,C,B_3,B_4,D,u,v,R_{12},R_{25},P,Q$ and the inverse distances
$D_{13},D_{14},D_{23},D_{24},D_{34},D_{35},D_{45}$; we write $x$ for all of them.  It is given by
\begin{gather*}
  z_1=A_1,\quad z_2=A_2,\quad z_3=C+m_5t^6u,\quad z_4=C+m_5t^6v,\\
  z_5=C-t^6(m_3u+m_4v),\quad w_3=t^{-6}B_3,\quad w_4=t^{-6}B_4,\\
  w_1=t^{-6}D+m_5t^4h^{-2}P,\quad w_2=t^{-6}D+m_5t^4h^{-2}Q,\\
  w_5=t^{-6}D-t^4h^{-2}(m_1P+m_2Q),\quad \delta_{15}=t^{-2}h,\\
  \delta_{12}=t^{-2}hR_{12},\quad \delta_{25}=t^{-2}hR_{25},\\
  \delta_{bc}=t^3D_{bc}\ \ (b\in\{1,2\},\ c\in\{3,4\}),\qquad
  \delta_{kl}=D_{kl}\ \ (k,l\in\{3,4,5\}).
\end{gather*}
For $th\ne0$ it is invertible: $C$ and $t^{-6}D$ are the centres of mass of $z_3,z_4,z_5$ and of
$w_1,w_2,w_5$, $h=t^2\delta_{15}$, $R_{12}=\delta_{12}/\delta_{15}$,
$R_{25}=\delta_{25}/\delta_{15}$, $u=(z_3-C)/(m_5t^6)$, $P=h^2(w_1-t^{-6}D)/(m_5t^4)$, and
similarly for the other coordinates.
The rows of the chart are the gauges $E_{20}:=A_2-A_1-1$ and $E_{21}:=B_4-B_3-1$ (that is,
$z_{12}=1$ and $w_4-w_3=t^{-6}$) and $E_0,\dots,E_{19}$, obtained by substituting the chart into
the twenty equations of \eqref{eq:system} other than $z_{12}=w_{12}$ and multiplying each by a
monomial in $t$ and $h^{1/2}$: $E_{2k-2}$ and $E_{2k-1}$ come from the $z$- and the $w$-equation
of body $k$, and $E_{10},\dots,E_{19}$ from the ties of the pairs
$12,13,14,15,23,24,25,34,35,45$.
In the corner coordinates
\[
  a:=th^{-1/2},\qquad \eta:=t^2h^{1/2}\qquad(\text{so } t^3=a\eta,\ h^{3/2}=\eta\,a^{-2})
\]
we use the monomials
\[
  \mu:=a^6=t^6h^{-3},\qquad \nu:=a^2\eta^2=t^6,\qquad T:=a\eta=t^3,\qquad \eta^2=t^4h,
\]
and we write $X:=A_1-C$, $Y:=A_2-C$, $M_{12}:=m_1+m_2$, $M:=m_3+m_4+m_5$, and
$C_5:=C-\nu(m_3u+m_4v)$, which is $z_5$.

\begin{lemma}[computer-assisted]\label{lem:d8-rows}
The following hold identically in the masses.
\begin{romanlist}
\item Every $E_j$ is a polynomial in $a$, $\eta$ and $x$, and its \emph{face}
$f_j:=E_j|_{a=\eta=0}$ is
\begin{align*}
  f_0&=-\bigl(m_2(X-Y)R_{12}^3+m_5X\bigr), & f_8&=m_1X+m_2YR_{25}^3,\\
  f_2&=-\bigl(m_1(Y-X)R_{12}^3+m_5YR_{25}^3\bigr), & f_{20}&=A_2-A_1-1,\\
  f_{10}&=m_5(X-Y)(P-Q)R_{12}^2-1, & f_{13}&=X\bigl((m_1+m_5)P+m_2Q\bigr)-1,\\
  f_{16}&=YR_{25}^2\bigl(m_1P+(m_2+m_5)Q\bigr)-1, & f_4&=f_6=C,\quad f_1=f_3=D,\\
  f_{11}&=X(D-B_3)D_{13}^2-1, & f_{12}&=X(D-B_4)D_{14}^2-1,\\
  f_{14}&=Y(D-B_3)D_{23}^2-1, & f_{15}&=Y(D-B_4)D_{24}^2-1,\\
  f_5&=B_3+m_4(B_4-B_3)D_{34}^3-m_5(B_3-D)D_{35}^3, & f_{21}&=B_4-B_3-1,\\
  f_7&=B_4-m_3(B_4-B_3)D_{34}^3-m_5(B_4-D)D_{45}^3,\\
  f_9&=D+m_3(B_3-D)D_{35}^3+m_4(B_4-D)D_{45}^3,\\
  f_{17}&=m_5(B_3-B_4)(u-v)D_{34}^2-1,\\
  f_{18}&=(B_3-D)\bigl((m_3+m_5)u+m_4v\bigr)D_{35}^2-1,\\
  f_{19}&=(B_4-D)\bigl(m_3u+(m_4+m_5)v\bigr)D_{45}^2-1.
\end{align*}
\item $m_1E_0+m_2E_2+m_5E_8+a^6(m_3E_4+m_4E_6)=a^6(m_1A_1+m_2A_2+MC)$.
\item $E_1=D+\eta^2\beta_1+T^3\tau_1+\sum_{i\ge6,\,j\ge2}a^i\eta^j(\cdots)$, and $E_3$ is the
same with $\beta_3,\tau_3$, where
\begin{align*}
  \beta_1&=-m_5\bigl(m_2(P-Q)R_{12}^3+(m_1+m_5)P+m_2Q\bigr),\\
  \beta_3&=-m_5\bigl(m_1(Q-P)R_{12}^3+(m_1P+(m_2+m_5)Q)R_{25}^3\bigr),\\
  \tau_1&=m_3(B_3-D)D_{13}^3+m_4(B_4-D)D_{14}^3,\\
  \tau_3&=m_3(B_3-D)D_{23}^3+m_4(B_4-D)D_{24}^3;
\end{align*}
the monomials of $E_1-E_3$ are $\eta^2$, $T^3$, $a^6\eta^2$ and $a^9\eta^5$.
\item For $k\in\{0,2,8,10,13,16,20\}$, $E_k=f_k(\mathcal I)+\sum_{i\ge6}a^i\eta^j(\cdots)$, where
$f_k$ is read as a function of $(X,Y,R_{12},R_{25},P,Q)$ and
$\mathcal I:=(A_1-C_5,\,A_2-C_5,\,R_{12},\,R_{25},\,P,\,Q)$.
\item $E_4=C+\nu L_3+T^3F_3+T^5G_3$ exactly, where
$L_3=-m_5\bigl(m_4(u-v)D_{34}^3+((m_3+m_5)u+m_4v)D_{35}^3-u\bigr)$,
$F_3=-\bigl(m_1(C-A_1)D_{13}^3+m_2(C-A_2)D_{23}^3\bigr)$ and
$G_3=-m_5u\bigl(m_1D_{13}^3+m_2D_{23}^3\bigr)$.
\end{romanlist}
\end{lemma}
The lemma is an exact computation over $\mathbb Q$, with the masses, $a$, $\eta$ and the chart
coordinates as indeterminates (Appendix~\ref{app:verification}).  For (iv) the computation
shows that the $\nu$-coefficient of $E_k$ is $-(\partial_Cf_k)(m_3u+m_4v)$ and that every other
non-face monomial has $i\ge6$; since each of these $f_k$ is affine in $(X,Y)$, this is (iv).

\begin{definition}\label{def:d8-ray}
A germ in the ratio chart is a \emph{ray germ} if along it $t\to0$, $x$ converges, $\ord a>0$
and $\ord\eta>0$.  Its \emph{ray} is $r:=\ord\eta/\ord a\in(0,\infty)$ and its \emph{exponent}
is $p:=\ord h/\ord t$.
\end{definition}
Since $t^3=a\eta$ and $h^{3/2}=\eta a^{-2}$, $p=(2r-4)/(r+1)$, an increasing bijection of
$(0,\infty)$ onto $(-4,2)$; the rays $r=\frac12,2,3$ are $p=-2,0,\frac12$, and the faces
$\{a=0\}$, $\{\eta=0\}$ of the corner chart are $p=-4$ and $p=2$.

\subsection{Coverage}

\begin{proposition}[Coverage]\label{prop:d8-cover}
Every Diagram~8 germ is a ray germ of the ratio chart with exponent $p\in(-4,2)$ and a finite
lower wing.  Precisely: let $\gamma$ be a Diagram~8 germ, rotated so that $z_{12}=1$, and put
$t:=(w_4-w_3)^{-1/6}$ and $h:=t^2\delta_{15}$ along $\gamma$.  Let $\theta$ be the order of
$z_{12}$ before the rotation, in the normalization $Z=W=\epsilon^{-2}$ with orders in $\epsilon$.
Then $-2<\theta<1$; every coordinate in $x$ converges along $\gamma$, and $R_{12}$, $R_{25}$,
$D_{13},D_{14},D_{23},D_{24}$, $D_{34},D_{35},D_{45}$ have nonzero limits; and
\begin{gather*}
  \ord t=\tfrac{2-\theta}6,\qquad \ord h=-\tfrac{2\theta}3,\qquad
  p=-\tfrac{4\theta}{2-\theta},\\
  \ord a=\tfrac{2+\theta}6,\qquad \ord\eta=\tfrac{2-2\theta}3,\qquad
  r=\tfrac{4(1-\theta)}{2+\theta}.
\end{gather*}
\end{proposition}

\begin{proof}
By Fact~\ref{fact:rules} and Lemma~\ref{lem:germ} the rules and estimates of \cite[\S3]{AK}
hold for the orders along $\gamma$.
\emph{Lower wing.}  $\{3,4,5\}$ is an isolated component of the $w$-diagram with $w$-circles,
so it has a maximal $w$-stroke \cite[Rule~1e, p.~546]{AK}.  By Rule~2e its three
$w$-separations have one order, so all three $w$-strokes are maximal, and the right equality
case of Estimate~2 \cite[p.~548]{AK}, in the colour $w$, gives $\ord w_{kl}=-2$,
$\ord z_{kl}=2$, $\ord r_{kl}=0$ for $kl\in\{34,35,45\}$.
\emph{Upper wing.}  $\{1,2,5\}$ is an isolated component of the $z$-diagram without circles, so
it has no maximal $z$-stroke (Rule~1e); it carries no $w$-stroke, hence no $zw$-edge.  So
neither equality case of Estimate~2 holds, and by Rule~2e the three $z$-separations have one
order $\theta$ with $-2<\theta<1$.  The $z$-strokes have $\ord Z_{kl}=-2$, so
$\ord r_{kl}=(\theta+2)/3$ and $\ord w_{kl}=(4-\theta)/3$.
\emph{Crossed pairs.}  For $b\in\{1,2\}$, $c\in\{3,4\}$, $z_{bc}=z_{b5}+z_{5c}$ has order
$\theta$ (as $\theta<2$) and $w_{bc}=w_{b5}+w_{5c}$ has order $-2$ (as $(4-\theta)/3>-2$), so
$\ord r_{bc}=(\theta-2)/2$.
\emph{Positions.}  Every $w_k$ has order $\ge-2$, with equality for $k=3,4$ ($w$-circles).  Let
$C$ be the centre of mass of $z_3,z_4,z_5$.  Then $(m_1+m_2+M)C=-m_1(z_1-C)-m_2(z_2-C)$, and
$z_b-C=z_{5b}+(z_5-C)$ with $\ord z_{5b}=\theta$ and $\ord(z_5-C)\ge2$; so $C$, $z_1$, $z_2$
have order $\ge\theta$, and $z_k-C$ has order $\ge2$ for $k=3,4,5$.
\emph{Chart.}  The rotation lowers all $z$-orders by $\theta$ and raises all $w$-orders by
$\theta$, so $\ord(w_4-w_3)=\theta-2$ and $\ord t=(2-\theta)/6>0$.  By the inverse formulas
of the chart: $A_1,A_2,C$ have order $\ge0$; $u,v$ have order $\ge(2-\theta)-6\ord t=0$; $B_3,B_4$
have order $0$ and $D$ order $\ge0$; $\ord h=2\ord t-\ord r_{15}=-2\theta/3$; $R_{12},R_{25}$
have order $0$; $P,Q$ have order $\ge2\ord h+\frac{4+2\theta}3-4\ord t=0$, because
$w_1-t^{-6}D=(m_2w_{21}+m_5w_{51})/(m_1+m_2+m_5)$ and similarly for $w_2$; $D_{bc}=t^{-3}\delta_{bc}$
has order $\frac{2-\theta}2-3\ord t=0$; and $D_{kl}=\delta_{kl}$ has order $0$.  Hence every
coordinate converges (Lemma~\ref{lem:germ}), the inverse-distance coordinates to nonzero limits,
and $p=\ord h/\ord t=-4\theta/(2-\theta)\in(-4,2)$.  Finally
$\ord a=\ord t-\frac12\ord h=(2+\theta)/6$ and $\ord\eta=2\ord t+\frac12\ord h=(2-2\theta)/3$,
both positive.
\end{proof}
In particular no Diagram~8 germ reaches the lower projective end of the chart, where the lower
inverse distances degenerate.

\subsection{The first condition}

\begin{proposition}[First condition]\label{prop:d8-first}
At positive masses every Diagram~8 germ satisfies \eqref{eq:d8-40}:
$m_5^2(m_1+m_2)^4=m_1^3m_2^3$.  Moreover, in the chart of Proposition~\ref{prop:d8-cover} its
limit $x_0$ has $A_1=-m_2/M_{12}$, $A_2=m_1/M_{12}$, $C=0$, $R_{25}=1$,
$R_{12}=-m_5M_{12}/(m_1m_2)$, $(P,Q)=(P_0,Q_0)$ the unique solution of $f_{10}=f_{13}=0$,
$D=0$, $B_3=-m_4/(m_3+m_4)$, $B_4=m_3/(m_3+m_4)$, and, with $L:=(M_{12}(m_3+m_4))^{1/2}>0$ and
signs $s_{bc}\in\{\pm1\}$,
\begin{gather*}
  D_{13}=\frac{is_{13}L}{\sqrt{m_2m_4}},\qquad D_{14}=\frac{s_{14}L}{\sqrt{m_2m_3}},\\
  D_{23}=\frac{s_{23}L}{\sqrt{m_1m_4}},\qquad D_{24}=\frac{is_{24}L}{\sqrt{m_1m_3}}.
\end{gather*}
\end{proposition}
We write $U:=(A_1,A_2,C,R_{12},R_{25},P,Q)$, $U_0$ for its limit, and
$\mathcal I_0:=(X_0,Y_0,R_{12}^0,1,P_0,Q_0)$, with $X_0=-m_2/M_{12}$, $Y_0=m_1/M_{12}$ and
$R_{12}^0=-m_5M_{12}/(m_1m_2)$, for the limit of $\mathcal I$ (note $C_5\to C\to0$).

\begin{proof}
Let $m\in\Rp^5$ and let $\gamma$ be a Diagram~8 germ, in the chart of
Proposition~\ref{prop:d8-cover}, with limit $x_0$.
By Lemma~\ref{lem:d8-rows}\,(i) and Proposition~\ref{prop:d8-cover},
$E_j=f_j(x)+O(|a|+|\eta|)$ along $\gamma$, so $f_j(x_0)=0$ for all $j$.
\emph{Positions.}  By Lemma~\ref{lem:d8-rows}\,(ii) and $a\ne0$, $m_1A_1+m_2A_2+MC=0$ along
$\gamma$.  With $f_4=C$ and $f_{20}$ this linear system, of determinant $M_{12}$, gives
$A_1=-m_2/M_{12}$, $A_2=m_1/M_{12}$, $C=0$ at $x_0$; so $X=A_1$, $Y=A_2$ and $X-Y=-1$ there.
\emph{Force rows.}  Then $f_8=\frac{m_1m_2}{M_{12}}(R_{25}^3-1)$ and $f_0=\frac{m_2}{M_{12}}F$
with $F:=M_{12}R_{12}^3+m_5$, so $R_{25}^3=1$ and $F=0$.
\emph{Closure of the $w$-triangle $125$.}  Identically,
\begin{align*}
  &(X-Y)R_{12}^2Xf_{16}-(X-Y)R_{12}^2YR_{25}^2f_{13}+XYR_{25}^2f_{10}\\
  &\qquad=(X-Y)R_{12}^2(YR_{25}^2-X)-XYR_{25}^2,
\end{align*}
which is free of $P,Q$.  At $x_0$ the right side is $-G/M_{12}^2$ with
$G:=M_{12}R_{12}^2(m_1R_{25}^2+m_2)-m_1m_2R_{25}^2$; so $G=0$.
\emph{Elimination.}  Let $H:=M_{12}(m_1+m_2R_{25})R_{12}^2-m_1m_2$ and
$\Phi(s):=m_5^2M_{12}(m_1+m_2s)^3-(m_1m_2)^3$.  Identically
\begin{align*}
  H&=R_{25}G+(1-R_{25}^3)\,m_1(M_{12}R_{12}^2-m_2),\\
  \Phi(R_{25})&=H\bigl(H^2+3m_1m_2H+3m_1^2m_2^2\bigr)
  -M_{12}(m_1+m_2R_{25})^3F(F-2m_5),
\end{align*}
so $H=0$ and $\Phi(R_{25})=0$ at $x_0$.
\emph{Phase.}  If $R_{25}\ne1$, then $R_{25}^2+R_{25}+1=0$ and $\Phi(R_{25})=\Phi_0+\Phi_1R_{25}$
with real $\Phi_0,\Phi_1$; as $R_{25}\notin\R$, $\Phi_0=\Phi_1=0$.  But
$\Phi_1=3m_1m_2m_5^2M_{12}(m_1-m_2)$ forces $m_1=m_2$, and then $\Phi_0=-m_1^4(m_1^2+2m_5^2)<0$.
Hence $R_{25}=1$, and $\Phi(1)=0$ is \eqref{eq:d8-40}.
\emph{The limit point.}  At $R_{25}=1$, $F=0$ and $H=M_{12}^2R_{12}^2-m_1m_2=0$ give
$R_{12}=R_{12}^3/R_{12}^2=-m_5M_{12}/(m_1m_2)$.  The faces $f_{10},f_{13}$ are linear in $(P,Q)$
with determinant $m_2m_5(m_1+m_2+m_5)R_{12}^2/M_{12}\ne0$.  Next $f_1=D$ gives $D=0$; the
identity $m_1f_1+m_2f_3+m_3f_5+m_4f_7+m_5f_9=(m_1+m_2+m_5)D+m_3B_3+m_4B_4$ and $f_{21}$ give
$B_3$, $B_4$; and $f_{11},f_{12},f_{14},f_{15}$ give
$D_{13}^2=-M_{12}(m_3+m_4)/(m_2m_4)$, $D_{14}^2=M_{12}(m_3+m_4)/(m_2m_3)$,
$D_{23}^2=M_{12}(m_3+m_4)/(m_1m_4)$, $D_{24}^2=-M_{12}(m_3+m_4)/(m_1m_3)$.
\end{proof}
Conversely, a direct computation shows that on \eqref{eq:d8-40} the whole face system has
solutions, with the lower wing at the phase $D_{45}=D_{35}$ and any crossed signs; so
\eqref{eq:d8-40} is exactly the solvability condition of the faces together with the centre of
mass.  We do not use this.

\begin{remark}\label{rem:d8-ak}
\cite[\S8.8]{AK} place the cluster $z_3,z_4,z_5$ at the centre of mass of $z_1,z_2$, invoke
their relation (34), $(m_1+m_2)W_{12}\sim m_5W_{51}+m_5W_{25}$, and refer to their \S8.7.3.  In
the chart $(m_1+m_2)W_{12}-m_5(W_{51}+W_{25})=t^{-2}h(\beta_1-\beta_3)=t^{-6}\eta^2
(\beta_1-\beta_3)$, and $\eta^2(\beta_1-\beta_3)$ is the $\eta^2$-part of $E_1-E_3$; so (34)
holds at the limit of every germ (Lemma~\ref{lem:d8-rows}\,(iii)), and given $R_{25}^3=1$, the
ties and the centre of mass it is equivalent to $F=0$.  The centre of mass cannot be dropped: at
$(m_1,m_2,m_5)=(4,1,\frac49)$, which violates \eqref{eq:d8-40}, the point
$A_1=R_{12}=-\frac23$, $A_2=\frac13$, $C=0$, $R_{25}=2$, with $P,Q$ from $f_{10}=f_{13}=0$,
satisfies every upper face and (34), but $m_1A_1+m_2A_2+MC=-\frac73$.  Relation (39) is not
used: at $D=0$, $B_3=-m_4/(m_3+m_4)$, $B_4=m_3/(m_3+m_4)$, with $D_{35}=g$, $D_{45}=\omega g$,
$D_{34}=\lambda g$, the face $f_5$ reads $g^3\bigl((m_3+m_4)\lambda^3+m_5\bigr)=1$, and the
escape $(m_3+m_4)\lambda^3+m_5=0$ of the lower scale forces $\omega=1$ and (39); by
Proposition~\ref{prop:d8-cover} no Diagram~8 germ escapes.
\end{remark}

\subsection{Exclusion}
The germ is excluded in three steps: an order-six pairing for $0<p<2$; a lemma on the
coefficients of $t^9=T^3$ in $E_1$ and $E_3$, which for $-4<p<\frac12$ forces the order-nine
pairing $\Xi$ below to vanish; and a second-order elimination on the locus $\Xi=0$.

\begin{lemma}[computer-assisted]\label{lem:d8-order-six}
Let $m\in\Rp^5$ satisfy \eqref{eq:d8-40}, and let $S$ be the vector of the rows
$E_0,E_2,E_4,E_{10},E_{13},E_{16},E_{20}$ and $E_Z:=m_1A_1+m_2A_2+MC$.
\begin{romanlist}
\item The rows $E_0$ and $E_2$ consist of their faces and multiples of $\mu$, $\nu$,
$a^9\eta^3$, $a^{11}\eta^5$; $E_4$ of its face and multiples of $\nu$, $T^3$, $T^5$; $E_{13}$ and
$E_{16}$ of their faces and multiples of $\nu$; and $E_{10}$, $E_{20}$, $E_Z$ equal their faces.
The faces of $S$ depend on $U$ only.
\item The faces of $S$ vanish at $U_0$, and their Jacobian $J$ with respect to $U$ there (an
$8\times7$ matrix) has rank $7$.
\item Let $S_6$ be the vector of $\mu$-coefficients of $S$ at $x_0$ and let $\ell\ne0$ satisfy
$\ell J=0$.  If $m_1\ne m_2$, scale the masses to $m_2=1$, $m_1=\alpha^2$ (so
$m_5=\alpha^3/(\alpha^2+1)^2$) and normalize $\ell$ by $\ell_{E_4}=\alpha^2-1$; then
\begin{gather*}
  \ell\cdot S_6=-\frac{2(\alpha^2+1)^4}{\alpha\,\Delta},\\
  \Delta:=(m_3+m_4)(\alpha^2+1)^2+\alpha^6+3\alpha^4+\alpha^3+3\alpha^2+1>0.
\end{gather*}
If $m_1=m_2$ $(=4m_5)$, then $\ell$ is a multiple of
$(1,-1,0,\frac32m_5,-\frac34m_5,-\frac34m_5,0,0)$, and $\ell\cdot S_6=-1$ for this vector.
\end{romanlist}
Consequently no Diagram~8 germ at the masses $m$ has $0<p<2$.
\end{lemma}
Parts (i)--(iii) are exact computations, with $\alpha,m_3,m_4$ as indeterminates if
$m_1\ne m_2$, and with $m_3,m_4,m_5$ as indeterminates if $m_1=m_2=4m_5$
(Appendix~\ref{app:verification}).  For $m_1\ne m_2$ the rank in (ii) and the null vector in
(iii) are computed over the field $\mathbb Q(\alpha,m_3,m_4)$, and they specialize to every
$\alpha>0$, $\alpha\ne1$.  The maximal minors of $J$ that involve $\alpha$ only have the greatest
common divisor $\alpha(\alpha^2+\alpha+1)(\alpha^4-\alpha^3+3\alpha^2-\alpha+1)$, which has no
positive zero, so $J$ has rank $7$ at every $\alpha>0$.  The denominators of the entries of $J$,
of the normalized $\ell$ and of $S_6$ are products of $\alpha$, $\alpha^2+1$ and $\Delta$, which
do not vanish for $\alpha>0$ and $m_3,m_4>0$.  Hence the identities $\ell J=0$ and the formula
for $\ell\cdot S_6$ hold at every such $\alpha$, where $\ell\ne0$ because
$\ell_{E_4}=\alpha^2-1\ne0$.

\begin{proof}[Proof of the last assertion]
Let $\gamma$ be a Diagram~8 germ with $0<p<2$, that is $r>2$.  Then $\nu/\mu=(\eta/a^2)^2$,
$T^3/\mu=(\eta/a)^3$, $T^5/\mu=\eta^5/a$, $a^9\eta^3/\mu$ and $a^{11}\eta^5/\mu$ tend to $0$,
because $\ord\eta>2\ord a$.  By (i) and the convergence of $x$, $S=S^{(0)}(U)+\mu S_6(x)+o(\mu)$
along $\gamma$, where $S^{(0)}$ is the vector of faces.  As $S=0$, $S^{(0)}(U)=O(\mu)$, and (ii)
gives $U-U_0=O(\mu)$.  Pairing with $\ell$ and using $\ell\cdot S^{(0)}(U)=O(|U-U_0|^2)$, we get
$0=\mu\bigl(\ell\cdot S_6(x_0)+o(1)\bigr)$, which contradicts (iii) since $\mu=a^6\ne0$.
\end{proof}

\begin{lemma}[computer-assisted]\label{lem:d8-kappa}
Let $m\in\Rp^5$ satisfy \eqref{eq:d8-40}; let $x_0$, $\mathcal I_0$, $U_0$ and the signs
$s_{bc}$ be as in Proposition~\ref{prop:d8-first}, $\beta:=\beta_1-\beta_3$,
$b:=(f_0,f_2,f_{10},f_{13},f_{16},f_{20})$ as a function of
$(X,Y,R_{12},R_{25},P,Q)$, and $B:=\partial b/\partial\mathcal I$ at $\mathcal I_0$.
\begin{alphalist}
\item $\beta(\mathcal I_0)=\dfrac{m_1^3m_2^3-m_5^2(m_1+m_2)^4}{m_1m_2m_5(m_1+m_2)^2}$, which is
$0$ by \eqref{eq:d8-40}.
\item If $m_1\ne m_2$, then $B$ is invertible.  If $m_1=m_2$ $(=4m_5)$, then $\ker B$ is spanned
by a vector $k_0$, equal to $(\frac{27}{272},\frac{27}{272},\frac9{272},-\frac9{68},-\frac{19}{17},1)$
at $(m_1,m_2,m_5)=(1,1,\frac14)$; $d\beta(\mathcal I_0)\,k_0=0$; and the Jacobian of the eight
functions $f_0,f_2,f_{10},f_{13},f_{16},f_{20}$ (evaluated at $X=A_1-C$, $Y=A_2-C$), $C$ and
$m_1A_1+m_2A_2+MC$ with respect to $U$ at $U_0$ has rank $7$.  (These three facts are computed
at $(m_1,m_2,m_5)=(1,1,\frac14)$ with $m_3,m_4$ free, which suffices by the mass scaling.)
\item With $\Xi:=s_{14}(m_1m_4)^{3/2}+s_{23}(m_2m_3)^{3/2}
+i\bigl(s_{13}(m_1m_3)^{3/2}+s_{24}(m_2m_4)^{3/2}\bigr)$,
\[
  (\tau_1-\tau_3)(x_0)=\frac{m_3m_4}{m_3+m_4}\,L^3\,(m_1m_2m_3m_4)^{-3/2}\,\Xi .
\]
\end{alphalist}
Consequently, if a Diagram~8 germ at the masses $m$ has $-4<p<\frac12$, then $\Xi=0$, that is,
$m_1=m_2=4m_5$, $m_3=m_4$, $s_{13}=-s_{24}$ and $s_{14}=-s_{23}$.
\end{lemma}
The lemma is an exact computation (Appendix~\ref{app:verification}); (a) is an identity in the
masses, and (c), for each choice of the signs, an identity in their square roots.  For
$m_1\ne m_2$, in the scaling $m_2=1$, $m_1=\alpha^2$ and after clearing denominators row by
row, $\det B$ factors as $(\alpha-1)(\alpha+1)\alpha^{13}(\alpha^2+\alpha+1)(\alpha^2+1)^6
(\alpha^4-\alpha^3+3\alpha^2-\alpha+1)$, and its only positive root is $\alpha=1$.

\begin{proof}[Proof of the last assertion]
Let $\gamma$ be such a germ.  Along $\gamma$, $\mathcal I\to\mathcal I_0$ and
$b(\mathcal I)=O(a^6)$ by Lemma~\ref{lem:d8-rows}\,(iv).  We claim
\begin{equation}\label{eq:d8-bracket}
  \beta(\mathcal I)=O(a^6+a^4\eta^4).
\end{equation}
If $m_1\ne m_2$, (b) and the inverse function theorem give $\mathcal I-\mathcal I_0=O(a^6)$, and
\eqref{eq:d8-bracket} follows from (a).  If $m_1=m_2$, the eight functions in (b) are
$O(a^6+\nu+T^3)$ along $\gamma$: the first six because
$f_k(A_1-C,A_2-C,\dots)-f_k(\mathcal I)=O(\nu)$, the seventh by Lemma~\ref{lem:d8-rows}\,(v),
and the last vanishes by Lemma~\ref{lem:d8-rows}\,(ii) as $a\ne0$.  So $U-U_0$ and
$\mathcal I-\mathcal I_0$ are $O(a^6+\nu+T^3)$.  Since
$B(\mathcal I-\mathcal I_0)=O(a^6+|\mathcal I-\mathcal I_0|^2)$, the component of
$\mathcal I-\mathcal I_0$ transverse to $k_0$ is $O(a^6+|\mathcal I-\mathcal I_0|^2)$; as
$\beta(\mathcal I_0)=0$ and $d\beta(\mathcal I_0)k_0=0$, we get
$\beta(\mathcal I)=O(a^6+|\mathcal I-\mathcal I_0|^2)=O(a^6+\nu^2+T^6)$, which is
\eqref{eq:d8-bracket} because $\nu^2=a^4\eta^4$ and $T^6=a^6\eta^6$.

Now $E_1=E_3=0$ along $\gamma$, so Lemma~\ref{lem:d8-rows}\,(iii) and \eqref{eq:d8-bracket}
give $T^3(\tau_1-\tau_3)(x)=O(a^6\eta^2+a^4\eta^6)$, that is,
$(\tau_1-\tau_3)(x)=O(a^3\eta^{-1}+a\eta^3)$.  As $-4<p<\frac12$ means $r<3$,
$\ord(a^3\eta^{-1})=(3-r)\ord a>0$; so $(\tau_1-\tau_3)(x_0)=0$, and $\Xi=0$ by (c).  Since all
$m_k>0$, the real part $s_{14}(m_1m_4)^{3/2}+s_{23}(m_2m_3)^{3/2}$ of $\Xi$ vanishes only if
$s_{14}=-s_{23}$ and $m_1m_4=m_2m_3$, and its imaginary part only if $s_{13}=-s_{24}$ and
$m_1m_3=m_2m_4$.  These give $m_3=m_4$ and $m_1=m_2$, and then \eqref{eq:d8-40} gives
$m_1=4m_5$.
\end{proof}

On the locus $\Xi=0$, which is the \emph{cancelling curve} $m=(4,4,y,y,1)$, $y>0$, up to
scaling, together with the four sign classes $s_{13}=-s_{24}$, $s_{14}=-s_{23}$, the order-nine
pairing gives nothing, and we expand to second order.

\begin{lemma}[computer-assisted]\label{lem:d8-cancel}
Let $m_1=m_2=4m_5$ and $m_3=m_4$, scaled to $(1,1,m,m,\frac14)$ with $m>0$.  Let $\nabla$ and
$\nabla^2$ denote gradient and Hessian with respect to $\mathcal I$ at $\mathcal I_0$, and let
$k_0$ be as in Lemma~\ref{lem:d8-kappa}\,(b).
\begin{alphalist}
\item \emph{Lower block.}  The faces $f_1,f_5,f_7,f_9,f_{17},f_{18},f_{19},f_{21}$ involve only
\[
  x':=(B_3,B_4,D,u,v,D_{34},D_{35},D_{45}).
\]
At each common zero with $g:=D_{35}\ne0$, writing
$D_{45}=\omega g$ and $D_{34}=\lambda g$: $D=0$, $B_3=-\frac12$, $B_4=\frac12$, $\omega^3=1$,
$2\lambda^2(\omega^2+1)=\omega^2$, $g^3(8m\lambda^3+1)=4$, $f_{18},f_{19}$ determine $u,v$, and
the Jacobian determinant of these eight faces with respect to $x'$ is
\[
  -9g^7\lambda m(8m+1)(\omega^2+1)(8m\lambda^3+1)/(128\,\omega)\ne0.
\]
\item \emph{Lower forcing.}  In $E_1,E_3,E_5,E_7,E_9,E_{17},E_{18},E_{19},E_{21}$ and in the
crossed rows $E_{11},E_{12},E_{14},E_{15}$, every monomial $a^i\eta^j$ other than the face has
$i\ge2$ and $i+2j\ge6$, except $\eta^2$ in $E_1,E_3,E_9$, whose coefficients vanish at
$\mathcal I_0$.
\item \emph{Upper rows.}  Below $a^8$ the rows $E_0,E_2,E_8,E_{10},E_{13},E_{16},E_{20}$ consist
of $f_k(\mathcal I)$ and the terms $\mu A_1$, $\mu A_2$, $\mu C$ of $E_0,E_2,E_8$.  The row $E_1$
(resp.\ $E_3$) consists of $D$, $\eta^2\beta_1$, $T^3\tau_1$, $m_5\mu\eta^2P$ (resp.\ $D$,
$\eta^2\beta_3$, $T^3\tau_3$, $m_5\mu\eta^2Q$) and a multiple of $a^9\eta^5$.  The Jacobian
$\nabla(f_0,f_2,f_8,f_{10},f_{13},f_{16},f_{20})$ has rank $5$ and kernel $\C k_0$.  The Jacobian
$B_5$ of $(f_2,f_{10},f_{13},f_{16},f_{20})$ with respect to $(X,Y,R_{12},R_{25},P)$ has
$\det B_5=-\frac{17}{16}$, and $\varphi:=m_1X+m_2Y$ satisfies $\varphi(k_0)=\frac{27}{136}$.
\item \emph{Second-order constants.}  Let $\delta_\mu,\delta_2\in\C^6$ have last entry $0$ and
first five entries $-B_5^{-1}(Y_0,0,0,0,0)^{\mathsf T}$ and
$-B_5^{-1}\bigl(\frac12k_0^{\mathsf T}\nabla^2f_k\,k_0\bigr)_{k=2,10,13,16,20}$, and put
\begin{gather*}
  c_0:=X_0+\nabla f_0\cdot\delta_\mu,\qquad
  q_2:=\tfrac12k_0^{\mathsf T}\nabla^2f_0\,k_0+\nabla f_0\cdot\delta_2,\\
  e_2:=\nabla\beta\cdot\delta_2+\tfrac12k_0^{\mathsf T}\nabla^2\beta\,k_0,\qquad
  c_{62}:=\nabla\beta\cdot\delta_\mu+m_5(P_0-Q_0).
\end{gather*}
Then $c_0=-1$, $q_2=\frac{243}{73984}$, $e_2=\frac{1215}{18496}$, $c_{62}=-4$ and
$c_0e_2-q_2c_{62}=-\frac{243}{4624}$.
\end{alphalist}
Consequently no Diagram~8 germ at these masses has a limit with $s_{13}=-s_{24}$ and
$s_{14}=-s_{23}$.
\end{lemma}
The lemma is an exact computation over the field $\mathbb Q(m)$; in (a) the Jacobian determinant
is reduced modulo $\omega^3-1$, and in (c) and (d) all constants are rational numbers
(Appendix~\ref{app:verification}).

\begin{proof}[Proof of the last assertion]
Let $\gamma$ be such a germ, with limit $x_0$.  Then $\Xi=0$, so $(\tau_1-\tau_3)(x_0)=0$ by
Lemma~\ref{lem:d8-kappa}\,(c), and $D_{35}(x_0)\ne0$ by Proposition~\ref{prop:d8-cover}.  Write
$\mathcal I=\mathcal I_0+\rho k_0+\delta$ with $\rho:=Q-Q_0$ and $\delta$ of last entry $0$.  All
estimates are along $\gamma$.

\emph{Upper rows.}  By (c), $E_2,E_{10},E_{13},E_{16},E_{20}$ give, for the first five entries
of $\delta$,
$B_5\delta=-\mu(Y_0,0,0,0,0)^{\mathsf T}-\frac12\rho^2(k_0^{\mathsf T}\nabla^2f_kk_0)_k
+o(\mu+\rho^2)$, so $\delta=\mu\delta_\mu+\rho^2\delta_2+o(\mu+\rho^2)$, and $E_0=0$ becomes
\begin{equation}\label{eq:d8-row0}
  c_0\mu+q_2\rho^2=o(\mu+\rho^2).
\end{equation}
\emph{Lower rows.}  By (a), (b) and the implicit function theorem (each crossed face
$f_{11},f_{12},f_{14},f_{15}$ contains one crossed unknown $D_{bc}$, with derivative
$2/D_{bc}\ne0$ at a zero), the lower and crossed coordinates differ from their limits by
$O(\nu+\mu+|\mathcal I-\mathcal I_0|)=O(\nu+\mu+|\rho|)$.

\emph{Amplitude.}  The centre of mass (Lemma~\ref{lem:d8-rows}\,(ii)), $E_4=0$
(Lemma~\ref{lem:d8-rows}\,(v)) and $C_5=C-\nu(m_3u+m_4v)$ give the identity
\[
  \varphi(\mathcal I)=(M_{12}+M)\bigl(\nu\Psi+T^3F_3+T^5G_3\bigr),
\]
where $\Psi:=L_3+M_{12}(m_3u+m_4v)/(M_{12}+M)$.
As $\varphi(\mathcal I_0)=0$ and $\varphi(\delta)=O(\mu+\rho^2)$, first $\rho=O(\nu+T^3+\mu)$,
and then, by the lower estimate, $\rho=\rho_\nu\nu+O(T^3+\nu^2+\mu)$ with
$\rho_\nu:=(M_{12}+M)\Psi(x_0)/\varphi(k_0)$.

\emph{Case $\Psi(x_0)=0$} (it contains the lower phase $\omega=1$).  Then
$\rho^2=O(T^6+\nu^4+\mu^2)=o(\mu)$, and \eqref{eq:d8-row0} reads $\mu(c_0+o(1))=0$, which is
impossible.

\emph{Case $\Psi(x_0)\ne0$.}  Then $\rho^2=\rho_\nu^2a^4\eta^4+o(a^6+a^4\eta^4)$, because
$\nu T^3=a^5\eta^5=o(a^6+a^4\eta^4)$; so \eqref{eq:d8-row0} reads
$c_0a^6+q_2\rho_\nu^2a^4\eta^4=o(a^6+a^4\eta^4)$.  Both coefficients are nonzero, so the two
monomials have the same order, that is $r=\frac12$ ($p=-2$), and
$c_0a_0^2+q_2\rho_\nu^2\eta_0^4=0$ for the leading coefficients $a_0,\eta_0\ne0$ of $a,\eta$.  On
this ray, (c) gives
$E_1-E_3=\eta^2\beta(\mathcal I)+T^3(\tau_1-\tau_3)(x)+m_5\mu\eta^2(P-Q)+O(a^9\eta^5)$, where
$\beta(\mathcal I)=\mu\,\nabla\beta\cdot\delta_\mu+e_2\rho^2+o(\mu+\rho^2)$ (because
$\beta(\mathcal I_0)=0$ and $\nabla\beta\cdot k_0=0$, Lemma~\ref{lem:d8-kappa}) and
$T^3(\tau_1-\tau_3)(x)=O\bigl(T^3(\nu+\mu+|\rho|)\bigr)=O(a^5\eta^5+a^6\eta^6+a^9\eta^3)$.  With
weight $1$ for $a$ and $\frac12$ for $\eta$, the monomials $a^6\eta^2$ and $a^4\eta^6$ have
weight $7$, while $a^5\eta^5$, $a^6\eta^6$, $a^9\eta^3$, $a^9\eta^5$ have weights
$\frac{15}2,9,\frac{21}2,\frac{23}2$.  Hence $E_1-E_3=0$ gives
$c_{62}a_0^2+e_2\rho_\nu^2\eta_0^4=0$.  Since $c_0e_2-q_2c_{62}\ne0$, the two linear equations
force $a_0=0$, a contradiction.
\end{proof}

\begin{theorem}\label{thm:d8}
No Diagram~8 germ exists at positive masses: for $m\in\Rp^5$ no curve in $A_m$ carries a
Diagram~8 germ, complex germs included.  In particular Theorem~\ref{thm:exclusion}\,(a) holds
for $r=8$.
\end{theorem}

\begin{proof}
By the symmetries at the beginning of this section we may use the numbering of
\cite[\S8.8]{AK}.  Let $\gamma$ be a Diagram~8 germ at $m\in\Rp^5$.  By
Proposition~\ref{prop:d8-cover} it is a ray germ with ray $r\in(0,\infty)$, and by
Proposition~\ref{prop:d8-first} the masses satisfy \eqref{eq:d8-40} and the limit has signs
$s_{bc}$.  If $r>2$, that is $0<p<2$, Lemma~\ref{lem:d8-order-six} is contradicted.  If $r\le2$,
that is $-4<p\le0$, Lemma~\ref{lem:d8-kappa} gives $m_1=m_2=4m_5$, $m_3=m_4$, $s_{13}=-s_{24}$
and $s_{14}=-s_{23}$, and Lemma~\ref{lem:d8-cancel} is contradicted.  So $\gamma$ does not exist.
A Diagram~8 germ with an unbounded four-product is in particular a Diagram~8 germ.
\end{proof}

\begin{remark}\label{rem:d8-window}
\begin{alphalist}
\item The window is sharp.  Along rays of the ratio chart with a finite lower wing, the strokes
and circles are exactly those of Diagram~8 for $-4<p<2$; at $p=2$ the $z$-circles at $1$ and
$2$ join the upper $z$-strokes, which is the pattern of Diagram~6; for $p>2$ there would be
$z$-circles at $1,2$ without $z$-strokes, against Rule~1a; and for $p\le-4$ the upper
$w$-strokes reach the order of the lower $w$-stroke $34$.  So the face $p=2$ of the chart is
never reached by a Diagram~8 germ (Proposition~\ref{prop:d8-cover}); nothing is excluded or
delegated there, and germs on that face are not Diagram~8 germs.
\item Along a ray germ the squared distances have orders $4-2p$ (upper wing), $-6$ (crossed
pairs) and $0$ (lower wing) in $t$, so the least four-product has order $-8-2p<0$: every
Diagram~8 germ has an unbounded four-product.  We do not use this.
\end{alphalist}
Both items follow from the orders in $t$ of the components of $\mathcal Z$ and $\mathcal W$
along the rays (Appendix~\ref{app:verification}).
\end{remark}



\appendix

\section{The computations}\label{app:verification}

The computer-assisted steps are carried out by programs written in Python~3.14 with the
libraries SymPy~1.14 and mpmath~1.3.  Each program has tests that compare its output with the
statement in the text.  Statements proved by hand in the text were also checked by computer
where this was easy; these checks are not listed here.

\emph{Exact algebra.}  Identities, resultants, greatest common divisors, Gr\"obner bases, ranks
and determinants are computed over $\mathbb Q$, over number fields such as $\mathbb Q(i)$,
$\mathbb Q(\omega)$ and $\mathbb Q(\sqrt3,i,3^{-1/3})$, and over fields of rational functions.
An identity is proved by expanding the difference of its two sides to the zero polynomial.  Real
roots are isolated by Sturm sequences, and signs of polynomials on boxes are obtained from
Bernstein coefficients.  A rank or a nonvanishing obtained over a field of rational functions is
used at a point only after the relevant denominators and minors have been shown not to vanish
there, as in Lemma~\ref{lem:d8-order-six}.  A computation modulo a prime is used to show that a
quantity is nonzero only through a minor that is a unit modulo that prime, as in
Lemma~\ref{lem:d3-order6}.

\emph{Interval and ball arithmetic.}  Real and complex numbers are enclosed by interval
arithmetic with outward rounding.  Roots are enclosed in discs by the Krawczyk and Pellet tests
and by Rouch\'e's theorem.  The covers of Section~\ref{sec:d2-cases} also use complex balls built
on IEEE binary64 arithmetic, with relative error bounds $2u+2u^2$ for the modulus, $2.5u$ for a
product, $3u+7u^2$ for an inverse and $3.5u+8u^2$ for a square root, where $u=2^{-53}$;
these bounds assume round-to-nearest and no overflow.  The ball operations use a common $8u$
allowance for these primitives.  Sums of $n$ terms are enlarged by the factors $(1+u)^n$ and
$\gamma_n=nu/(1-nu)$, every radius is increased by $10^{-300}$ to account for underflow, and
nonzero arguments of modulus below $2^{-1000}$ are rejected.

\emph{Exact linear programming.}  The systems of rational linear equations and inequalities of
Section~\ref{sec:d5} are decided by the simplex method over $\mathbb Q$, with Bland's rule and free
variables; strict inequalities are encoded with a common slack $0<s\le1$.

\emph{Stored results.}  The tables of Lemmas~\ref{d5-lem-census} and~\ref{d5-lem-four}, the
zero classes of Lemma~\ref{lem:d4-census}, the representatives of Lemma~\ref{lem:d2-K} and the
jets of the curves $C_p$ of Lemma~\ref{d6-lem-corank} are stored with the programs.  Designated
verification checks compare pinned witnesses with their recorded SHA-256 digests; an individual
helper may read a witness before that separate check runs.  A matching digest establishes the
identity of the stored data, not the correctness of its calculation.  The pinned checks test
selected identities and witness entries; a separate full replay recomputes the tables of
Section~\ref{sec:d5} and the representatives of Lemma~\ref{lem:d2-K} from scratch.  The zero
classes are checked by root isolation, and the jets by substitution into the equations.

\emph{Programs and running times.}  Table~\ref{tab:verification} lists the programs used for
the computer-assisted statements, including supplementary remarks not needed for the main
theorem; their names begin with \code{ak\_}, which is omitted.  The selected proof verification
checks do not replay every supplementary remark or the cross-check of Appendix~\ref{app:census}.
On a computer with eight logical processors, the tests of these programs take about $15$ hours
of processor time, and the recomputation of the tables of Section~\ref{sec:d5} about $4$ more
hours.

\begin{table}[ht]
\centering
\footnotesize
\begin{tabular}{p{4.6cm}p{8.4cm}}
\hline
Statement & Programs \\
\hline
Lemmas \ref{lem:d1-chart}, \ref{lem:d1-virial} & \code{diagram1\_separated\_center\_cone},
\code{diagram1\_virial\_exclusion}, \code{diagram5\_valuation\_reduction} \\
Lemma \ref{lem:d1-torque} & \code{central\_body\_slice} \\
Lemma \ref{lem:d1-shape} & \code{isolated\_triangle\_shape\_derivation} \\
Lemma \ref{lem:d1-mb} & \code{diagram5\_valuation\_reduction}, \code{diagram1\_edge\_triangle} \\
\hline
Lemma \ref{lem:d2-fibre}, Definition \ref{def:d2-mass} & \code{blowup}, \code{resultants},
\code{p\_source\_completeness}, \code{p\_general}, \code{diagram2\_symmetry} \\
Lemma \ref{lem:d2-second} & \code{l\_block\_symmetric\_germs}, \code{p\_sheet\_second\_order} \\
Lemma \ref{lem:d2-empty} & \code{l4\_same\_sign}, \code{resultants}, \code{exceptional},
\code{l5\_obstruction} \\
Lemmas \ref{lem:d2-L5}, \ref{lem:d2-corank-two} & \code{l5\_obstruction} \\
Lemmas \ref{lem:d2-rays}, \ref{lem:d2-Lray} & \code{exceptional\_branch\_positive\_feasibility},
\code{l1\_complex\_fibre}, \code{l\_block\_symmetric\_germs} \\
Lemma \ref{lem:d2-K} & \code{k1\_curve}, \code{k1\_complex\_fibre}, \code{p4\_curve},
\code{p5\_completion}, \code{resultants}, \code{diagram2\_symmetry} \\
Lemmas \ref{lem:d2-P}, \ref{lem:d2-fold} & \code{p\_sheet\_second\_order},
\code{p\_sheet\_special\_ends}, \code{p\_sheet\_e5\_corner}, \code{p\_sheet\_euler\_closed\_form},
\code{p\_sheet\_p5\_exact} \\
Remark \ref{rem:d2-scope} & \code{p\_general}, \code{diagram2\_symmetry} \\
\hline
Lemmas \ref{lem:d3-fibre}--\ref{lem:d3-census} & \code{diagram3\_blowup} \\
Lemma \ref{lem:d3-order5} & \code{diagram3\_fifth\_order} \\
Lemmas \ref{lem:d3-order6}, \ref{lem:d3-twisted} & \code{diagram3\_equal\_shape} \\
\hline
Lemma \ref{lem:d4-torque} & \code{diagram4\_valuation\_coverage} \\
Lemma \ref{lem:d4-census} & \code{diagram4\_first\_order\_formula}, \code{diagram4\_sign\_census},
\code{diagram4\_sign\_orbits}, \code{diagram4\_zero\_uniqueness},
\code{diagram4\_third\_order\_formula} \\
\hline
Lemmas \ref{d5-lem-mb}, \ref{d5-lem-consequences} & \code{diagram5\_morse\_bott\_lemma},
\code{diagram5\_morse\_bott\_uniformity} \\
Lemmas \ref{d5-lem-census}, \ref{d5-lem-four} & \code{diagram5\_outer\_census},
\code{diagram5\_census\_model}, \code{diagram5\_first\_limit\_relevance}, \code{tropical\_census} \\
\hline
Lemmas \ref{d6-lem-chart}, \ref{d6-lem-torque} & \code{diagram6\_blowup},
\code{diagram6\_valuation\_coverage}, \code{diagram6\_relation32\_derivation} \\
Lemma \ref{d6-lem-atlas}, Definition \ref{d6-def-labels}, Remarks \ref{d6-rem-regular},
\ref{d6-rem-projective} & \code{diagram6\_blowup} \\
Lemma \ref{d6-lem-corank} & \code{diagram6\_blowup}, \code{diagram6\_primitive\_cubic},
\code{diagram6\_primitive\_quadratic}, \code{diagram6\_primitive\_opposite} \\
\hline
Lemma \ref{lem:d8-rows} & \code{diagram8\_relation40\_derivation}, \code{diagram8\_kappa\_row} \\
Lemma \ref{lem:d8-order-six} & \code{diagram8\_order\_six\_pairing},
\code{diagram8\_equal\_pair\_germs} \\
Lemma \ref{lem:d8-kappa} & \code{diagram8\_kappa\_row} \\
Lemma \ref{lem:d8-cancel} & \code{diagram8\_cancelling\_curve} \\
Remark \ref{rem:d8-window} & \code{diagram8\_ray\_diagrams}, \code{diagram8\_pattern\_boundaries} \\
\hline
\end{tabular}
\caption{Programs used for the computer-assisted statements.}
\label{tab:verification}
\end{table}

\section{Cross-check through the exceptional set}\label{app:census}

The exceptional set of \cite{AK} is not used in the proof.  As an independent check of the
case analysis we rebuilt it.  The pairs of first and second limits in the proof of
\cite[Thm.~6]{AK}, with all labellings of the bodies, give $171$ components of codimension two in
the mass space and $661$ orbits of pairs of diagrams under the symmetries.  Of these orbits,
$255$ contain no positive mass vector, and each of the other $406$ is excluded by
Sections~\ref{sec:d1}--\ref{sec:d8} applied to its first limit; no orbit requires an argument
about its second limit.


\begin{thebibliography}{99}

\bibitem{AK}
A.~Albouy and V.~Kaloshin, \emph{Finiteness of central configurations of five bodies in the
plane}, Ann. of Math. (2) \textbf{176} (2012), 535--588.

\bibitem{BCR}
J.~Bochnak, M.~Coste and M.-F.~Roy, \emph{Real Algebraic Geometry}, Ergeb. Math. Grenzgeb. (3)
\textbf{36}, Springer, 1998.

\bibitem{CC1}
K.-M.~Chang and K.-C.~Chen, \emph{Toward finiteness of central configurations for the planar
six-body problem by symbolic computations. (I) Determine diagrams and orders}, J. Symbolic
Comput. \textbf{123} (2024), 102277; arXiv:2303.02853.

\bibitem{CC2}
K.-M.~Chang and K.-C.~Chen, \emph{Toward finiteness of central configurations for the planar
six-body problem by symbolic computations. (II) Determine mass relations}, SIAM J. Appl. Dyn.
Syst. \textbf{24} (2025), 2369--2404.

\bibitem{DH}
Y.~Deng and M.~Hampton, \emph{Equilateral chains and cyclic central configurations of the
planar $5$-body problem}, arXiv:2112.06755, 2021.

\bibitem{HJ}
M.~Hampton and A.~N.~Jensen, \emph{Finiteness of spatial central configurations in the
five-body problem}, Celestial Mech. Dynam. Astronom. \textbf{109} (2011), 321--332.

\bibitem{HM}
M.~Hampton and R.~Moeckel, \emph{Finiteness of relative equilibria of the four-body problem},
Invent. Math. \textbf{163} (2006), 289--312.

\bibitem{JL}
A.~N.~Jensen and A.~Leykin, \emph{Smale's 6th problem for generic masses}, J. Exp. Math.
\textbf{1} (2025), 306--314; arXiv:2301.02305.

\bibitem{LS}
T.-L.~Lee and M.~Santoprete, \emph{Central configurations of the five-body problem with equal
masses}, Celestial Mech. Dynam. Astronom. \textbf{104} (2009), 369--381,
doi:10.1007/s10569-009-9219-0.

\bibitem{MZ19}
M.~Moczurad and P.~Zgliczy\'nski, \emph{Central configurations in planar $n$-body problem with
equal masses for $n=5,6,7$}, Celestial Mech. Dynam. Astronom. \textbf{131} (2019),
doi:10.1007/s10569-019-9920-6; arXiv:1812.07279.

\bibitem{MZ26}
M.~Moczurad and P.~Zgliczy\'nski, \emph{Central configurations with unequal masses: finiteness
in several exceptional cases of five bodies}, arXiv:2601.01165, 2026.
\bibitem{Moulton}
F.~R.~Moulton, \emph{The straight line solutions of the problem of $n$ bodies}, Ann. of Math. (2)
\textbf{12} (1910), 1--17.

\bibitem{Roberts}
G.~E.~Roberts, \emph{A continuum of relative equilibria in the five-body problem}, Phys. D
\textbf{127} (1999), 141--145.

\bibitem{Smale}
S.~Smale, \emph{Mathematical problems for the next century}, Math. Intelligencer \textbf{20}
(1998), no.~2, 7--15.

\bibitem{Wintner}
A.~Wintner, \emph{The Analytical Foundations of Celestial Mechanics}, Princeton University
Press, 1941.

\bibitem{YZ}
X.~Yu and S.~Zhu, \emph{On the finiteness of four-body central configurations},
arXiv:2103.08906, 2021.

\end{thebibliography}

\end{document}
